\documentclass[11pt, a4paper, openany]{book}
\usepackage[scale=0.8, centering]{geometry}

\usepackage{bm}
\usepackage{mathtools, amsthm, amssymb, amsmath, tikz-cd, mathrsfs, derivative, stmaryrd, bbm}
\usepackage[dvipsnames]{xcolor}
\usepackage[hidelinks]{hyperref}

\newtheorem{thm}{Theorem}\newtheorem{lem}{Lemma}\newtheorem{crl}{Corollary}
\theoremstyle{definition}
\newtheorem{define}{Definition}\newtheorem{exam}{Example}\newtheorem{remark}{Remark}

\DeclareMathOperator{\Corr}{Corr}
\DeclareMathOperator{\id}{id}
\DeclareMathOperator{\src}{s}
\DeclareMathOperator{\tgt}{t}
\DeclareMathOperator{\Hom}{Hom}
\DeclareMathOperator{\End}{End}
\DeclareMathOperator{\Aut}{Aut}
\DeclareMathOperator{\dif}{d}
\DeclareMathOperator{\im}{im}
\DeclareMathOperator{\symdif}{\triangle}
\DeclareMathOperator{\vol}{vol}
\DeclareMathOperator{\diam}{diam}
\DeclareMathOperator{\RE}{Re}
\DeclareMathOperator{\IM}{Im}
\DeclareMathOperator{\supp}{supp}
\DeclareMathOperator{\cl}{cl}
\DeclareMathOperator{\Var}{Var}
\DeclareMathOperator{\Cov}{Cov}
\DeclareMathOperator{\sd}{sd}

\newcommand{\mf}[1]{\overline{\mathscr{#1}}}
\newcommand{\ph}{\underline{\hspace{0.3cm}}}
\newcommand{\sph}{\underline{\hspace{0.2cm}}}
\newcommand{\cat}[1]{\textup{\textsf{#1}}}
\newcommand{\op}[1]{{#1}^\textup{op}}
\newcommand{\er}{\overline{\mathbb{R}}}
\newcommand{\ud}{\textup{u.d.}}
\newcommand{\ul}{\amalg}\newcommand{\one}{\mathbbm{1}}

\DeclarePairedDelimiterX\braket[2]{\langle}{\rangle}{#1\delimsize\vert\mathopen{}#2}
\DeclarePairedDelimiterX\bracket[2]{\langle}{\rangle}{#1,\mathopen{}#2}

\newenvironment{tight_enumerate}{\begin{enumerate}\itemsep -0.2em}{\end{enumerate}}

\title{Measure and Probability Theory}
\author{}
\date{}

\begin{document}
\frontmatter\maketitle\tableofcontents\mainmatter
\setlength{\abovedisplayskip}{6pt}
\setlength{\belowdisplayskip}{6pt}

Let $\alpha\in\mathbb{R}$. $]0,\infty[\,\,\rightarrow \,\,]0,\infty[ \,,\,\,x\mapsto x^{\alpha}$ can be continuously extended to a function $[0,\infty]\rightarrow [0,\infty]$.

\chapter{Systems of Sets}
    \begin{define}
        A set $\mathscr{C}\subset\mathscr{P}X$ of subsets of $X$ is said to be
        \begin{tight_enumerate}
            \item \textit{$\cup$-stable} if $A,B\in \mathscr{C}\Rightarrow A\cup B\in\mathscr{C}$.
            \item \textit{$\cap$-stable} if $A,B\in \mathscr{C}\Rightarrow A\cap B\in\mathscr{C}$.
            \item \textit{$\setminus$-stable} if $A,B\in \mathscr{C}\Rightarrow A\setminus B\in\mathscr{C}$.
            \item \textit{$\sqcup$-stable} if $A,B\in \mathscr{C}\text{ disjoint}\Rightarrow A\sqcup B\in\mathscr{C}$.
            \item \textit{complement-stable} if $A\in\mathscr{C}\Rightarrow A^c\in\mathscr{C}$.
            \item \textit{stable under countable union} if
            \[
                (A_n)\in\mathscr{C}^\mathbb{N}\Longrightarrow \bigcup_{n=1}^\infty A_n\in\mathscr{C}
            \]
            and similarly for it to be \textit{stable under countable intersection}.
            \item \textit{stable under countable disjoint union} if
            \[
                (A_n)\in \mathscr{C}^\mathbb{N}\text{ pairwise disjoint}\Longrightarrow \bigsqcup_{n=1}^\infty A_n\in\mathscr{C}
            \]
        \end{tight_enumerate}
    \end{define}
    \begin{thm}\label{setminus stable implies cap stable}
        Let $\mathscr{C}\subset\mathscr{P}X$.
        \[
            \mathscr{C}\text{ is $\setminus$-stable}\Longrightarrow\mathscr{C}\text{ is $\cap$-stable}
        \]
    \end{thm}
    \begin{proof}
        $A\cap B=A\setminus(A\setminus B)$.
    \end{proof}
    \begin{thm}\label{countable union implies binary union}
        Let $\mathscr{C}\subset\mathscr{P}X$.
        \begin{gather*}
            \mathscr{C}\text{ is stable under countable union}\Longrightarrow \mathscr{C}\text{ is $\cup$-stable}\\
            \varnothing\in\mathscr{C}\text{ and }\mathscr{C}\text{ is stable under countable disjoint union}\Longrightarrow \mathscr{C}\text{ is $\sqcup$-stable}
        \end{gather*}
    \end{thm}
    \begin{proof}
        $A\cup B=A\cup B\cup B\cup\cdots$. $A\sqcup B=A\sqcup B\sqcup \varnothing\sqcup\varnothing\sqcup \cdots$.
    \end{proof}
\section{\texorpdfstring{$\sigma$}{Sigma}-rings and \texorpdfstring{$\sigma$}{sigma}-algebras}
    Recall that $(\mathscr{P}X,\cup,\cap)$ is a Boolean algebra and $(\mathscr{P}X,\symdif,\cap)$ is the corresponding Boolean ring.
\subsection{Rings of sets}
    \begin{define}
        A \textit{ring} (resp. \textit{algebra}) \textit{of sets} on $X$ is a subrng (resp. subring) of $(\mathscr{P}X,\symdif,\cap)$.
    \end{define}
    \begin{thm}\label{characterization of rings}
        Let $\varnothing\neq \mathscr{R}\subset\mathscr{P}X$.
        \[
            \mathscr{R}\text{ is a ring}\,\,\Longleftrightarrow\,\, (\mathscr{R}\text{ is $\setminus$-stable})\wedge (\mathscr{R}\text{ is $\cup$-stable})  %closed under finite colimits and difference?
        \]
    \end{thm}
    \begin{proof}
        Recall that in a Boolean ring, a subset is an additive subgroup if and only if it is nonempty and closed under addition. Thus, $\mathscr{R}$ is a subrng if and only if it is nonempty and closed under $\symdif$ and $\cap$, which implies closure under $\cup$ and $\setminus$ since $A\cup B=(A\symdif B)\symdif(A\cap B)$ and $A\setminus B=A\symdif(A\cap B)$. Conversely, $A\symdif B=(A\setminus B)\cup (B\setminus A)$ and recall Theorem \ref{setminus stable implies cap stable}.
    \end{proof}
    \begin{thm}\label{every countable union is a disjoint union in a ring}
        Every countable union of sets in a ring $\mathscr{R}$ can be written as a countable disjoint union of sets in $\mathscr{R}$:
        \[
            \bigcup_{n=1}^\infty A_n=\bigsqcup_{n=1}^\infty A_n'
        \]
        where $A_n'=A_n\setminus(A_1\cup\cdots \cup A_{n-1})\in\mathscr{R}$ and $A_n'\subset A_n$.
    \end{thm}
\subsection{Algebras of sets}
    \begin{define}
        A \textit{ring} (resp. \textit{algebra}) \textit{of sets} on $X$ is a subrng (resp. subring) of $(\mathscr{P}X,\symdif,\cap)$.
    \end{define}
    \begin{thm}\label{characterization of algebras}
        Let $\varnothing\neq \mathscr{A}\subset\mathscr{P}X$.
        \begin{align*}
            \mathscr{A}\text{ is an algebra}\,\,&\Longleftrightarrow\,\,(\mathscr{A}\text{ is complement-stable})\wedge (\mathscr{A}\text{ is $\cup$-stable})\\
            &\Longleftrightarrow\,\,(\mathscr{A}\text{ is complement-stable})\wedge (\mathscr{A}\text{ is $\cap$-stable})
        \end{align*}
    \end{thm}
    \begin{proof}
            Since subring = subrng + 1, by Theorem \ref{characterization of rings}, an algebra is equivalently a $\setminus$-stable and $\cup$-stable system that contains $X$. Hence $\Rightarrow$ is immediate. Conversely, since $\mathscr{A}\neq\varnothing$, there exists $A\in\mathscr{A}$ such that $X=A\cup A^c\in\mathscr{A}$. By De Morgan's law,$A\setminus B=(A^c\cup B)^c$, and $A\cup B=(A^c\cap B^c)^c$.
    \end{proof}
\subsection{\texorpdfstring{$\sigma$}{sigma}-rings}
    \begin{define}
        A $\sigma$\textit{-ring} is a ring stable under countable disjoint union.
    \end{define}
    \begin{thm}\label{property of sigma rings}
        If $\mathscr{R}$ is a $\sigma$-ring, then
        \[
            (A_n)\in \mathscr{R}^\mathbb{N}\Longrightarrow \bigcup_{n=1}^\infty A_n,\,\bigcap_{n=1}^\infty A_n,\,\varlimsup_{n\to\infty}A_n,\,\varliminf_{n\to\infty}A_n\in\mathscr{R}
        \]
    \end{thm}
    \begin{proof}
        Theorem \ref{every countable union is a disjoint union in a ring}. Let $A=\bigcup_{n=1}^\infty A_n$. By De Morgan's law, $\bigcap_{n=1}^\infty A_n=A\setminus \bigcup_{n=1}^\infty (A\setminus A_n)$.
    \end{proof}
\subsection{\texorpdfstring{$\sigma$}{sigma}-algebras}
    \begin{define}\label{definition of sigma algebra}
        A $\sigma$\textit{-algebra} is an algebra stable under countable disjoint union.
    \end{define}
    \begin{thm}\label{sigma algebras}
        $\mathscr{A}\subset\mathscr{P}X$ is a $\sigma$-algebra if and only if the following hold:
        \begin{tight_enumerate}
            \item $X\in \mathscr{A}$,
            \item $\mathscr{A}$ is complement-stable,
            \item $\mathscr{A}$ is stable under countable union.%or $\mathscr{A}$ is stable under countable intersection.
        \end{tight_enumerate}
    \end{thm}
\subsection{Pullback by functions}
    \begin{thm}
        If $\mathscr{F}$ is a $\sigma$-ring (resp. $\sigma$-algebra) on $Y$, then for any function $f:X\rightarrow Y$,
        \[
            f^{-1}(\mathscr{F}):=%\{A\in\mathscr{P}X\,|\,\exists B\in \mathscr{F}(A=f^{-1}(B))\}
            \bigcup_{F\in\mathscr{F}}f^{-1}(F)
            =\mathfrak{i}\mathfrak{p}f(\mathscr{F})
        \]
        is a $\sigma$-ring (resp. $\sigma$-algebra) over $X$.

        In particular, if $X\subset Y$ and $\iota:X\rightarrow Y$ is the inclusion map, then
        \[
            \iota^{-1}(\mathscr{F})=%\{A\in\mathscr{P}X\,|\,\exists B\in\mathscr{F}(A=B\cap X)\}
            \bigcup_{F\in\mathscr{F}}X\cap F
            =:\mathscr{F}|_X
        \]
        is a $\sigma$-ring (resp. $\sigma$-algebra), called the trace of $\mathscr{F}$, on $X$.
    \end{thm}
    \begin{proof}
        Preimages preserve arbitrary set operations.
    \end{proof}
\section{Monotone classes and Dynkin systems}
\subsection{Monotone classes}
    \begin{define}
        $\mathscr{M}\subset\mathscr{P}X$ is called a \textit{monotone class} if $\mathscr{M}$ contains the limit of every monotone (increasing or decreasing) sequence of sets in $\mathscr{M}$.
    \end{define}
    \begin{thm}\label{characterization of sigma rings and sigma algebras via monotone classes}
        \[
            \text{$\sigma$-ring (resp. $\sigma$-algebra)}\Longleftrightarrow \text{ring (resp. algebra)}\wedge\text{monotone class}
        \]
    \end{thm}
    \begin{proof}
        $\Rightarrow$: Theorem \ref{property of sigma rings}. $\Leftarrow$: In a $\cup$-stable system, every countable union can be written as the limit of an increasing sequence.
    \end{proof}
\subsection{Dynkin systems}
    \begin{define}\label{definition of Dynkin system}
         A \textit{Dynkin system} on $X$ is a set $\mathscr{D}\subset\mathscr{P}X$ such that (c.f. Theorem \ref{sigma algebras})
         \begin{tight_enumerate}
             \item $X\in\mathscr{D}$,
             \item $\mathscr{D}$ is complement-stable,
             \item $\mathscr{D}$ is stable under countable disjoint union.
         \end{tight_enumerate}
    \end{define}
    \begin{thm}
        If $\mathscr{D}$ is a Dynkin system on $X$, then for any fixed $D_0\in \mathscr{D}$,
        \[
            \{D\in\mathscr{D}\,|\,D\cap D_0\in\mathscr{D}\}
        \]
        is again a Dynkin system on $X$.
    \end{thm}
    \begin{proof}
        If $D\in\mathscr{D}$ with $D\cap D_0\in\mathscr{D}$, then $D^c\in \mathscr{D}$ with $D^c\cap D_0=((D\cap D_0)\sqcup D_0^c)^c\in\mathscr{D}$.
    \end{proof}
    \begin{thm}\label{characterization of Dynkin systems via monotone classes}
        $\mathscr{D}\subset\mathscr{P}X$ is a Dynkin system if and only if the following holds:
        \begin{tight_enumerate}
            \item $X\in \mathscr{D}$.
            \item $(A,B\in\mathscr{D})\wedge (A\subset B)\Rightarrow B\setminus A\in\mathscr{D}$.
            \item $\mathscr{D}$ contains the limit of every increasing sequence.
        \end{tight_enumerate}
        In particular, a Dynkin system is necessarily a monotone class.
    \end{thm}
    \begin{proof}
        Necessity. A Dynkin system is $\sqcup$-stable by Theorem \ref{countable union implies binary union}. If $(A,B\in\mathscr{D})\wedge (A\subset B)$, then $A,B^c$ are disjoint and $B\setminus A=(A\sqcup B^c)^c\in\mathscr{D}$. With \textit{2}, the limit of every increasing sequence can be written as a countable disjoint union. Sufficiency. \textit{1} and \textit{2} implies complement-stability. If $A,B\in\mathscr{D}$ are disjoint, then $A\subset B^c$ and $A\sqcup B=(B^c\setminus A)^c\in\mathscr{D}$. Thus, every countable disjoint union can be written as the limit of an increasing sequence.
    \end{proof}
\subsection{Characterization of sigma-algebras}
    \begin{thm}\label{characterization of sigma-algebras via Dynkin systems and monotone classes}
        Let $\mathscr{A}\subset\mathscr{P}X$. The following are equivalent:
        \begin{tight_enumerate}
            \item $\mathscr{A}$is a $\sigma$-algebra.
            \item $\mathscr{A}$ is a $\cap$-stable Dynkin system.
            \item $\mathscr{A}$ is a $\cup$-stable Dynkin system.
            \item $\mathscr{A}$ is a $\setminus$-stable monotone class that contains $X$.
        \end{tight_enumerate}
    \end{thm}
    \begin{proof}
        \textit{2}$\vee$\textit{3}$\Rightarrow$\textit{1}. A $\cap$-stable or $\cup$-stable Dynkin system is, by Theorem \ref{characterization of algebras}, an algebra, and by Definition \ref{definition of sigma algebra} and Definition \ref{definition of Dynkin system}, a $\sigma$-algebra.
        \quad
        \textit{4}$\Rightarrow$\textit{2}. Theorem \ref{characterization of Dynkin systems via monotone classes}. Theorem \ref{setminus stable implies cap stable}.
    \end{proof}
\section{Generators}
    \begin{thm}
        The intersection of an arbitrary set of rings (resp. algebras, $\sigma$-rings, $\sigma$-algebras, monotone classes, Dynkin systems) on $X$ is also a ring (resp. algebra, $\sigma$-ring, $\sigma$-algebra, monotone class, Dynkin system) on $X$.
    \end{thm}
    \begin{define}
        Let $\mathscr{G}\subset\mathscr{P}X$. The sets
        \begin{align*}
            \mathfrak{r}(\mathscr{G})&:=\bigcap\,\{\text{rings in }X\text{ that contains }\mathscr{G}\text{ as a subset}\}\\
            \mathfrak{a}(\mathscr{G})&:=\bigcap\,\{\text{algebras in }X\text{ that contains }\mathscr{G}\text{ as a subset}\}\\
            \mathfrak{s}(\mathscr{G})&:=\bigcap\,\{\sigma\text{-rings in }X\text{ that contains }\mathscr{G}\text{ as a subset}\}\\
            \sigma(\mathscr{G})&:=\bigcap\,\{\sigma\text{-algebras in }X\text{ that contains }\mathscr{G}\text{ as a subset}\}\\
            \mathfrak{m}(\mathscr{G})&:=\bigcap\,\{\text{monotone classes in }X\text{ that contains }\mathscr{G}\text{ as a subset}\}\\
            \mathfrak{d}(\mathscr{G})&:=\bigcap\,\{\text{Dynkin systems in }X\text{ that contains }\mathscr{G}\text{ as a subset}\}
        \end{align*}
        are called resp. the ring, algebra, $\sigma$-ring, $\sigma$-algebra, monotone class, Dynkin system \textit{generated} by $\mathscr{G}$.
    \end{define}
    \begin{thm}\label{generators}
        \begin{tight_enumerate}
            \item If $\mathscr{R}$ is a ring, then $\mathfrak{s}(\mathscr{R})=\mathfrak{m}(\mathscr{R})$.
            \item If $\mathscr{A}$ is an algebra, then $\sigma(\mathscr{A})=\mathfrak{m}(\mathscr{A})$.
            \item If $\mathscr{G}$ is $\cap$-stable, then $\sigma(\mathscr{G})=\mathfrak{d}(\mathscr{G})$.
            \item If $\mathscr{G}$ is $\cap$-stable and $\setminus$-stable, then $\sigma(\mathscr{G})=\mathfrak{m}(\mathscr{G})$.
        \end{tight_enumerate}
    \end{thm}
    \begin{proof}
        
    \end{proof}
    \begin{thm}
        Let $\mathscr{G}\subset\mathscr{P}X$.
        \[
            \mathfrak{m}(\mathscr{G})=\mathfrak{s}(\mathscr{G})\Longleftrightarrow \mathfrak{r}(\mathscr{G})\subset\mathfrak{m}(\mathscr{G})
        \]
    \end{thm}
    \begin{proof}
        What is the analogue for $\sigma(\mathscr{G})$?
    \end{proof}
\subsection{Borel \texorpdfstring{$\sigma$}{sigma}-algebras}
    \begin{define}
        The $\sigma$-algebra generated by the open sets in a topological space $X$ is called the \textit{Borel $\sigma$-algebra} on $X$, denoted by $\mathscr{B}(X)$, whose elements are called \textit{Borel sets}.
    \end{define}
    It follows from Theorem \ref{sigma algebras} that \textit{$G_\delta$ sets and $F_\sigma$ sets are Borel sets}.
    \begin{thm}\label{generators of a Borel sigma algebra}
        Let $X$ be a topological space, with $\mathscr{C}$ the set of closed sets, $\mathscr{K}$ the set of compact sets.
        \begin{tight_enumerate}
            \item $\mathscr{B}(X)=\sigma(\mathscr{C})$.
            \item $\mathscr{B}(X)=\sigma(\mathscr{K})$ if $X$ is a $\sigma$-compact Hausdorff space.
        \end{tight_enumerate}
    \end{thm}
    \begin{proof}
        Recall that $\mathscr{K}\subset \mathscr{C}$ in Hausdorff space, and compactness is weakly hereditary. If $F$ is a closed set, then $F=F\cap X=F\cap (\bigcup_{n}K_n)=\bigcup_n(F\cap K_n)\in\sigma(\mathscr{K})$. Thus, $\mathscr{C}\subset\sigma(\mathscr{K})$.
    \end{proof}
    \begin{thm}
        $\mathscr{B}(\mathbb{R}^p)$ can be generated by any of the following that are closed under binary intersection:
        \begin{gather*}
            \rrbracket \mathbb{R}^p,\mathbb{R}^p\llbracket\,\,:=\left\{%\rrbracket a,b\llbracket\,\,:=\,
            \, ]a^1,b^1[\,\times\cdots\times\,]a^p,b^p[\,\big|\,(a,b\in\mathbb{R}^p)\wedge(a^i\leqslant b^i, \,i=1,\dots,p)\right\},\\
            \rrbracket \mathbb{Q}^p,\mathbb{Q}^p\llbracket\,\,:=\left\{%\rrbracket a,b\llbracket\,\,:=\,
            \, ]a^1,b^1[\,\times\cdots\times\,]a^p,b^p[\,\big|\,(a,b\in\mathbb{Q}^p)\wedge(a^i\leqslant b^i, \,i=1,\dots,p)\right\},\\
            \llbracket \mathbb{R}^p,\mathbb{R}^p\llbracket\,,\quad \llbracket \mathbb{Q}^p,\mathbb{Q}^p\llbracket\,,\\
            \rrbracket \mathbb{R}^p,\mathbb{R}^p\rrbracket\,,\quad \rrbracket \mathbb{Q}^p,\mathbb{Q}^p\rrbracket\,,\\
            \{\varnothing\}\cup\llbracket \mathbb{R}^p,\mathbb{R}^p\rrbracket\,,\quad \llbracket \mathbb{Q}^p,\mathbb{Q}^p\rrbracket\cup\{\varnothing\}\,\,
        \end{gather*}
    \end{thm}
    \begin{proof}
        I show that $\mathscr{O}\subset\sigma(\rrbracket \mathbb{Q}^p,\mathbb{Q}^p\llbracket)$. Let $U$ be an open set. $\rrbracket \mathbb{Q}^p,\mathbb{Q}^p\llbracket$ is countable, and a fortiori, $\mathscr{I}=\{I\in\,\,\rrbracket \mathbb{Q}^p,\mathbb{Q}^p\llbracket\,|\, I\subset U\}$ is countable. It suffices to show that $\mathscr{I}$ covers $U$. Indeed, for all $p\in U$, there is an open ball $B_p(r)\subset U$, and hence an open interval $I$ with rational endpoints such that $p\in I\subset B_p(r)\subset U$. Together with $\rrbracket \mathbb{Q}^p,\mathbb{Q}^p\llbracket\,\,\subset\,\,\rrbracket \mathbb{R}^p,\mathbb{R}^p\llbracket\,\,\subset \mathscr{O}$ I have $\sigma(\mathscr{O})=\sigma(\rrbracket \mathbb{R}^p,\mathbb{R}^p\llbracket)=\sigma(\rrbracket \mathbb{Q}^p,\mathbb{Q}^p\llbracket)$.
        
        Moreover, every open interval with rational endpoints can be written as a countable union of half-open  intervals with rational endpoints, as well as a countable union of closed intervals with rational endpoints. Conversely, every half-open or closed interval with rational endpoints can be written as a countable intersection of open intervals with rational endpoints.
    \end{proof}
    \begin{crl}
        $\mathscr{B}(\mathbb{R}^p)=\mathfrak{d}(\mathscr{O})=\mathfrak{d}(\mathscr{C})=\mathfrak{d}(\mathscr{K})$.
    \end{crl}
\subsection{Pullback of generated \texorpdfstring{$\sigma$}{sigma}-algebras}
    \begin{thm}\label{sigma commutes with preimage}
        Let $f:X\rightarrow Y$ be a function and $\mathscr{G}\subset\mathscr{P}(Y)$.
        \[f^{-1}(\sigma(\mathscr{G}))=\sigma(f^{-1}(\mathscr{G}))\]
    \end{thm}
    \begin{proof}
        $f^{-1}(\sigma(\mathscr{G}))$ is a $\sigma$-algebra that contains $f^{-1}(\mathscr{G})$, so $f^{-1}(\sigma(\mathscr{G}))\supset\sigma(f^{-1}(\mathscr{G}))$. Conversely, it can be verified that $\Sigma:=\{B\subset Y\,|\,f^{-1}(B)\in\sigma(f^{-1}(\mathscr{G}))\}$ is a $\sigma$-algebra in $Y$, and since $\mathscr{G}\subset \Sigma$, I have $\sigma(\mathscr{G})\subset \Sigma$ and $f^{-1}(\sigma(\mathscr{G}))\subset f^{-1}(\Sigma)\subset\sigma(f^{-1}(\mathscr{G}))$. Nothing is actually mysterious about $\Sigma$, for $\Sigma=\mathfrak{p}\mathfrak{p}f(\sigma(f^{-1}(\mathscr{G})))$ and recall $\mathfrak{i}f(A)\subset B\Longleftrightarrow A\subset \mathfrak{p}f(B)$ that holds for all functions $f$.
    \end{proof}
    \begin{crl}
        Let $X\subset Y$ and $\mathscr{A}$ a $\sigma$-algebra on $Y$. If $\mathscr{A}=\sigma(\mathscr{G})$, then $\mathscr{A}|_X=\sigma(\mathscr{G}|_X)$. %generates the trace $\sigma$-algebra $\mathscr{A}|_X$ over $X$.
    \end{crl}
    If $A\in\mathscr{A}|_X$ and $X\in\mathscr{A}$, then $A\in\mathscr{A}$.
\subsection{Borel \texorpdfstring{$\sigma$}{sigma}-algebras on subspaces}
    \begin{crl}\label{Borel Sigma Algebra of Topological Subspace}
        Let $Y$ be a topological space. For $X$ any subspace of $Y$,
        \[\mathscr{B}(X)=\mathscr{B}(Y)|_X\]
    \end{crl}
\section{Semirings}
    \begin{define}
        A \textit{semiring} on $X$ is a set $\mathscr{H}\subset\mathscr{P}X$ such that
        \begin{tight_enumerate}
            \item $\varnothing\in\mathscr{H}$.
            \item $\mathscr{H}$ is closed under $\cap$.
            \item $(S_1,S_2\in\mathscr{H})\Longrightarrow(S_2\setminus S_1\text{ equals a finite disjoint union of some elements of }\mathscr{H})$.
        \end{tight_enumerate}
    \end{define}
    A ring is necessarily a semiring.
    \begin{thm}
        If $A,B_1,\dots,B_n$ are elements of a semiring $\mathscr{H}$, then there exists finitely many disjoint sets $C_1,\dots, C_m\in\mathscr{H}$ such that\footnote{Square cups are used to denote the union of \textit{disjoint} sets}
        \[
            A\setminus\bigcup_{i=1}^n B_i=\bigsqcup_{j=1}^mC_j
        \]
    \end{thm}
    \begin{proof}
        Induction on $n\in\mathbb{N}$.
    \end{proof}
\subsection{Product of semirings}
    \begin{thm}
        If $\mathscr{A}$ and $\mathscr{B}$ are semirngs over $X$ and $Y$ respectively, then 
            \[\mathscr{A}* \mathscr{B}:=\{A\times B\,|\,A\in\mathscr{A},B\in\mathscr{B}\}\] 
        is a semirng over $X\times Y$.
    \end{thm}
    \begin{proof}
        
    \end{proof}
    In fact, $\mathscr{A}* \mathscr{B}$ is also a product of $\mathscr{A}$ and $\mathscr{B}$ in \cat{Set}. The Cartesian product $\mathscr{A}\times \mathscr{B}$ is not a natural consideration in this context because it is not a subset of $\mathscr{P}(X\times Y)$.
    \begin{crl}
        $\llbracket \mathbb{R}^p,\mathbb{R}^p\llbracket$ and $\llbracket \mathbb{Q}^p,\mathbb{Q}^p\llbracket$ are semirngs over $\mathbb{R}^p$ for $p\in\mathbb{N}$.
    \end{crl}
\subsection{The ring generated by a semiring}
    \begin{thm}\label{ring generated by semiring}
        If $\mathscr{H}$ is a semiring, then $\mathfrak{r}(\mathscr{H})=\{\text{finite disjoint unions of sets in }\mathscr{H}\}$.
    \end{thm}
    \begin{proof}
        
    \end{proof}
\subsection{Countable union in a semiring}
    \begin{thm}\label{every countable union is a disjoint union in a semiring}
        A countable union of sets in a semiring $\mathscr{H}$ can be written as a countable disjoint union of sets in $\mathscr{H}$:
        \[
            \bigcup_{n=1}^\infty A_n=\bigsqcup_{n=1}^\infty A_n'
        \]
        such that for every $n\in\mathbb{N}$ there exists $m\in\mathbb{N}$ such that $A_n'\subset A_m$.
    \end{thm}
    \begin{proof}
        Consider $\mathfrak{r}(\mathscr{H})$ and use Corollary \ref{every countable union is a disjoint union in a ring}.
    \end{proof}
    % Thus,
    % \[
    %     \mathfrak{r}(\llbracket \mathbb{Q}^p,\mathbb{Q}^p\llbracket)=\left\{\bigcup_{i=1}^n I_i\,\bigg|\,I_1,I_2,\dots,I_n\in\,\,\llbracket \mathbb{Q}^p,\mathbb{Q}^p\llbracket\text{ are disjoint}, n\geqslant 0\right\}
    % \]
\subsection{Borel \texorpdfstring{$\sigma$}{sigma}-algebra on \texorpdfstring{$\mathbb{R}^p$}{Rp}}
    \begin{thm}
        $\mathscr{B}(\mathbb{R}^p)=\mathfrak{d}(\mathfrak{r}(\llbracket \mathbb{Q}^p,\mathbb{Q}^p\llbracket))=\mathfrak{m}(\mathfrak{r}(\llbracket \mathbb{Q}^p,\mathbb{Q}^p\llbracket))$.
    \end{thm}
    \begin{proof}
        $\mathscr{B}(\mathbb{R}^p)=\mathfrak{s}(\mathfrak{r}(\llbracket \mathbb{Q}^p,\mathbb{Q}^p\llbracket))=\mathfrak{m}(\mathfrak{r}(\llbracket \mathbb{Q}^p,\mathbb{Q}^p\llbracket))$ by Theorem \ref{generators}.
    \end{proof}
    \begin{crl}
        $\mathscr{B}(\mathbb{R}^p)=\mathfrak{m}(\mathscr{O})=\mathfrak{m}(\mathscr{C})=\mathfrak{m}(\mathscr{K})$.
    \end{crl}
    \begin{proof}
        
    \end{proof}
\section{The Extended Real Line}
    The \textit{extended real line} $\er:=\mathbb{R}\cup\{-\infty,\infty\}$ extends the total order on $\mathbb{R}$ to include $-\infty<x<\infty$ for all $x\in\mathbb{R}$. 
    
    $\er$ and $[0,1]$ are order isomorphic tosets. Consequently,
    \begin{enumerate}
        \item $\er$ is homeomorphic to $[0,1]$ in the order topology. Thus, $\er$ is a compactification of $\mathbb{R}$.
        \item $\er$ is a linear continuum that has $\infty$ as the greatest element and $-\infty$ as the least element. Thus, $\er$ is a complete lattice.
    \end{enumerate}
    \begin{thm}
        Every increasing sequence $(a_n)\in\er^\mathbb{N}$ converges to $\sup_{n\in\mathbb{N}}a_n$, i.e.
        \[
            \exists \lim_{n\to\infty} a_n=\sup_{n\in\mathbb{N}}a_n\in\overline{\mathbb{R}}
        \]
        In this case, we write $a_n\nearrow \sup_{n\in\mathbb{N}}a_n$ as $n\to \infty$.
    \end{thm}
    Thus, limit superior and limit inferior of any sequence in $\er$ exist, and the limit superior (resp. inferior) is the largest (resp. smallest) limit point of the sequence.
    % The Limes inferior and the Limes superior of $(a_n)$ are defined as:
    % \[\varliminf_{n\to \infty}a_n:=\lim_{n\to\infty}\inf_{k\geqslant n}a_k=\sup_{n\in\mathbb{N}}\inf_{k\geqslant n}a_k,\quad\quad\varlimsup_{n\to \infty}a_n:=\lim_{n\to\infty}\sup_{k\geqslant n}a_k=\inf_{n\in\mathbb{N}}\sup_{k\geqslant n}a_k\]
    \begin{thm}
        Let $(a_n)\in \mathbb{R}^\mathbb{N}$. $\exists \lim\limits_{n\to \infty}a_n\in\er\Longleftrightarrow \varliminf\limits_{n\to\infty}a_n=\varlimsup\limits_{n\to \infty}a_n$.
    \end{thm}
    % \begin{thm}
    %     Let $(s_n)\in{\mathbb{R}}^\mathbb{N}$ and $(t_n)\in{\overline{\mathbb{R}}^*}^\mathbb{N}$. If $s_n\to s$ and $t_n\to t$, where $t\neq 0$, then
    %     \[\frac{s_n}{t_n}\to \frac{s}{t}\]
    % \end{thm}
    \begin{define}
        $\mathbb{R}_+:=\{x\in\mathbb{R}\,|\,x\geqslant 0\}$. $\er_+:=\{x\in\er\,|\,x\geqslant 0\}=\mathbb{R}_+\cup\{\infty\}$.
    \end{define}
\subsection{Extension of the arithmetic}
    \begin{define}
        Let $+:\er\cup\{\ud\}\times\er\cup\{\ud\}\rightarrow \er\cup\{\ud\}$ be defined by
        \begin{tight_enumerate}
            \item $+|_{\mathbb{R}\times\mathbb{R}}$ is the usual addition on the real line.
            \item $\infty+a=a+\infty=\infty$ for all $a\in \,]\!-\!\infty,\infty]$.
            \item $(-\infty)+a=a+(-\infty)=-\infty$ for all $a\in [-\infty,\infty[$.
            \item $\infty+(-\infty)=(-\infty)+\infty=\ud$
            \item $\ud + x=x+\ud=\ud$ for all $x\in \er\cup\{\ud\}$.
        \end{tight_enumerate}
        Thus, $+$ is a commutative and associative binary operation on $\er\cup\{\ud\}$.
    \end{define}
    \begin{define}
        Let $\cdot :\er\cup\{\ud\}\times\er\cup\{\ud\}\rightarrow \er\cup\{\ud\}$ be defined by
        \begin{tight_enumerate}
            \item $\cdot|_{\mathbb{R}\times\mathbb{R}}$ is the usual multiplication on the real line.
            \item $(\pm\infty)\cdot a=a\cdot (\pm\infty)=\pm\infty$ for all $a\in\,]0,\infty]$.
            \item $\pm\infty\cdot 0=0\cdot \pm\infty=0$.
            \item $(\pm\infty)\cdot a=a\cdot (\pm\infty)=\mp\infty$ for all $a\in [-\infty,0[$.
            \item $x\cdot \ud=\ud\cdot x=\ud$ for all $x\in\er\cup\{\ud\}\setminus\{0\}$.
            \item $\ud\cdot 0=0\cdot \ud=0$.
        \end{tight_enumerate}
        Thus, $\cdot$ is a commutative and associative binary operation on $\er\cup\{\ud\}$, and for all $x,y\in \er\cup\{\ud\}$,
        \[
            x\cdot 0=0\cdot x=0
        \]
        and
        \[
            x\cdot y=0\Longrightarrow x=0\vee y=0
        \]
    \end{define}
    $\cdot$ is not distributive over $+$ since $\infty\cdot (-1+2)=\infty\cdot 1=\infty\neq \ud=\infty\cdot (-1)+\infty\cdot 2$.
    \begin{thm}
        Let $a\in\mathbb{R}$ and $x,y\in \er\cup\{\ud\}$.
        \[
            a\cdot (x+y)=a\cdot x+a\cdot y
        \]
    \end{thm}
    \begin{define}
        For $x\in\er$, let $-x:=(-1)\cdot x$.
    \end{define}
    \begin{thm}
        $-(a+b)=(-a)+(-b)$ for all $a,b\in\er\cup\{\ud\}$.
    \end{thm}
\subsection{The Borel \texorpdfstring{$\sigma$}{sigma}-algebra}
    \begin{thm}\label{Borel sigma algebra on R}
        $\mathscr{B}(\mathbb{R})=\sigma((\mathbb{Q},\infty))=\sigma([\mathbb{Q},\infty))=\sigma((-\infty,\mathbb{Q}))=\sigma((-\infty,\mathbb{Q}])$.
    \end{thm}
    \begin{thm}\label{Borel sigma algebra on ER}
        $\mathscr{B}(\overline{\mathbb{R}})=\sigma((\mathbb{Q},\infty])=\sigma([\mathbb{Q},\infty])=\sigma([-\infty,\mathbb{Q}))=\sigma([-\infty,\mathbb{Q}])$.
    \end{thm}
    \begin{crl}
        $\mathscr{B}(\overline{\mathbb{R}})=\{B\cup S\,|\,B\in\mathscr{B}(\mathbb{R}),\,S\subset \{-\infty,\infty\}\}$.
    \end{crl}
    \begin{proof}
        By Corollary \ref{Borel Sigma Algebra of Topological Subspace} and $\mathbb{R}\in \mathscr{B}(\overline{\mathbb{R}})$, we have $\mathfrak{B}={\overline{\mathfrak{B}}}|_{\mathbb{R}}\subset \overline{\mathfrak{B}}$.
    \end{proof}
    \begin{thm}\label{Borel sigma algebra of a countable subset}
        If $C\subset\er$ is a countable subset of $\er$, then
        \[
            \mathscr{B}(C)=\sigma(\{\{x\}\,|\,x\in C\})
        \]
    \end{thm}
\subsection{Commutation of suprema}
    \begin{thm}\label{commutation of suprema}
        Let $(a_{m,n}\in \er)_{(m,n)\in\mathbb{N}^2}$ a double sequence in $\er$.
        \[
            \sup_{m\in\mathbb{N}}\sup_{n\in\mathbb{N}}a_{m,n}=\sup_{(m,n)\in\mathbb{N}^2}a_{m,n}=\sup_{n\in\mathbb{N}}\sup_{m\in\mathbb{N}}a_{m,n}
        \]
    \end{thm}
\chapter{Measures}
    \begin{define}
        A \textit{set function over }$X$ is a function $\varphi:\mathscr{C}\rightarrow [0,\infty]$ with $\mathscr{C}\subset \mathscr{P}X$.
    \end{define}
    \begin{define}
        A set function $\varphi:\mathscr{C}\rightarrow [0,\infty]$ over $X$ is
        \begin{tight_enumerate}
            \item \textit{isotone} if $A\subset B\Rightarrow \varphi(A)\leqslant \varphi(B)$ for all $A,B\in\mathscr{C}$.
            \item \textit{finite} if $\varphi(A)<\infty$ for all $A\in\mathscr{C}$.
            \item $\sigma$-\textit{finite}\footnote{Admittedly, this is somewhat artificial since $X$ may not be in $\mathscr{H}$, i.e., $\varphi(X)$ is undefined.} if $\{\varphi<\infty\}$ contains a countable cover of $X$.
            \item \textit{trivial} if $\varphi(A)=0$ for all $A\in\mathscr{C}$.
        \end{tight_enumerate}
    \end{define}
    \begin{thm}
        Let $\varphi:\mathscr{C}\rightarrow [0,\infty]$ be an isotone $\sigma$-finite set function over $X$.
    
        If $\mathscr{C}$ is $\cup$-stable, then there exists an increasing sequence $(A_n)_{n\in\mathbb{N}}$ of sets in $\mathscr{C}$ such that
        \begin{gather*}
            A_n\uparrow X\\
            \varphi(A_n)<\infty\quad\forall n\in\mathbb{N}
        \end{gather*}

        If $\mathscr{C}$ is a semiring, then $\{\varphi<\infty\}$ contains a countable partition of $X$. 
    \end{thm}
    \begin{proof}
        Theorem \ref{every countable union is a disjoint union in a semiring}.
    \end{proof}
    \begin{define}
        Let $\varphi:\mathscr{C}\rightarrow [0,\infty]$ be a set function with $\mathscr{C}$ complement-stable. 
        
        A $\varphi$\textit{-null set} is a set $N\in\mathscr{C}$ with $\varphi(N)=0$ and a $\varphi$\textit{-full set} is a set $F\in\mathscr{C}$ with $\varphi(F^c)=0$.
    \end{define}
    % The set of $\mu$-null sets is denoted by $\mathscr{N}_\mu:=\mu^{-1}(\{0\})$.
\section{Contents}
    \begin{define}
        A \textit{content} is a set function $\mu:\mathscr{C}\rightarrow [0,\infty]$ satisfying
        \begin{tight_enumerate}
            \item $\varnothing\in\mathscr{C}$ and $\mu(\varnothing)=0$,
            \item \textit{finite additivity}, i.e., for all $n\in\mathbb{N}$,
            \[
                A_1,\dots,A_n\in\mathscr{C}\text{ disjoint and }\bigsqcup_{i=1}^nA_i\in\mathscr{C}\Longrightarrow \mu\left(\bigsqcup_{i=1}^nA_i\right)=\sum_{i=1}^n\mu(A_i)
            \]
        \end{tight_enumerate}
        % $\mu(\varnothing)=0$ is not redundant because $\mu(\varnothing)=\infty$ also satisfies $\mu(\varnothing)=\mu(\varnothing)+\mu(\varnothing)$.
    \end{define}
\subsection{Isotonicity of contents on a semiring}
    \begin{thm}\label{monotonicity of a content on a semiring}
        If $\mu$ is a content on a semiring $\mathscr{H}$, then for all $A,B\in\mathscr{H}$
        \[A\subset B\Longrightarrow \mu(A)\leqslant\mu(B)\]
    \end{thm}
    \begin{proof}
        $B\setminus A$ may not be in $\mathscr{H}$, but can be written as a finite disjoint union $\bigsqcup_{i=1}^n H_i$ of sets in $\mathscr{H}$ since $\mathscr{H}$ is a semiring. Then $\mu(B)=\mu(A)+\mu(H_1)+\cdots+\mu(H_n)$.
    \end{proof}
\subsection{The first extension theorem}
    \begin{thm}\label{first extension theorem}
        For every content $\mu$ on a semiring $\mathscr{H}$, there exists a unique extension of $\mu$ to a content $\overline{\mu}$ on $\mathfrak{r}(\mathscr{H})$, and for all $A\in \mathfrak{r}(\mathscr{H})$ there exist $S_1,\dots,S_n\in\mathscr{H}$ such that
        \[
          \overline{\mu}(A)=\mu(S_1)+\cdots+\mu(S_n)  
        \]
        
        Furthermore, 
        \begin{tight_enumerate}
            \item $\overline{\mu}$ is a premeasure if and only if $\mu$ is a premeasure.
            \item $\overline{\mu}$ is $\sigma$-finite if and only if $\mu$ is $\sigma$-finite.
        \end{tight_enumerate}
    \end{thm}
    \begin{proof}

    \end{proof}
\subsection{Contents on a ring}
    \begin{thm}\label{properties of contents}
        If $\mu$ is a content on a ring $\mathscr{R}$, then
        \begin{tight_enumerate}
            \item for all $A,B\in\mathscr{R}$,
            \[
                \mu(A\setminus B)=0\wedge \mu(B\setminus A)=0 \Longrightarrow \mu(A)=\mu(B)
            \]
            \item for all $A,B\in\mathscr{R}$ such that $A\subset B$ and $\mu(A)<\infty$,
            \[\mu(B\setminus A)=\mu(B)-\mu(A)\]
            \item for all $A,B\in\mathscr{R}$,
            \[\mu(A)+\mu(B)=\mu(A\cup B)+\mu(A\cap B)\]
            \item \label{subadditivity of contents on a ring} for all $A_1,\dots,A_n\in\mathscr{R}$, where $n\in\mathbb{N}$,
            \[\mu(A_1\cup \cdots\cup A_n)\leqslant \mu(A_1)+\cdots+\mu(A_n)\]
            \item for all $(A_n)\in\mathscr{R}^\mathbb{N}$ disjoint,
            \[
                \bigcup_{n=1}^\infty A_n\subset B\in\mathscr{R}\Longrightarrow \sum_{n=1}^\infty \mu(A_n)\leqslant \mu(B)
            \]
            and in particular,
            \[
                \bigcup_{n=1}^\infty A_n\in\mathscr{R}\Longrightarrow \sum_{n=1}^\infty \mu(A_n)\leqslant \mu\left(\bigcup_{n=1}^\infty A_n\right)
            \]
        \end{tight_enumerate}
    \end{thm}
    \begin{proof}
        1. $\mu(A)=\mu(A\cap B)+\mu(A\setminus B)=\mu(A\cap B)=\mu(B\cap A)+\mu(B\setminus A)=\mu(B)$. 2.
    \end{proof}
    \begin{crl}
        If $\mu$ is a content on a ring $\mathscr{R}$, then a finite union of $\mu$-null sets is $\mu$-null.
    \end{crl}
    \begin{crl}\label{subsets of a finite union of null sets of a content}
        If $\mu$ is a content on a semiring $\mathscr{H}$ and $N_1,\dots,N_n\in\mathscr{N}_\mu$, then for all $A\in\mathscr{H}$,
        \[
            A\subset N_1\cup\cdots\cup N_n\Longrightarrow A\in\mathscr{N}_\mu
        \]
    \end{crl}
    \begin{proof}
        Let $\overline{\mu}$ be the extension to $\mathfrak{r}(\mathscr{H})$. Clearly, $\mathscr{N}_\mu\subset \mathscr{N}_{\overline{\mu}}$ and $\mu(A)=\overline{\mu}(A)$.
    \end{proof}
\section{Premeasures and measures}
    \begin{define}
        A \textit{premeasure} is a set function $\mu:\mathscr{C}\rightarrow [0,\infty]$ satisfying
        \begin{tight_enumerate}
            \item $\varnothing\in\mathscr{C}$ and $\mu(\varnothing)=0$,
            \item \textit{$\sigma$-additivity}, i.e.\footnote{Since the union $\bigsqcup_{n=1}^\infty A_n$ is invariant under permutations of the sequence $A_\bullet$, so should be the series $\sum_{n=1}^\infty\mu(A_n)$. Indeed, recall that absolute convergence implies unconditional convergence.},
            \[
                A_\bullet\in\mathscr{C}^\mathbb{N}\text{ disjoint and }\bigsqcup_{n=1}^\infty A_n\in\mathscr{C}\Longrightarrow \mu\left(\bigsqcup_{n=1}^\infty A_n\right)=\sum_{n=1}^\infty\mu(A_n)
            \]
        \end{tight_enumerate}
        
        A \textit{measure} is a premeasure whose domain is a $\sigma$-algebra. 
        
        A \textit{measure space} is a triple $(E,\mathscr{E},\mu)$ where $\mathscr{E}$ is a $\sigma$-algebra on $E$ and $\mu$ a measure on $\mathscr{E}$; if $\mu(E)=1$ then $\mu$ is called a \textit{probability measure} and $(E,\mathscr{E},\mu)$ a \textit{probability space}.
    \end{define}
    % Both finite additivity and $\sigma$-additivity are \textit{hereditary}, i.e., any restriction of a finite additive (resp. $\sigma$-additive) function is also finite additive (resp. $\sigma$-additive).
    \begin{thm}
        A premeasure is necessarily a content.
    \end{thm}
    \begin{proof}
        $\mu(A_1\sqcup \cdots\sqcup A_n)=\mu(A_1\sqcup\cdots\sqcup A_n\sqcup\varnothing\sqcup\varnothing\sqcup\cdots)=\mu(A_1)+\cdots+\mu(A_n)$ since $\mu(\varnothing)=0$.
    \end{proof}
\subsection{The counting measure}
    \begin{thm}
        $(E,\mathscr{P}E,\#)$ is a measure space for any set $E$, where $\#$ is the {counting measure}:
        \[
            \#:\mathscr{P}E\rightarrow [0,\infty],\,\,A\mapsto \begin{cases}
                \text{the number of elements of }A & \text{if $A$ is a finite set},\\
                \infty &\text{otherwise}.
            \end{cases}
        \]
    \end{thm}
\subsection{\texorpdfstring{$\sigma$}{sigma}-subadditivity}
    \begin{thm}
        If $\mu$ is a premeasure on a ring $\mathscr{R}$, then for all $A\in\mathscr{R}$ and $(A_n)\in\mathscr{R}^\mathbb{N}$,
        \[
            A\subset \bigcup_{n=1}^\infty A_n\Longrightarrow\mu(A)\leqslant \sum_{n=1}^\infty \mu(A_n)
        \]
        In particular, if $\mu$ is a measure on a $\sigma$-algebra $\mathscr{A}$, then for all $(A_n)\in\mathscr{A}^\mathbb{N}$,
        \[
            \mu\left(\bigcup_{n=1}^\infty A_n\right)\leqslant \sum_{n=1}^\infty \mu(A_n)
        \]
    \end{thm}
    \begin{proof}
        By Corollary \ref{every countable union is a disjoint union in a ring}, $A=A\cap \bigcup_{n=1}^\infty A_n=\bigsqcup_{n=1}^\infty A\cap A_n'$ with $A\cap A_n'\in\mathscr{R}$ and $A\cap A_n'\subset A_n$.
    \end{proof}
    \begin{crl}\label{measurable subsets of a countable union of null sets is null}
        If $\mu$ is a premeasure on a semiring $\mathscr{H}$ and $(N_n)\in\mathscr{N}_\mu^\mathbb{N}$, then for all $A\in\mathscr{H}$,
        \[
            A\subset \bigcup_{n=1}^\infty N_n \Longrightarrow A\in\mathscr{N}_\mu
        \]
        In particular, for $\mu$ any measure, 
        \begin{tight_enumerate}
            \item a countable union of $\mu$-null sets is a $\mu$-null set,
            \item a countable intersection of $\mu$-full sets is a $\mu$-full set.
        \end{tight_enumerate}
        % \[
        %     (N_n)\in\mathscr{N}_{\mu}^\mathbb{N}\Longrightarrow \bigcup_{n=1}^\infty N_n\in \mathscr{N}_\mu
        % \]
    \end{crl}
    \begin{proof}
        See proof of Corollary \ref{subsets of a finite union of null sets of a content}.
    \end{proof}
\subsection{Characterization of premeasures by monotone convergence of contents}
    \begin{thm}
    \label{characterization of premeasures by monotone convergence}
        Let $\mu$ be a content on a ring $\mathscr{R}$.
        \[
            \begin{tikzcd}[column sep={20pt}]
                \text{$\mu$ is a premeasure } \ar[r, Leftrightarrow] \ar[d, Rightarrow, shorten <= 0.13cm, shorten >= 0.14cm] & \begin{tabular}{c}
                    $\forall A_\bullet\in\mathscr{R}^\mathbb{N}$   \\
                    $(A_\bullet\uparrow A\in\mathscr{R})\Longrightarrow \mu(A_\bullet)\uparrow\mu(A)$
                \end{tabular}
                \ar[d, Rightarrow]\\
                \begin{tabular}{c}
                      $\forall A_\bullet\in\mathscr{R}^\mathbb{N}$ \\
                      $(A_\bullet\downarrow \varnothing)\wedge(\mu(A_1)<\infty)\Longrightarrow \mu(A_\bullet)\downarrow 0$
                \end{tabular} \ar[r, Leftrightarrow]  & \begin{tabular}{c}
                      $\forall A_\bullet\in\mathscr{R}^\mathbb{N}$ \\
                      $(A_\bullet\downarrow A\in\mathscr{R})\wedge(\mu(A_1)<\infty)\Longrightarrow \mu(A_\bullet)\downarrow\mu(A)$
                \end{tabular}
            \end{tikzcd}
        \]
        
        If $\mu$ is finite, then all four statements are equivalent.
    \end{thm}
    \begin{proof}
        
    \end{proof}
    Theorem \ref{characterization of premeasures by monotone convergence} does not hold for contents on semirings.
\subsection{s-finite measures}
    \begin{thm}\label{mu(A cap ) is a measure}
        If $\mu:\mathscr{A}\rightarrow [0,\infty]$ is a measure, then 
        \[
            \mu_{\cap B}:\mathscr{A}\rightarrow [0,\infty],\,\,A\mapsto \mu(A\cap B)
        \]
        is a measure for every fixed $B\in\mathscr{A}$.
    \end{thm}
    \begin{proof}
        Since $\mathscr{A}$ is $\cap$-stable, $\mu_{\cap B}$ is well-defined.
        For $A_\bullet\in\mathscr{A}^\mathbb{N}$ holds $\mu_{\cap B}(\bigsqcup_{n=1}^\infty A_n)=\mu(\bigsqcup_{n=1}^\infty A_n\cap B)=\sum_{n=1}^\infty\mu(A_n\cap B)=\sum_{n=1}^\infty\mu_{\cap B}(A_n)$.
    \end{proof}
    \begin{thm}
        If $(\mu_n)_{n\in\mathbb{N}}$ is a sequence of measures on a $\sigma$-algebra $\mathscr{A}$, then
        \[
            \sum_{n=1}^\infty \mu_n
        \]
        is also a measure on $\mathscr{A}$.
    \end{thm}
    \begin{define}
        A measure $\mu:\mathscr{A}\rightarrow [0,\infty]$ is \textit{$\Sigma$-finite} if there exists a sequence $(\mu_n)_{n\in\mathbb{N}}$ of finite measures on $\mathscr{A}$ such that
        \[
            \mu=\sum_{n=1}^\infty \mu_n
        \]
    \end{define}
\subsection{Complete measure spaces}
    \begin{define}
        A measure space $(X,\mathscr{A},\mu)$ is \textit{complete} if
        \[N\in\mathscr{N}_\mu\Longrightarrow \mathscr{P}N\subset \mathscr{A}\]
        The \textit{completion} of a measure space is its smallest complete extension.
    \end{define}
    \begin{thm}
        For $(X,\mathscr{A},\mu)$ any measure space,
        \[\overline{\mathscr{A}}:=\{A\cup N'\,|\,A\in\mathscr{A}\wedge (N'\text{ is a subset of a }\mu\text{-null set})\}\]
        is a $\sigma$-algebra on $X$, and there exists a unique content $\overline{\mu}$ on $\overline{\mathscr{A}}$ that extends $\mu$.
        \[
            \begin{tikzcd}[column sep = 60pt]
                \overline{\mathscr{A}} \ar[r, "\overline{\mu}", "\exists!\,\text{content}"'] & {[0,\infty]}\\
                \mathscr{A} \ar[ur, "\mu"'] \ar[u, "\iota"] &
            \end{tikzcd}
        \]
        
        The unique content $\overline{\mu}$ is a measure on $\overline{\mathscr{A}}$, and $(X,\overline{\mathscr{A}},\overline{\mu})$ is the completion of $(X,\mathscr{A},\mu)$.
    \end{thm}
    \begin{proof}
        That $\overline{\mathscr{A}}$ is a $\sigma$-algebra follows from Corollary \ref{measurable subsets of a countable union of null sets is null}. If $\overline{\mu}$ is a content on $\overline{\mathscr{A}}$ that extends $\mu$, then by Theorem \ref{properties of contents}.\ref{subadditivity of contents on a ring} and  Theorem \ref{monotonicity of a content on a semiring}, for all $A\in\mathscr{A}$ and $N'$ a subset of a $\mu$-null set, \begin{equation}\label{completion of a measure}
            \overline{\mu}(A\cup N')=\mu(A)
        \end{equation}
        Thus, the extension $\overline{\mu}$ is unique if it exists at all, i.e., if $\overline{\mu}$ is well-defined by \eqref{completion of a measure}. In fact, if $A_1\cup N_1'=A_2\cup N_2'$, where $N_1'\subset N_1\in\mathscr{N}_\mu$ and $N_2'\subset N_2\in\mathscr{N}_\mu$,
        %where $A_1,A_2\in\mathscr{A}$ and $N_1,N_2$ are subsets of $\mu$-null sets, 
        then by Theorem \ref{monotonicity of a content on a semiring} and Theorem \ref{properties of contents}.\ref{subadditivity of contents on a ring},
        %$\mathscr{A}\ni A_1\setminus A_2\subset N_2'$ is a subset of a $\mu$-null set, and hence a $\mu$-null set, that is, $\mu(A_1)=\mu(A_1\cap A_2)=\mu(A_2)$.
        $\mu(A_1)\leqslant\mu(A_2\cup N_2)\leqslant \mu(A_2)+\mu(N_2)=\mu(A_2)$. By symmetry, $\mu(A_1)=\mu(A_2)$. It then suffices to verify that $\overline{\mu}$ defined by \eqref{completion of a measure} is a content. Indeed, $\overline{\mu}$ is a measure as $\sigma$-additivity of $\overline{\mu}$ follows immediately from that of $\mu$. It is clear that $(X,\overline{\mathscr{A}},\overline{\mu})$ is the smallest complete extension.
    \end{proof}
    \begin{thm}
        If $(X,\overline{\mathscr{A}},\overline{\mu})$ is the completion of $(X,\mathscr{A},\mu)$, then
        \[
            \overline{\mathscr{A}}=\{B\in\mathscr{P}X\,|\,\exists A_1,A_2\in\mathscr{A}(A_1\subset B\subset A_2)\wedge(\mu(A_2\setminus A_1)=0)\}
        \]
    \end{thm}
\section{Extension theorems}
\subsection{The uniqueness theorem}
    \begin{thm}\label{uniqueness theorem}
        Let $\varphi:\mathscr{G}\rightarrow [0,\infty]$ be a set function over $E$. If
        \begin{tight_enumerate}
            \item $\mathscr{G}$ is $\cap$-stable
            \item $\varphi$ is $\sigma$-finite% i.e., $\{\varphi<\infty\}$ contains a countable cover of $X$
        \end{tight_enumerate}
        then there exists at most one measure $\mu$ on $\sigma(\mathscr{G})$ that extends $\varphi$.
        \[
            \begin{tikzcd}[column sep = 60pt]
                \sigma(\mathscr{G}) \ar[r, "\mu"] & {[0,\infty]} \\
                \mathscr{G} \ar[u, "\iota"] \ar[ur, "\varphi"'] &
            \end{tikzcd}
        \]
    \end{thm}
    \begin{proof}
        Let $\mu,\nu$ be measures on $\sigma(\mathscr{G})$ that extends $\varphi$. For $G\in\mathscr{G}$ with $\varphi(G)<\infty$, the set \[
        \{A\in\sigma(\mathscr{G})\,|\,\mu_{\cap G}(A)=\nu_{\cap G}(A)\}\]
        is a Dynkin system that contains $\mathscr{G}$, and hence equal to $\sigma(\mathscr{G})$ by Theorem \ref{generators}. Let $(G_n)_{n\in\mathbb{N}}$ be a countable cover of $E$ from $\{\varphi<\infty\}$. Let $F_0=\varnothing$, $F_n=G_1\cup\cdots\cup G_n$. Then $F_\bullet\uparrow E$ and $F_n=\bigsqcup_{k=1}^n G_k\cap F_{k-1}^c$. Hence $\mu(A\cap F_n)=\sum_{k=1}^n\mu_{\cap G_k}(A\cap F_{k-1}^c)=\sum_{k=1}^n\nu_{\cap G_k}(A\cap F_{k-1}^c)=\nu(A\cap F_n)$ for all $A\in\sigma(\mathscr{G})$. Letting $n\to\infty$ yields the claim by Theorem \ref{characterization of premeasures by monotone convergence}.
    \end{proof}
    \begin{thm}\label{uniqueness theorem for finite measures}
        Let $\mu,\nu:\mathscr{E}\rightarrow [0,\infty[$ be two finite measures over $E$ with $\mu(E)=\nu(E)$. If $\mu$ and $\nu$ agree on a $\cap$-stable generator of $\mathscr{E}$, then $\mu=\nu$.

        In particular, two probability measures that agree on a $\cap$-stable generator are identical.
    \end{thm}
    \begin{proof}
        Let the $\cap$-stable generator be $\mathscr{G}$. Then $\mathscr{G}\cup\{E\}$ is also $\cap$-stable. Then $\mu|_{\mathscr{G}\cup\{E\}}=\nu|_{\mathscr{G}\cup\{E\}}$ is $\sigma$-finite. Apply Theorem \ref{uniqueness theorem}.
    \end{proof}
\subsection{The comparison theorem}
    \begin{thm}\label{comparison theorem}
        Let $\mathscr{H}\subset \mathscr{P}X$ be a semiring and $\mu,\nu$ two measures on $\sigma(\mathscr{H})$. If
        \begin{tight_enumerate}
            \item $\mu(S)\leqslant \nu(S)$ for all $S\in\mathscr{H}$,
            \item $\nu|_\mathscr{H}$ is $\sigma$-finite (and hence $\mu|_\mathscr{H}$ is $\sigma$-finite),
        \end{tight_enumerate}
        then $\mu(A)\leqslant\nu(A)$ for all $A\in\sigma(\mathscr{H})$.
    \end{thm}
\subsection{Outer measures}
    \begin{define}
        An \textit{outer measure} is a function $\omega:\mathscr{P}X\rightarrow \overline{\mathbb{R}}_+$ satisfying
        \begin{gather*}
            \omega(\varnothing)=0\\
            A\subset B\Longrightarrow \omega(A)\leqslant \omega(B)\\
            (A_n)\in (\mathscr{P}X)^\mathbb{N}\Longrightarrow  \omega\left(\bigcup_{n=1}^\infty A_n\right)\leqslant\sum_{n=1}^\infty\omega(A_n)
        \end{gather*}
    \end{define}
    $\sigma$-subadditivity of an outer measure $\omega$, together with $\omega(\varnothing)=0$, yields finite subadditivity.

    The sum of inifinitely or finitely many outer measures is an outer measure.
\subsection{Carath\'{e}odory-measurable sets}
    \begin{define}
        Let $\omega:\mathscr{P}X\rightarrow \overline{\mathbb{R}}_+$ be an outer measure.
        \[\mathscr{C}(\omega):=\{A\in\mathscr{P}X\,|\,\omega(Q)\geqslant\omega(Q\cap A)+\omega (Q\cap A^c)\text{ for all }Q\in\mathscr{P}X\}\]
    \end{define}
    \begin{thm}
        \begin{tight_enumerate}
            \item $\omega(A)=0\vee \omega(A^c)=0\Longrightarrow A\in \mathscr{C}(\omega)$.
            \item $\mathscr{C}(\omega)=\{A\in\mathscr{P}X\,|\,\omega(Q)\geqslant\omega(Q\cap A)+\omega (Q\cap A^c)\text{ for all }Q\in\mathscr{P}X\text{ with }\omega(Q)<\infty\}$.
            \item $\mathscr{C}(\omega)=\{A\in\mathscr{P}X\,|\,\omega(Q)=\omega(Q\cap A)+\omega (Q\cap A^c)\text{ for all }Q\in\mathscr{P}X\}$.
        \end{tight_enumerate}
    \end{thm}
    \begin{proof}
        
    \end{proof}
\subsection{Restriction of an outer measure to a measure}
    \begin{thm}
        $\mathscr{C}(\omega)$ is a $\sigma$-algebra and $(X,\mathscr{C}(\omega),\omega|_{\mathscr{C}(\omega)})$ is a measure space. 
        
        That is, in the commutative diagram
        \[
            \begin{tikzcd}[column sep = 60pt]
                \mathscr{C}(\omega) \ar[r, "\iota", hook] \ar[rr, ":=\,\mu"', bend right=18] & \mathscr{P}(X) \ar[r, "\omega"] & {[0,\infty]}
            \end{tikzcd}
        \]
        $\mu$ is a measure.
    \end{thm}
    \begin{proof}
        
    \end{proof}
\subsection{The Lebesgue outer measure}
    \begin{thm}
        Let $\varnothing\in\mathscr{E}\subset\mathscr{P}X$. For any function $\varphi:\mathscr{E}\rightarrow [0,\infty]$ with $\varphi(\varnothing)=0$, the function
        \[
            \varphi^*:\mathscr{P}X\rightarrow [0,\infty],\quad A\mapsto\inf\left\{\sum_{n=1}^\infty \varphi(S_n)\,\bigg|\,(S_n)\in \mathscr{E}^\mathbb{N},\,A\subset \bigcup_{n=1}^\infty S_n\right\}
        \]
        is an outer measure with $\varphi^*|_\mathscr{E}\leqslant \varphi$. Moreover, $\varphi^*$ is maximal: for all outer measures $\omega$ on $\mathscr{P}X$,
        \[
            \omega|_{\mathscr{E}}\leqslant \varphi \Longrightarrow \omega\leqslant \varphi^*
        \]
    \end{thm}
    \begin{proof}
        
    \end{proof}
% \subsection{Carath\'eodory's Theorem}
%     \begin{thm}\label{existence of measure}
%         Be $\mathscr{S}\subset\mathscr{P}(X)$ a semiring. Be $m:\mathscr{S}\rightarrow \overline{\mathbb{R}}_+$ a premeasure. Then
%         \begin{tight_enumerate}
%             \item For $A\in\mathscr{P}(X)$ the following are equivalent
%             \begin{tight_enumerate}
%                 \item $m^*(A)=0$.
%                 \item $A\in\mathscr{N}_{m^*|_{\mathscr{C}(m^*)}}$.
%                 \item $A\subset N\in\mathscr{N}_{m^*|_{\sigma(\mathscr{S})}}$.
%             \end{tight_enumerate}
%             \item $(X,\mathscr{C}(m^*),m^*|_{\mathscr{C}(m^*)})$ is the completion of $(X,\sigma(\mathscr{S}),m^*|_{\sigma(\mathscr{S})})$.
%         \end{tight_enumerate}
%         %Then $m$ has an extension to a measure $\mu$ on $(X,\sigma  (\mathscr{S}))$. If, moreover, $m$ is $\sigma$-finite, then the extension $\mu$ is unique.
%     \end{thm}
%     \begin{proof}
%         \begin{tight_enumerate}
%             \item[5.] (a)$\Leftrightarrow$(b): Necessity. By subadditivity of outer measure, $m^*(Q)\leqslant m^*(Q\cap A)+m^*(Q\setminus A)$. By monotonicity of outer measure, $m^*(Q)\geqslant m^*(Q\setminus A)=m^*(A)+m^*(Q\setminus A)\geqslant m^*(Q\cap A)+m^*(Q\setminus A)$. Thus,     $A\in \mathscr{L}(m^*)$. Sufficiency is obvious. (a)$\Leftrightarrow$(c): Necessity. For each $k\in\mathbb{N}$, there exists a sequence $(S_{i})\in \mathscr{S}^\mathbb{N}$ whose union $S_k=\bigcup_{i\in\mathbb{N}}S_{i}\in\sigma(\mathscr{S})$ with $S_k\supset A$ and $m^*(S_k)\leqslant\sum_{i\in\mathbb{N}}m(S_{i})<1/k$. Then $A\subset\bigcap_{k\in\mathbb{N}}S_k\in \sigma(\mathscr{S})$ with $m^*(\bigcap_{k\in\mathbb{N}}S_k)=0$.
%             \item[6.] From 5 we conclude that $\overline{\sigma(\mathscr{S})}\subset\mathscr{L}(m^*)$. Conversely, let $A\in \mathscr{L}(m^*) $. Find for each $k\in\mathbb{N}$ a sequence $(S_i)\in\mathscr{S}^\mathbb{N}$ such that $S_k=\bigcup_{i\in\mathbb{N}}S_i\supset A$ and $m^*(S_k)-m^*(A)\leqslant\sum_{i\in\mathbb{N}}m(S_i)-m^*(A)<1/k$. Let $S=\bigcap_{k\in\mathbb{N}}S_k\supset A$.  Then $m^*(S\setminus A)=m^*(S)-m^*(A)\leqslant m^*(S_k)-m^*(A)<1/k$ for all $k\in\mathbb{N}$. Thus, $m^*(S\setminus A)=0$ and by 5 we conclude that $S\setminus A$ is a subset of a $m^*|_{\sigma(\mathscr{S})}$-null set and hence in $\overline{\sigma(\mathscr{S})}$.
%         \end{tight_enumerate}
%     \end{proof}
\subsection{Extension of premeasures to measures}
    \begin{thm}\label{extension of premeasures on a semiring to measures}
       Let $\gamma$ be a content on a semiring $\mathscr{H}\subset\mathscr{P}X$.
        \begin{tight_enumerate}
            \item $\mathscr{H}\subset \mathscr{C}(\gamma^*)$, where $\gamma^*$ is the Lebesgue outer measure associated with $\gamma$.
            \item If $\gamma$ is a premeasure, then the diagram
            \[
                \begin{tikzcd}[column sep = {40pt}]
                    \mathscr{H} \ar[r, hook, "\iota"' ] \ar[rrrr, "\gamma", bend left=14] &\sigma(\mathscr{H}) \ar[r, hook, "\iota"'] \ar[rrr, ":=\,\nu"', bend right=24] & \mathscr{C}(\gamma^*) \ar[r, hook, "\iota"'] \ar[rr, ":=\,\mu"', bend right=16]& \mathscr{P}X \ar[r, hook, "\gamma^*"'] & {[0,\infty]}
                \end{tikzcd}
            \]
            commutes, i.e., $\gamma^*|_\mathscr{H}=\gamma$. Moreover, $\mu$ is a measure, and if $\gamma$ is $\sigma$-finite, then $(X,\mathscr{C}(\gamma^*),\mu)$ is the completion of $(X,\sigma(\mathscr{H}),\mu|_{\sigma(\mathscr{H})})$, $\nu$ is the unique measure on $\sigma(\mathscr{H})$ that extends $\gamma$ and $\mu$ is the unique measure on $\mathscr{C}(\gamma^*)$ that extends $\gamma$.
            \item If $\gamma$ is not a premeasure, then there exists $S\in\mathscr{H}$ with $\gamma^*(S)<\gamma(S)$.
        \end{tight_enumerate}
    \end{thm}
    Remark: %If $\mathscr{H}$ is a semiring and $\gamma$ a content, then
    \begin{align*}
        \gamma^*(A)&=\inf\left\{\sum_{n=1}^\infty \overline{\gamma}(B_n)\,\bigg|\,(B_n)\in \mathfrak{r}(\mathscr{H})^\mathbb{N},\,\,A\subset \bigcup_{n\in\mathbb{N}}B_n\right\}\quad \text{by Theorem \ref{first extension theorem}}\\
        &=\inf\left\{\sum_{n=1}^\infty \overline{\gamma}(B_n)\,\bigg|\,(B_n)\in \mathfrak{r}(\mathscr{H})^\mathbb{N} \text{ disjoint},\,\,A\subset \bigcup_{n\in\mathbb{N}}B_n\right\} \quad\text{by Corollary \ref{every countable union is a disjoint union in a ring}}\\
        &=\inf\left\{\sum_{n=1}^\infty \gamma(S_n)\,\bigg|\,(S_n)\in \mathscr{H}^\mathbb{N} \text{ disjoint},\,\,A\subset \bigcup_{n\in\mathbb{N}}S_n\right\} \quad\text{by Theorem \ref{ring generated by semiring}}
    \end{align*}
    \begin{proof}
        
    \end{proof}
    % \begin{crl}
    %     Every $\sigma$-finite premeasure on a semiring $\mathscr{H}$ can be extended in exactly one way to a measure on $\sigma(\mathscr{H})$.
    % \end{crl}
    % \begin{crl}
    %     If $\gamma$ is a $\sigma$-finite premeasure on a semiring $\mathscr{H}\subset \mathscr{P}X$, then there exists a unique measure $\mu$ on $\mathscr{C}(\gamma^*)$ that extends $\gamma$.
    %     \[
    %         \begin{tikzcd}[column sep = 60pt]
    %             \mathscr{C}(\gamma^*) \ar[r, "\overline{\mu}", "\exists!\,\text{measure}"'] & {[0,\infty]}\\
    %             \mathscr{H} \ar[ur, "\gamma"'] \ar[u, "\iota"] &
    %         \end{tikzcd}
    %     \]
    % \end{crl}
    \begin{crl}
        Let $\mathscr{L}^p:=\mathscr{C}((\vol^p)^*)$. The diagram
        \[
            \begin{tikzcd}[column sep = {40pt}]
                \llbracket\mathbb{R}^p,\mathbb{R}^p\llbracket \ar[r, hook, "\iota"' ] \ar[rrrr, "\vol", bend left=14] &\mathscr{B}^p \ar[r, hook, "\iota"']  \ar[rrr, ":=\,\beta^p"', bend right=24] & \mathscr{L}^p \ar[rr, ":=\,\lambda^p"', bend right=16] \ar[r, hook, "\iota"']& \mathscr{P}X \ar[r, hook, "(\vol^p)^*"'] & {[0,\infty]}
            \end{tikzcd}
        \]
        commutes, with $\lambda^p$ and $\beta^p$ called respectively the Lebesgue measure and the Borel measure on $\mathbb{R}^p$.

        The Borel measure is the unique measure on $\mathscr{B}(\mathbb{R}^p)$ that extends $\vol$ and the Lebesgue measure is the unique measure on $\mathscr{L}(\mathbb{R}^p)$ that extends $\vol$. The Lebesgue measure is the completion of the Borel measure.
    \end{crl}
    If $A\subset\mathbb{R}^p$ is countable, then $A$ is $\beta^p$-null (and a fortiori, $\lambda^p$-null).
    \begin{thm}
        Hyperplanes in $\mathbb{R}^p$ are $\beta^p$-null.
    \end{thm}
    \begin{remark}[4.8]
        \textup{a)} Ein intuitiv naheliegender Weg zur Fortsetzung von Pr\"a\-ma{\ss}en zu Ma{\ss}en wird von D.~Maharam (Port.\ Math.\ \textbf{44}, 265--282 (1987)) vorgeschlagen:
        Sie betrachtet ein auf einer Algebra $\mathcal{A}$ \"uber $X$ definiertes Pr\"ama{\ss} $\mu : \mathcal{A} \to \overline{\mathbb{R}}$
        und definiert mithilfe des \"au{\ss}eren Ma{\ss}es $\eta$ eine Topologie auf $\mathcal{P}(X)$.
        Sodann zeigt sie, dass der Abschluss $\overline{\mathcal{A}}$ von $\mathcal{A}$ bez.\ dieser Topologie eine
        $\sigma$-Algebra ist und $\eta|_{\overline{\mathcal{A}}}$ eine Ma{\ss}fortsetzung von $\mu$.
        Die $\sigma$-Algebra $\overline{\mathcal{A}}$ ist gleich $\mathcal{A}_\eta$
        (s.\ loc.\ cit., Theorem~4).
        
        \medskip
        \textup{b)} Im Hinblick auf den Fortsetzungssatz stellt sich die Frage nach weiteren
        Ma{\ss}fortsetzungen von $\mu$. Zu diesem Problem gibt es eine umfangreiche Literatur;
        s.\ W.~Hackenbroch, \emph{Ann.\ Univ.\ Sarav., Ser.\ Math.}\ \textbf{2}, No.~2,
        137--158 (1989).
    \end{remark}
\section{The Lebesgue measure}
\subsection{Contents and premeasures on \texorpdfstring{$\mathbb{R}$}{R}}
    \begin{thm}
        The following diagram in \cat{Set}
        % \[
        %     \begin{tikzcd}
        %         \{\text{increasing functions } \mathbb{R}\rightarrow\mathbb{R}\}/\{f-g=\text{const}\}
        %          & \{\text{increasing functions } \mathbb{R}\rightarrow\mathbb{R}\} \ar[dl, shift left=4pt, "\,\,f\,\mapsto\, (\mu_f:\,[a{,}b[\,\,\mapsto\, f(b)-f(a))"]
        %         \\
        %         \{\text{finite contents on }\mathscr{I}\} \ar[u, shift left=4pt, "\mu\,\mapsto\, \left(f_\mu:\,x\,\mapsto\, (x\geqslant 0\,?\, \mu([0{,}x[) \,:\, -\mu([x{,}0[))\right)\,\,"] &
        %     \end{tikzcd}
        % \]
        \[
            \begin{tikzcd}
                &\{\text{finite contents on }]\mathbb{R},\mathbb{R}]\} \ar[dr, "\id"] \ar[dl, hook, "\mu\,\mapsto\, I_\mu"'] & \\
                \{\text{increasing functions }\mathbb{R}\rightarrow\mathbb{R}\} \ar[rr, twoheadrightarrow, "f\,\mapsto\,\mu_f"]& & \{\text{finite contents on }]\mathbb{R},\mathbb{R}]\}
            \end{tikzcd}
        \]
        commutes, where $\mu_f:[a,b[\,\,\mapsto f(b)-f(a)$ and $I_\mu:x\mapsto (x\geqslant 0\,?\, \mu(]0{,}x]) : -\mu(]x{,}0])))$.

        % given by
        % \begin{align*}
        %     f&\mapsto (\mu_f:[a,b[\,\,\mapsto f(b)-f(a))\\
        %     \mu&\mapsto \left(f_\mu:x\mapsto \begin{cases}
        %         \mu([0,x[)&\text{for }x\geqslant 0\\
        %         -\mu([x,0[)& \text{for }x<0
        %     \end{cases}\right)
        % \end{align*}
        Furthermore, for $\mu$ any finite content on $]\mathbb{R},\mathbb{R}]$,
        \[
            \mu\text{ is a premeasure} \Longleftrightarrow I_\mu\text{ is right-continuous}
        \]
    \end{thm}
    \begin{proof}
        
    \end{proof}
    % \begin{thm}
    %     If $f:\mathbb{R}\rightarrow\mathbb{R}$ is an increasing function, then \[\mu_f:\mathscr{I}\rightarrow\mathbb{R},\,\,[a,b[\,\,\mapsto f(b)-f(a)\] is a finite content. If both $f,g:\mathbb{R}\rightarrow\mathbb{R}$ are increasing functions, then \[\mu_f=\mu_g\Longleftrightarrow f-g=\text{const}\]
    %     Furthermore, if $\mu:\mathscr{I}\rightarrow\mathbb{R}$ is a finite content, then
    %     \[
    %         f:\mathbb{R}\rightarrow \mathbb{R},\,\, x\mapsto \begin{cases}
    %             \mu([0,x[)&\text{for }x\geqslant 0\\
    %             -\mu([x,0[)& \text{for }x<0
    %         \end{cases}
    %     \]
    %     is an increasing function such that $\mu_f=\mu$.
    % \end{thm}
    \begin{crl}
        $\{\text{increasing functions }\mathbb{R}\rightarrow\mathbb{R}\}/\{f-g=\text{const}\}\cong \{\text{finite contents on }]\mathbb{R},\mathbb{R}]\}$ in \cat{Set}.
    \end{crl}
    \begin{proof}
        $\gamma:f\mapsto \mu_f$ has a right inverse and is therefore surjective and by canonical decomposition there is a unique bijection $\tilde{\gamma}:\End_\cat{Tos}(\mathbb{R})/\!\sim\,\,\rightarrow \{\text{finite contents on }\mathscr{I}\}$, where $\sim$ is the fiber equivalence relation of $\gamma$. Thus it suffices to verify that $\sim\,=\{f-g=\text{const}\}$, i.e., $\gamma(f)=\gamma(g)\Leftrightarrow f-g=\text{const}$.
    \end{proof}
    \begin{thm}
        For every increasing right-continuous function $f:\mathbb{R}\rightarrow\mathbb{R}$, there exists an increasing right-continuous jump function $g:\mathbb{R}\rightarrow\mathbb{R}$ and an increasing continuous function $h:\mathbb{R}\rightarrow\mathbb{R}$ such that $f=g+h$. The functions $g$ and $h$ are uniquely determined by $f$ up to addtive constants. For associated premeasures, one has $\gamma(f)=\gamma(g)+\gamma(h)$. (does $\gamma$ in general preserve addition?)
    \end{thm}
    \begin{proof}
        
    \end{proof}
\subsection{Contens and premeasures on \texorpdfstring{$\mathbb{R}^p$}{Rp}}
    \begin{thm}
        The function 
        \[\vol^p:\llbracket\mathbb{R}^p,\mathbb{R}^p\llbracket\,\rightarrow \mathbb{R},\,\,\llbracket a,b\llbracket\,\mapsto \prod_{i=1}^p (b_i-a_i)\]
        is a premeasure on the semiring $\llbracket\mathbb{R}^p,\mathbb{R}^p\llbracket$.
    \end{thm}
    \begin{proof}
        
    \end{proof}
     To be continued. Skipped temporarily.
\subsection{*Lebesgue-Stieltjes Measures}
\subsection{*The Lebesgue Measure}
    \begin{define}
        Let $\mathscr{L}^p:=\mathscr{L}(\mathbb{R}^p)$. A set $A\subset\mathbb{R}^p$ is said to be \textit{Lebesgue measurable} if $A\in\mathscr{L}^p$.
    \end{define}
    \begin{thm}
        For every $A\in\mathscr{L}^p$ and every $\varepsilon>0$, there exists an open superset $U\supset A$ with $\lambda^p(U\setminus A)<\varepsilon$, and a closed subset $F\subset A$ with $\lambda^p(A\setminus F)<\varepsilon$.
    \end{thm}
    \begin{proof}
        
    \end{proof}
    \begin{crl}
        For every $A\in\mathscr{L}^p$,
        \begin{align*}
            \lambda^p(A)&=\inf\{\lambda^p(U)\,|\,A\subset U,\, U\text{ open}\}\\
            &=\sup\{\lambda^p(F)\,|\, F\subset A,\, F\text{ closed}\}\\
            &=\sup\{\lambda^p(K)\,|\, K\subset A,\,K\text{ compact}\}
        \end{align*}
    \end{crl}
    \begin{proof}
        
    \end{proof}
    \begin{crl}
        For every $A\in\mathscr{L}^p$, there exists a $G_\delta$ set $B\supset A$ and a $F_\sigma$ set $C\subset A$ with $\lambda^p(B\setminus A)=\lambda^p(A\setminus C)=0$. 
    \end{crl}
    \begin{proof}
        
    \end{proof}
    \begin{thm}
        Let $A\subset \mathbb{R}^p$. The following are equivalent:
        \begin{tight_enumerate}
            \item $A\in\mathscr{L}^p$.
            \item for every $\varepsilon>0$ there exists an open set $U$ and a closed set $F$ such that $F\subset A\subset U$ with $\lambda^p(U\setminus F)<\varepsilon$.
            \item there exists a $G_\delta$ set $B\supset A$ and a $F_\sigma$ set $C\subset A$ with $\lambda^p(B\setminus C)=0$.
        \end{tight_enumerate}
    \end{thm}
    \begin{proof}
        
    \end{proof}
    The above theorems and corollaries can be extended to Lebesgue-Stieltjes measures.
    \begin{thm}
        If $A\in\mathscr{L}^p$ and $\lambda^p(A)>0$, then $A-A:=\{x-y\,|\,x,y\in A\}$ is a neighborhood of $0$, i.e., there exists $\delta>0$ with $B_\delta(0)\subset A-A$.
    \end{thm}
    \begin{proof}
        
    \end{proof}
    \begin{thm}
        The boundary $\partial A$ of a convex set $A\subset \mathbb{R}^p$ is a Lebesgue null set. 
        
        In particular, every convex set is Lebesgue measurable.
    \end{thm}
    \begin{proof}
        Recall that $A\setminus\partial A=A^\circ$ and $A\cap \partial A$ is also Lebesgue null by completeness of the Lebesgue measure.
    \end{proof}
    % \begin{thm}
    %     A compact subset of $\mathbb{R}^p$ is Lebesgue null if and only if it is Jordan null.
    % \end{thm}
    % \begin{crl}
    %     Every bounded convex subset of $\mathbb{R}^p$ is Jordan measurable.
    % \end{crl}
    % \begin{proof}
    %     Jordan measurable $\Longleftrightarrow$ Jordan-null boundary?
    % \end{proof}
    % Alternatively, for a bounded set, the inner Jordan measure is the Lebesgue measure of topological interior and the outer Jordan measure is the Lebesgue measure of the closure. Thus, a bounded set is Jordan measurable if and only if it has Lebesgue-null boundary.
\subsection{*The Cantor Discontinuum}
\section{Metric Outer Measures}
    Let $(X,d)$ be a metric space. Recall that for $A,B\subset X$, $d(A,B):=\inf\{d(a,b)\,|\,a\in A,b\in B\}$ and $d(A):=\sup\{d(x,y)\,|\,x,y\in A\}$. If $A$ or $B$ is empty, then $d(A,B)=\inf \varnothing=\infty$.
    \begin{define}
        A \textit{metric outer measure} is an outer measure $\omega:\mathscr{P}X\rightarrow [0,\infty]$ on a metric space $(X,d)$ satisfying for all $A,B\subset X$
        \[d(A,B)>0\Longrightarrow \omega(A\cup B)=\omega(A)+\omega(B)\]
    \end{define}
    \begin{thm}
        Let $\varnothing\in\mathscr{E}\subset\mathscr{P}X$. For any function $\varphi:\mathscr{E}\rightarrow [0,\infty]$ with $\varphi(\varnothing)=0$, the function
        % If
        % \[\forall x\in X\,\forall\varepsilon>0\,\exists S\in\mathscr{S}(x\in S\wedge \diam S<\varepsilon)\]
        \[{\varphi^*_0}:=\sup_{\delta>0}(\varphi|_{\mathscr{E}_\delta})^*=\lim_{\delta\to 0}(\varphi|_{\mathscr{E}_\delta})^*\]
        where $\mathscr{E}_\delta:=\{E\in\mathscr{E}\,|\,d (E)\leqslant\delta\}$, is a metric outer measure.
    \end{thm}
    For $X=\mathbb{R}$, $\mathscr{E}$ the set of bounded subsets of $\mathbb{R}$ and $\varphi=d$ the diameter, $d^*_0$ is equal to the Lebesgue outer measure on $\mathbb{R}$.
    \begin{lem}
        Be $\mu^*$ a metric outer measure on $X$. Let $\mathscr{P}(X)\ni A_n\uparrow A$. Then
        \[\forall n\in\mathbb{N}(d(A_n,A\setminus A_{n+1})>0)\Longrightarrow \mu^*(A)=\lim_{n\to\infty}\mu^*(A_n)\]
    \end{lem}
    \begin{thm}
        An outer measure $\omega:\mathscr{P}X\rightarrow [0,\infty]$ is a metric outer measure if and only if $\mathscr{B}(X)\subset \mathscr{C}(\omega)$.
    \end{thm}
\section{*Hausdorff Measures}
    Let $(X,d)$ be a metric space. Let $f:\overline{\mathbb{R}}_+\rightarrow\overline{\mathbb{R}}_+$ be an increasing, right-continuous function with $f(0)=0$ and $f(\infty)=\infty$. The metric outer measure $\eta^f:=(f\circ \diam)^*_0$ is called Hausdorff measure w.r.t. $f$. %The restriction $\mathcal{H}^\phi:=\overline{\mathcal{H}}^\phi|_{\mathscr{B}(X)}$ is a measure on $(X,\mathscr{B}(X))$.
    The restriction $\eta^f|_{\mathscr{B}(X)}$ is a measure on $(X,\mathscr{B}(X))$.
    \begin{lem}
        Denote by $\mathscr{T}$ the topology on $X$ and $\mathscr{T}^c$ the collection of closed sets. For all $\varepsilon>\delta$
        \[(\phi\circ\diam)|_{\mathscr{T}_\varepsilon}^*\leqslant (\phi\circ\diam)|_{\mathscr{T}^c_\delta}^*\leqslant (\phi\circ\diam)|_{\mathscr{P}_\delta}^*\leqslant (\phi\circ\diam)|_{\mathscr{T}_\delta}^*\]
        In particular, $\overline{\mathcal{H}}^\phi=(\phi\circ\diam|_{\mathscr{T}})_0^*=(\phi\circ\diam|_{\mathscr{T}^c})_0^*$.
    \end{lem}
    Recall that an isometry or motion of $(X,d)$ is a bijection that preserves distance. If $g$ is a motion, then $\eta^f(g(A))=\eta^f(A)$ for all $A\in\mathscr{P}X$. In particular, $A\in\mathscr{C}(\eta^f)$ if and only if $g(A)\in\mathscr{C}(\eta^f)$. Moreover, since an isometry is a homeomorphism, $A\in\mathscr{B}(X)$ if and only if $g(A)\in\mathscr{B}(X)$. Thus, the restriction of $\eta^f$ to $\mathscr{C}(\eta^f)$ and $\mathscr{B}(X)$ are both motion invariant measures.
\section{*Rectifiable Curves}
\section{*Hausdorff Dimension}
\section{*Regular Measure}
    % Be $(X,\mathscr{T})$ a topological space. Be $(X,\mathscr{A},\mu)$ a measure space. $A\in\mathscr{A}$ is said to be
    % \begin{enumerate}
    %     \item inner regular if
    %     \[\mu(A)=\sup\{\mu(K)\,|\,K\subset A,K\in\mathscr{A},K\text{ is compact}\}\]
    %     \item outer regular if 
    %     \[\mu(A)=\inf\{\mu(G)\,|\,G\supset A,G\in\mathscr{A},G\text{ is open}\}\]
    % \end{enumerate}
    % \par $\mu$ is called an inner (resp. outer) regular measure if every $A\in\mathscr{A}$ is inner (resp. outer) regular.
    Be $X$ a Hausdorff topological space. Be $(X,\mathscr{B}(X),\mu)$ a measure space. $B\in\mathscr{B}(X)$ is said to be
    \begin{tight_enumerate}
        \item inner regular if
        \[\mu(B)=\sup\{\mu(K)\,|\,\text{compact }K\subset B\}\]
        \item outer regular if 
        \[\mu(B)=\inf\{\mu(G)\,|\,\text{open }G\supset B\}\]
    \end{tight_enumerate}
    $\mu$ is called an inner (resp. outer) regular measure if every $B\in\mathscr{B}(X)$ is inner (resp. outer) regular.
\section{Probability measures}
    \begin{define}
        A \textit{probability space} is a measure space $(\Omega,\mathscr{H},\mathbb{P})$ with $\mathbb{P}(\Omega)=1$; the set $\Omega$ is called the \textit{sample space}, whose elements are called \textit{outcomes}; the $\sigma$-algebra $\mathscr{H}$ is called the \textit{history}, whose elements are called \textit{events}.
    \end{define}
    \begin{thm}
        Let $(\Omega,\mathscr{H},\mathbb{P})$ be a probability space.
        \begin{tight_enumerate}
            \item For events $A,B\in \mathscr{H}$,
            \begin{gather*}
                A\subset B\Longrightarrow \mathbb{P}(A)\leqslant \mathbb{P}(B)\\
                \mathbb{P}(A)+\mathbb{P}(B)=\mathbb{P}(A\cup B)+\mathbb{P}(A\cap B)
            \end{gather*}
            \item For a sequence $A_\bullet\in\mathscr{H}^\mathbb{N}$ of events,
            \begin{gather*}
                \mathbb{P}\left(\bigcup_{n=1}^\infty A_n\right)\leqslant \sum_{n=1}^\infty \mathbb{P}(A_n)\\
                A_n\uparrow A\Longrightarrow \mathbb{P}(A_n)\uparrow \mathbb{P}(A)\\
                A_n\downarrow A\Longrightarrow \mathbb{P}(A_n)\downarrow \mathbb{P}(A)
            \end{gather*}
        \end{tight_enumerate}
    \end{thm}
\subsection{Cumulative distribution functions}
    \begin{define}
        The \textit{cumulative distribution function} or \textit{c.d.f.} of a probability measure $\mu$ on $\er$ is
        \[
            c:\mathbb{R}\rightarrow [0,1],\,\,x\mapsto \mu([-\infty,x])
        \]
        and the \textit{c.d.f.} of a random variable $X:(\Omega,\mathscr{H})\rightarrow (\er,\mathscr{B}(\er))$ is that of its distribution.
        % \[
        %     \mathbb{R}\rightarrow [0,1],\,\,x\mapsto \mathbb{P}(\{X\leqslant  x\})
        % \]
    \end{define}
    \begin{thm}
        The c.d.f. of any probability measure on $\er$ is increasing right-continuous.
    \end{thm}
    \begin{thm}\label{limits of c.d.f.}
        Let $c$ be the c.d.f. of a probability measure $\mu$ on $\er$.
        \begin{align*}
            \lim_{x\to -\infty}c(x)&=\mu(\{-\infty\})\\
            \lim_{x\to\infty}c(x)&=1-\mu(\{\infty\})\\
            c(x-):=\lim_{t\uparrow x}c(t)&=\mu([-\infty,x[)
        \end{align*}
        In particular, $c$ is discontinuous at $x\in\mathbb{R}$ if and only if $\mu(\{x\})>0$.
    \end{thm}
    \begin{thm}
        Mapping probability measures on $\er$ to their c.d.f.s yields a bijection $\textup{c.d.f.}$
        \[
            \begin{tikzcd}
                \{\text{probability measures on }\er\} \ar[r, "\textup{c.d.f.}"] & \{\text{increasing right-continuous function }f:\mathbb{R}\rightarrow [0,1]\}\\
                \{\text{probability measures on }\mathbb{R}\} \ar[u, "\iota"] \ar[ur]
            \end{tikzcd}
        \]
        and mapping each probability measure $\nu$ on $\mathbb{R}$ to the measure
        \[
            \mu:\mathscr{B}(\er)\rightarrow [0,\infty],\,\,\overline{B}\mapsto \nu(\overline{B}\cap\mathbb{R})
        \]
        yields an injection $\iota$ whose image is the set of probability measures $\mu$ on $\er$ such that $\mu(\{-\infty\})=\mu(\{\infty\})=0$. Thus, probability measures on $\mathbb{R}$ correspond to increasing right-continuous functions $f:\mathbb{R}\rightarrow [0,1]$ such that
        \[
            \lim_{x\to-\infty}f(x)=0,\quad\lim_{x\to\infty}f(x)=1
        \]
    \end{thm}
    \begin{proof}
        For $f:\mathbb{R}\rightarrow [0,1]$ increasing right-continuous, $]a,b]\mapsto f(b)-f(a)$, $\{\infty\}\mapsto 1-\lim_{t\to\infty}f(t)$, $\{-\infty\}\mapsto \lim_{t\to-\infty}f(t)$ defines a premeasure $\gamma_f$ on the semiring $]\mathbb{R},\mathbb{R}]\cup\{\{\infty\}\}\cup\{\{-\infty\}\}$ as can be verified by Theorem \ref{characterization of premeasures by monotone convergence}. By Theorem \ref{extension of premeasures on a semiring to measures}, $\gamma_f$ extends to a unique measure $\varphi(f)$ on $\mathscr{B}(\er)$ and $\varphi(f)$ is a probability measure by definition of $\gamma_f$. I claim that $\varphi$ is inverse to $\textup{c.d.f.}$. %Then $\mu_f(]\!-\!\infty,x])=\lim_{t\to-\infty}f(x)-f(t)=f(x)-\lim_{t\to-\infty}f(t)$ and $\mu_f([-\infty,x])=f(x)$; it is clear that $\mu_f$ is a probability measure. Now let $c$ be the c.d.f. of a probability measure $\mu$ on $\er$. By definition, $\gamma_c(]a,b])=c(b)-c(a)=\mu(]a,b])$ and $\gamma_c(\{-\infty\})=\lim_{t\to-\infty}c(t)=\mu(\{-\infty\})$ by Theorem \ref{limits of c.d.f.}. By Theorem \ref{uniqueness theorem for finite measures}, $\mu_c=\mu$.

        Let $\psi$ map each probability measure $\mu$ on $\mathscr{B}(\er)$ to its restriction $\mu|_{\mathscr{B}(\mathbb{R})}$. I claim $\psi\circ\iota=\id$.
    \end{proof}
\subsection{Quantile functions}
    Let $f:\mathbb{R}\rightarrow [0,1]$ be an increasing function.
    \begin{define}
        The \textit{quantile function} of $f$ is
        \[
            q:\,\,]0,1[\,\,\rightarrow \er,\,\,u\mapsto \inf{\{x\in\mathbb{R}\,|\,f(x)>u\}} 
        \]
    \end{define}
    \begin{thm}
        \[
            q(u)=\sup{\{x\in\mathbb{R}\,|\,f(x)\leqslant u\}}
        \]
    \end{thm}
    \begin{proof}
        By isotonicity of $f$, for $x\in\mathbb{R}$, if $f(x)\leqslant u$, then $x$ is a lower bound of $\{f>u\}$ and hence $x\leqslant q(u)$. Thus, $q(u)$ is an upper bound of $\{f\leqslant u\}$.
    \end{proof}
    \begin{thm}\label{lemma for quantile functions}
        Let $x\in\mathbb{R}$ and $u\in\,\,]0,1[$.
        \begin{gather*}
            x<q(u)\Longrightarrow f(x)\leqslant u\\
            x>q(u)\Longrightarrow f(x)>u
        \end{gather*}
    \end{thm}
    \begin{proof}
        If $f(x)>u$, then by definition $q(u)\leqslant x$, and the first $\Rightarrow$ follows by contrapositive.

        For the second $\Rightarrow$: by isotonicity of $f$, if $x\notin \{f>u\}$ then $x$ is a lower bound of $\{f>u\}$.
    \end{proof}
    \begin{thm}
        The quantile function $q$ is real-valued if and only if
        \[
            \lim_{x\to-\infty}f(x)=0,\quad\lim_{x\to\infty}f(x)=1
        \]
        Consequently, the quantile function of the c.d.f. of a real-valued random variable is real-valued.
    \end{thm}
    \begin{thm}\label{quantile function is increasing right-continous}
        The quantile function $q$ is increasing right-continuous.
    \end{thm}
    \begin{proof}
        If $q(u)<x<q(u+)$, then $u<f(x)\leqslant u+\varepsilon$ for all $\varepsilon>0$, contradiction.
    \end{proof}
    % \begin{thm}
    %     For $x\in\mathbb{R}$, $u\in \,\,]0,1[$,
    %     \[
    %         q(u)<x\Longrightarrow u<f(x)\Longrightarrow q(u)\leqslant x\Longrightarrow u\leqslant f(x)
    %     \]
    % \end{thm}
    % \begin{proof}
    %     The first two $\Rightarrow$ follow from the isotonicity of $f$. Indeed, if $x\notin \{f>u\}$ then $x$ is a lower bound for $\{f>u\}$. The third $\Rightarrow$ follows from the right-continuity of $f$. Indeed, if $q(u)=x$, then $f(x+\varepsilon)>u$ for all $\varepsilon>0$. Taking limit $\varepsilon\downarrow 0$ yields $f(x)\geqslant u$. %(In view of the first $\Rightarrow$, it suffices to show $q(u)=x\Longrightarrow u\leqslant f(x)$)
    % \end{proof}
    \begin{thm}\label{duality of quantile functions and c.d.f.s}
        If $f$ is right-continuous, then for all $x\in\mathbb{R}$,
        \[
            f(x)=\inf{\{u\in\,\,]0,1[\,|\,q(u)>x\}}
        \]
    \end{thm}
    \begin{proof}
        Fix $x\in\mathbb{R}$. For $u\in \,\,]0,1[$, if $q(u)>x$, then $f(x)\leqslant u$. Thus, $f(x)\leqslant \inf{\{q>x\}}$. If $f(x)<u<\inf{\{q>x\}}$, then $x\leqslant q(u)$. If $x=q(u)$, then $f(x+\varepsilon)>u$ and by the assumption that $f$ is right-continuous, $f(x)=\lim_{\varepsilon\downarrow0}f(x+\varepsilon)\geqslant u$, contradiction. Thus, $u\in \{q>x\}$ and $u\geqslant \inf{\{q>x\}}$.
    \end{proof}
    \begin{thm}
        For $x\in\mathbb{R}$, $u\in \,\,]0,1[$,
        \[
            x\leqslant q(u)\Longleftrightarrow f(x-)\leqslant u
        \]
        and if $f$ is right-continuous, then
        \[
            u\leqslant f(x)\Longleftrightarrow q(u-)\leqslant x
        \]
    \end{thm}
\subsection{The correspondence theorem}
    \begin{thm}
        Mapping probability measures on $\er$ to their c.d.f.s yields a bijection $\textup{c.d.f.}$
        \[
            \begin{tikzcd}
                \{\text{probability measures on }\er\} \ar[r, "\textup{c.d.f.}"] & \{\text{increasing right-continuous function }f:\mathbb{R}\rightarrow [0,1]\}\\
                \{\text{probability measures on }\mathbb{R}\} \ar[u, "\iota"] \ar[ur]
            \end{tikzcd}
        \]
        and mapping each probability measure $\nu$ on $\mathbb{R}$ to the measure $\mu:\mathscr{B}(\er)\rightarrow [0,\infty],\,\,\overline{B}\mapsto \nu(\overline{B}\cap\mathbb{R})\textbf{}$ yields an injection $\iota$.
        
        The inverse of $\textup{c.d.f.}$ maps each increasing right-continuous function $f$ to the distribution $q_*\lambda$ of the quantile function $q:\,\,]0,1[\,\,\rightarrow \er$ as a random variable on the probability space $(]0,1[,\mathscr{B}(]0,1[),\lambda)$, where $\lambda$ is the Lebesgue measure.
        
        The image of $\iota$ is the set of probability measures $\mu$ on $\er$ such that $\mu(\{-\infty\})=\mu(\{\infty\})=0$. Furthermore, for an increasing right-continuous function $f:\mathbb{R}\rightarrow [0,1]$, the following are equivalent:
        \begin{tight_enumerate}
            \item $f$ is the c.d.f. of a probability measure on $\mathbb{R}$, or equivalently, of a real-valued random variable.
            \item \[
                \lim_{x\to-\infty}f(x)=0,\quad\lim_{x\to\infty}f(x)=1
            \]
            \item The quantile function $q$ of $f$ is real-valued.
        \end{tight_enumerate}
    \end{thm}
    \begin{proof}
        % For $f:\mathbb{R}\rightarrow [0,1]$ increasing right-continuous, $]a,b]\mapsto f(b)-f(a)$, $\{\infty\}\mapsto 1-\lim_{t\to\infty}f(t)$, $\{-\infty\}\mapsto \lim_{t\to-\infty}f(t)$ defines a premeasure $\gamma_f$ on the semiring $]\mathbb{R},\mathbb{R}]\cup\{\{\infty\}\}\cup\{\{-\infty\}\}$ as can be verified by Theorem \ref{characterization of premeasures by monotone convergence}. By Theorem \ref{extension of premeasures on a semiring to measures}, $\gamma_f$ extends to a unique measure $\varphi(f)$ on $\mathscr{B}(\er)$ and $\varphi(f)$ is a probability measure by definition of $\gamma_f$. I claim that $\varphi$ is inverse to $\textup{c.d.f.}$. 
        
        %Then $\mu_f(]\!-\!\infty,x])=\lim_{t\to-\infty}f(x)-f(t)=f(x)-\lim_{t\to-\infty}f(t)$ and $\mu_f([-\infty,x])=f(x)$; it is clear that $\mu_f$ is a probability measure. Now let $c$ be the c.d.f. of a probability measure $\mu$ on $\er$. By definition, $\gamma_c(]a,b])=c(b)-c(a)=\mu(]a,b])$ and $\gamma_c(\{-\infty\})=\lim_{t\to-\infty}c(t)=\mu(\{-\infty\})$ by Theorem \ref{limits of c.d.f.}. By Theorem \ref{uniqueness theorem for finite measures}, $\mu_c=\mu$.
        I verify the statement about the inverse of \textup{c.d.f.} Let $f:\mathbb{R}\rightarrow[0,1]$ be an increasing right-continuous function. Then $]0,f(x)[\,\,\subset\{q\leqslant x\}$ is by definition of $q$. If $q(u)\leqslant x$, then by Theorem \ref{duality of quantile functions and c.d.f.s}, $u\leqslant f(x)$. Thus, $\{q\leqslant x\}\subset \,\,]0,f(x)]$ and $q_*\lambda([-\infty,x])=\lambda(\{q\leqslant x\})=f(x)$. Let $\mu$ be a probability measure on $\er$ and $c$ its c.d.f. Let $q$ be the quantile function of $c$. Then $q_*\lambda([-\infty,x])=c(x)=\mu([-\infty,x])$ for all $x\in\mathbb{R}$ by the preceding argument. This implies that the probability measures $q_*\lambda$ and $\mu$ coincide on the $\cap$-stable generator $]\mathbb{R},\mathbb{R}]\cup\{\{-\infty\}\}$ of $\mathscr{B}(\er)$. By Theorem \ref{uniqueness theorem for finite measures}, $\mu=q_*\lambda$.  
            
        Let $\psi$ map each probability measure $\mu$ on $\mathscr{B}(\er)$ to its restriction $\mu|_{\mathscr{B}(\mathbb{R})}$. I claim $\psi\circ\iota=\id$.
    \end{proof}
\subsection{*Probability measures on \texorpdfstring{$\mathbb{R}^n$}{Rn}}
    See Elstrodt p66.
\chapter{Measurable Functions}
\section{The category Mes}
    \begin{define}
        A \textit{measurable space} consists of a set $E$ and a $\sigma$-algebra $\mathscr{E}$ on $E$.
    \end{define}
    \begin{define}
        Let $(E,\mathscr{E})$, $(F,\mathscr{F})$ be measurable spaces. Let 
        \[
            \mathscr{E}/\mathscr{F}:=\{f:E\rightarrow F\,|\,f^{-1}\mathscr{F}\subset \mathscr{E}\}
        \]
        functions in which are said to be $\mathscr{E}/\mathscr{F}$-measurable.
    \end{define}
    \begin{thm}\label{measurable functions}
        Let $(E,\mathscr{E})$, $(F,\mathscr{F})$ be measurable spaces. A function $f:E\rightarrow F$ is $\mathscr{E}/\mathscr{F}$\textit{-measurable} if and only if $f^{-1}\mathscr{C}\subset \mathscr{E}$ for some generator $\mathscr{C}$ of $\mathscr{F}$.
    \end{thm}
    \begin{proof}
        Theorem \ref{sigma commutes with preimage}. 
    \end{proof}
    \begin{thm}
        Measurable spaces and measurable functions form a category \cat{Mes}.
    \end{thm}
    \begin{proof}
        If $f\in \mathscr{E}/\mathscr{F}$ and $g\in \mathscr{F}/\mathscr{G}$, then $(g\circ f)^{-1}\mathscr{G}=f^{-1}g^{-1}\mathscr{G}\subset f^{-1}\mathscr{F}\subset \mathscr{E}$. 
        % $\circ$ is the usual function composition
    \end{proof}
    % For $X',X''$ two measurable spaces with the same underlying set $X$, the identity map $\id_X$ is measurable if and only if $\Sigma X''\subset \Sigma X'$.
\subsection{Measurability of composite functions}
    \begin{thm}
        Let $(E,\mathscr{E})$, $(G,\mathscr{G})$ be measurable spaces. For a diagram
        \[
            \begin{tikzcd}
                E \ar[r,"f"] & F\ar[r, "g"]& G
            \end{tikzcd}
        \]
        in \cat{Set}, $g\circ f$ is $\mathscr{E}/\mathscr{G}$-measurable if and only if $f$ is $\mathscr{E}/g^{-1}\mathscr{G}$-measurable.
    \end{thm}
    \begin{thm}
        Let $X,Y$ be measurable spaces. For any function $f:X\rightarrow Y$ and any set $Z$ such that $f(X)\subset Z\subset Y$, if $Z$ is equipped with the trace $\sigma$-algebra $\mathscr{Y}|_Z$, then the inclusion map $\iota:Z\rightarrow Y$ is measurable, and $f:X\rightarrow Y$ is measurable if and only if $\hat{f}:X\rightarrow Z$ is measurable.
        \[
            \begin{tikzcd}
                X \ar[dr,"f"'] \ar[r,"\hat{f}"]& Z \ar[d,"\iota"] \\
                & Y
            \end{tikzcd}
        \]
    \end{thm}
    \begin{thm}\label{measurability w.r.t. codomain restricted to a topological subspace}
        Let $X$ be a measurable space and $Y$ a topological space. Let $Z$ be a topological subspace of $Y$ and $\iota:Z\rightarrow Y$ the inclusion map. For any function $f:X\rightarrow Z$ is 
        \[
            f\text{ is $\mathscr{X}/\mathscr{B}(Z)$-measurable}\Longleftrightarrow \iota\circ f\text{ is $\mathscr{X}/\mathscr{B}(Y)$-measurable}
        \]
        % For any function $f:X\rightarrow Y$ and any topological subspace $Z$ of $Y$ with $f(X)\subset Z\subset Y$,
        % \[f\text{ is $\mathscr{X}/\mathscr{B}(Y)$-measurable}\Longleftrightarrow f\text{ is $\mathscr{X}/\mathscr{B}(Z)$-measurable}\]
    \end{thm}
% \section{Subspaces}
%     \begin{define}
%         A \textit{subspace} of a measurable space $X$ is a measurable space $Y$ such that $\ul Y\in \Sigma X$ and $\Sigma Y=\Sigma X|_{Y}$.
%     \end{define}
%     \begin{thm}
%     If $A\in\mathscr{A}$, then 
%     \[  
%         (A,\mathscr{A}|_A,\mu|_A):=(A,\mathscr{A}|_A,\mu|_{\mathscr{A}|_A})
%     \]
%     is a measure space, and may be called the $A$-subspace of $(X,\mathscr{A},\mu)$.
%     \end{thm}
%     \begin{proof}
%         $A\in\mathscr{A}\Longrightarrow \mathscr{A}|_A\subset\mathscr{A}$. Thus, $\mu|_{\mathscr{A}|_A}$ is well-defined and inherits $\sigma$-additivity of $\mu$.
%     \end{proof}
\subsection{Products}
    Let $(X_i,\mathscr{X}_i)_{i\in I}$ be an indexed family of measurable spaces. Let the product of $(X_i)_{i\in I}$ in \cat{Set} be
    \[
        \begin{tikzcd}
            X \ar[d, "\pi_i"]\\
            X_i
        \end{tikzcd}_{i\in I}
    \]
    \begin{define}
        The \textit{product $\sigma$-algebra} on $X$ is
        \[
            \mathscr{X}=\bigotimes_{i\in I}\mathscr{X}_i:=\sigma \left(\bigcup_{i\in I}\pi_{i}^{-1}(\mathscr{X}_i)\right)=\sigma\{\pi_i^{-1}(A_i)\,|\,A_i\in\mathscr{X}_i,\,i\in I\}
        \]
        i.e., the $\sigma$-algebra generated by the projections $(\pi_i)_{i\in I}$.
    \end{define}
    \begin{thm}\label{product measurable space}
        $(X,\mathscr{X})$ together with $(\pi_i)_{i\in I}$ is a product of $(X_i,\mathscr{X}_i)_{i\in I}$ in \cat{Mes}.
        % In particular, for any measurable space $Z$ if $u:Z\rightarrow X$ and $v:Z\rightarrow Y$ are measurable functions, then $(u,v):Z\rightarrow X\times Y,\,\,z\mapsto (u(z),v(z))$ is measurable.
    \end{thm}
    \begin{proof}
        The projections $(\pi_i)_{i\in I}$ are clearly measurable.
        Since $X$ is a product of $(X_i)_{i\in I}$ in \cat{Set}, it suffices to show that if $u_i:Z\rightarrow X_i$ is a measurable function for each $i\in I$, then the function $u:Z\rightarrow X,\,\,x\mapsto (u_i(x))_{i\in I}$ is measurable. Indeed, by Theorem \ref{measurable functions}, it suffices to check that $u^{-1}(\pi_i^{-1}(A_i))=(\pi_i\circ u)^{-1}(A_i)=u_i^{-1}(A_i)\in\mathscr{Z}$ for all $A_i\in\mathscr{X}_i$ and any $i\in I$.
    \end{proof}
    \begin{thm}\label{sections are measurable}
        For any $j\in I$, fix $a_i\in X_i$ for each $i\in I\setminus \{j\}$. The function
        \[
            \iota:X_j\rightarrow X,\,\,x_j\mapsto \left(i\mapsto \begin{cases}
                a_i&\text{ for }i\neq j\\
                x_j&\text{ for }i=j
            \end{cases}\right)
        \]
        is $\mathscr{X}_j/\mathscr{X}$-measurable.
    \end{thm}
    \begin{proof}
        By Theorem \ref{measurable functions}, it suffices to show $\iota^{-1}(\pi_i^{-1}(A_i))=(\pi_i\circ \iota)^{-1}(A_i)\in\mathscr{X}_j$ for any $i\in I$, $A_i\in\mathscr{X}_i$. Indeed, $\pi_i\circ \iota:X_j\rightarrow X_i$ is the constant map to $a_i$ if $i\neq j$ and the identity map of $X_j$ if $i=j$. In the former case, $(\pi_i\circ \iota)^{-1}(A_i)$ is either $\varnothing$ or $X_j$, and in the latter, $(\pi_j\circ \iota)^{-1}(A_j)=A_j\in\mathscr{X}_j$.
    \end{proof}
    The converse is not true, see Exercise 6.28 of GTM261.
    \begin{thm}\label{generator of arbitrary product sigma algebra}
        The product $\sigma$-algebra $\mathscr{X}$ is also generated by the $\cap$-stable system
        \[
            \mathscr{G}=\left\{\bigtimes_{i\in I}A_i\,\bigg|\,A_i\in\mathscr{X}_i\text{ for all }i\in I\text{ and }A_i=X_i\text{ except for finitely many }i\in I\right\}
        \]
    \end{thm}
    \begin{proof}
        $\pi_{i}^{-1}(\mathscr{X}_i)\subset\mathscr{G}$ for all $i\in I$. Thus, $\bigcup_{i\in I}\pi_{i}^{-1}(\mathscr{X}_i)\subset\mathscr{G}$.
        Conversely, every $A_{i_1}\times\cdots\times A_{i_n}\times\prod_{i\in I\setminus\{i_1,\dots,i_n\}}X_i$ in $\mathscr{G}$ is equal to $\pi_{i_1}^{-1}(A_{i_1})\cap \cdots\cap\pi_{i_n}^{-1}(A_{i_n})\in\sigma \left(\bigcup_{i\in I}\pi_{i}^{-1}(\mathscr{X}_i)\right)$.
    \end{proof}
    \begin{thm}\label{Product of Generator is Generator of Product Sigma Algebra}
        Let $\mathscr{X}_1=\sigma(\mathscr{F})$ and $\mathscr{X}_2=\sigma(\mathscr{G})$. If $\mathscr{F}$ contains a countable cover of $X_1$, and $\mathscr{G}$ contains a countable cover of $X_2$, then
         \[\mathscr{X}_1\otimes \mathscr{X}_2=\sigma(\mathscr{F}* \mathscr{G})\]
         where $\mathscr{F}* \mathscr{G}$ denotes $\{A\times B\,|\,A\in\mathscr{F},\,B\in\mathscr{G}\}$.
     \end{thm}
     \begin{proof}
         $\mathscr{X}_1\otimes \mathscr{X}_2=\sigma(\mathscr{X}_1*\mathscr{X}_2)\supset\sigma(\mathscr{F}* \mathscr{G})$ is obvious. Conversely, the system
         \[
         \Sigma=\{A\in\mathscr{X}_1\,|\,A\times G\in \sigma(\mathscr{F}* \mathscr{G})\text{ for all }G\in\mathscr{G}\}
         \]
         is a $\sigma$-algebra on $X_1$ that contains $\mathscr{F}$. Indeed, $X_1\in\Sigma$ since $\mathscr{F}$ contains a countable cover of $X_1$, $A\in\Sigma\Longrightarrow A^c\in\Sigma$ since $A^c\times G=(X\times G)\setminus(A\times G)$, and stability under countable union is clear. Thus, $\mathscr{X}_1=\Sigma$, i.e., $\mathscr{X}_1* \mathscr{G}\subset \sigma(\mathscr{F}* \mathscr{G})$.  Similarly, $\mathscr{F}* \mathscr{X}_2\subset \sigma(\mathscr{F}* \mathscr{G})$. Thus, $\pi_1^{-1}(\mathscr{X}_1),\,\pi_2^{-1}(\mathscr{X}_2)\subset \sigma(\mathscr{F}* \mathscr{G})$ since $A\times X_2$ can be written as a countable union of sets from $\mathscr{X}_1*\mathscr{G}$, and $X_1\times B$ as a countable union of sets from $\mathscr{F}*\mathscr{X}_2$.
     \end{proof}
    \begin{thm}\label{Borel sigma algebra of product space}
        Let $Y_1,Y_2$ be topological spaces.
        \[
            \mathscr{B}(Y_1)\otimes \mathscr{B}(Y_2)\subset\mathscr{B}(Y_1\times Y_2)
        \]
        If $Y_1$ is second-countable, then
        \[
            \mathscr{B}(Y_1)\otimes \mathscr{B}(Y_2)=\mathscr{B}(Y_1\times Y_2)
        \]
        In particular, if $u_1:X\rightarrow Y_1$ is $\mathscr{X}/\mathscr{B}(Y_1)$-measurable and $u_2:X\rightarrow Y_2$ is $\mathscr{X}/\mathscr{B}(Y_2)$-measurable, then $(u_1,u_2):X\rightarrow Y_1\times Y_2$ is $\mathscr{X}/\mathscr{B}(Y_1\times Y_2)$-measurable.
    \end{thm}
    \begin{proof}
        By Theorem \ref{Product of Generator is Generator of Product Sigma Algebra}, $\mathscr{B}(Y_1)\otimes \mathscr{B}(Y_2)$ is generated by $\{\text{open sets in }Y_1\}*\{\text{open sets in }Y_2\}$ which is a subset (base) of $\{\text{open sets in }Y_1\times Y_2\}$. If $Y_1$ admits a countable base $(B_n)_{n\in\mathbb{N}}$, fix $U$ open in $Y_1\times Y_2$, since topology is stable under arbitrary union, one can choose a largest open set $C_n$ in $Y_2$ such that $B_n\times C_n\subset U$ for each $n$. Then $U=\bigcup_{n\in\mathbb{N}} B_n\times C_n$. Indeed, for $(x,y)\in U$, there exists oepn set $B$ in $Y_1$ and $C$ in $Y_2$ with $(x,y)\in B\times C\subset U$, and hence $B_n$, with $(x,y)\in B_n\times C\subset U$. Then $C\subset C_n$ and $(x,y)\in B_n\times C_n$. Thus, $U$ is a countable union of sets from $\{\text{open sets in }Y_1\}*\{\text{open sets in }Y_2\}$.
    \end{proof}
    \textcolor{Aquamarine}{Bogachev seems to require $Y_1,Y_2$ be Hausdorff in his "Measure Theory II", why? What about arbitrary product instead of binary product? GTM261 asserts that the equality holds for the product of a countable family of second-countable spaces, e.g., $\mathbb{R}^\mathbb{N}$.}
    \begin{thm}
        Let $(F_{\bullet},\mathscr{F}_{\bullet}):I\rightarrow \cat{Mes}$ be an indexed family of measurable spaces and $(E,f_{\bullet})$ a cone over $F_{\bullet}$ in \cat{Set}. Then the initial $\sigma$-algebra generated by the family $f_\bullet$ on $E$ is precisely the preimage $\sigma$-algebra generated by $f:E\rightarrow (\prod_{i\in I}F_i,\bigotimes_{i\in I}\mathscr{F}_i),\,\,x\mapsto (f_i(x))_{i\in I}$.
    \end{thm}
\subsection{Measurability on a product space}
    \begin{thm}\label{measurability on a product space}
        Let $(F,\mathscr{F})$ be a measurable space and let
        $f:\mathbb{R}\times F\rightarrow\overline{\mathbb{R}}$.
        Suppose that
        \begin{enumerate}
            \item $y\mapsto f(x,y)$ is $\mathscr{F}$-measurable
            for every $x\in\mathbb{R}$;
            \item $x\mapsto f(x,y)$ is right-continuous
            for every $y\in F$.
        \end{enumerate}
        Then $f$ is
        $\mathscr{B}(\mathbb{R})\otimes\mathscr{F}$-measurable.
    
        The same conclusion holds if right-continuity is replaced
        with left-continuity.
    \end{thm}
    \begin{proof}
        Suppose first that $x\mapsto f(x,y)$ is right-continuous
        for every $y\in F$.
        For each $n\in\mathbb{N}$, define
        \[
            f_n(x,y):=
            f\left(\frac{\lfloor 2^n x\rfloor+1}{2^n},y\right).
        \]
        To verify measurability, observe that, for every
        $B\in\mathscr{B}(\overline{\mathbb{R}})$,
        \[
            f_n^{-1}(B)
            =
            \bigcup_{k\in\mathbb{Z}}
            \left(
                \left[\frac{k}{2^n},\frac{k+1}{2^n}\right[
                \times
                \left\{
                    y\in F:
                    f\left(\frac{k+1}{2^n},y\right)\in B
                \right\}
            \right).
        \]
        Each set in the second coordinate belongs to $\mathscr{F}$
        by assumption. Hence $f_n$ is
        $\mathscr{B}(\mathbb{R})\otimes\mathscr{F}$-measurable.
    
        For every $x\in\mathbb{R}$,
        \[
            0<
            \frac{\lfloor 2^n x\rfloor+1}{2^n}-x
            \leqslant 2^{-n}.
        \]
        Right-continuity therefore gives
        \[
            \lim_{n\to\infty}f_n(x,y)=f(x,y)
            \qquad
            \forall (x,y)\in\mathbb{R}\times F.
        \]
        Thus $f$ is measurable as a pointwise limit of measurable
        functions.
    
        If instead $x\mapsto f(x,y)$ is left-continuous for every
        $y\in F$, apply the preceding argument to
        $g(x,y):=f(-x,y)$.
        Then $g$ is measurable, and so is
        $f(x,y)=g(-x,y)$, since $(x,y)\mapsto(-x,y)$ is measurable.
    \end{proof}
\subsection{Stochastic processes and filtrations}
    Let $(E,\mathscr{E})$, $(F,\mathscr{F})$ be measurable spaces.
    \begin{define}
        A \textit{stochastic process} on $(E,\mathscr{E})$ with \textit{state space} $(F,\mathscr{F})$ and \textit{parameter set} $T$ is a measurable function $\varphi:(E,\mathscr{E})\rightarrow (F,\mathscr{F})^T=(F^T,\mathscr{F}^{\otimes T})$.
    \end{define}
    \begin{thm}
        A function $\varphi:E\rightarrow F^T$ is a stochastic process if and only if for all $t\in T$
        \[
            \pi_t\circ \varphi:E\rightarrow F
        \]
        is measurable.
    \end{thm}
    \begin{define}
        A \textit{filtration} of a measurable space $(E,\mathscr{E})$ is a family of $\sigma$-algebras $(\mathscr{E}_t)_{t\in T}$ on $E$ indexed by an upward directed set $(T,\leqslant)$ such that for $t_1,t_2\in T$
        \[
            t_1\leqslant t_2\Longrightarrow \mathscr{E}_{t_1}\subset \mathscr{E}_{t_2}\subset\mathscr{E}
        \]

        A \textit{filtered measurable space} is a measurable space equipped with a filtration.
    \end{define}
    \begin{thm}
        For a filtration $(\mathscr{E}_t)_{t\in T}$ of a measurable space $(E,\mathscr{E})$,
        \[
            \mathscr{E}_\vee:=\sigma\left(\bigcup_{t\in T}\mathscr{E}_t\right)=\mathfrak{m}\left(\bigcup_{t\in T}\mathscr{E}_t\right)
        \]
    \end{thm}
    \begin{proof}
        $\bigcup_{t\in T}\mathscr{E}_t$ is an algebra. Theorem \ref{generators}.
    \end{proof}
    \begin{define}
        Let $(E,(\mathscr{E}_t)_{t\in T},\mathscr{E})$ be a filtered measurable space with $T\subset \er$. A \textit{stopping time} is a function $\tau:E\rightarrow T\cup\{\infty\}$ satisfying for all $t\in T$
        \[
            \{\tau\leqslant t\}\in\mathscr{E}_t
        \]
        The $\sigma$-algebra associated with a stopping time $\tau$ is
        \[
            \mathscr{E}_\tau:=\{A\in\mathscr{E}\,|\,A\cap \{\tau\leqslant t\}\in \mathscr{E}_t\text{ for all }t\in T\}
        \]
    \end{define}
    Clearly, for any stopping time $\tau$,
    \[
        \bigcap_{t\in T}\mathscr{E}_t\subset \mathscr{E}_\tau
    \]
    \begin{define}
        Let $(E,(\mathscr{E}_t)_{t\in T},\mathscr{E})$ be a filtered measurable space.
        A \textit{stochastic process} $\varphi:(E,\mathscr{E})\rightarrow (F,\mathscr{F})^T$ is \textit{adapted} to the filtration $(\mathscr{E}_t)_{t\in T}$ if for all $t\in T$,
        \[
            \pi_t\circ \varphi\in \mathscr{E}_t/\mathscr{F}
        \]
    \end{define}
\subsection{Pushforward measures}
    \begin{thm}
        Let $u:E_1\rightarrow E_2$ be a measurable function. If $\mu$ is a measure on $\mathscr{E}_1$, then 
        \[
            u_*\mu:\mathscr{E}_2\rightarrow [0,\infty],\,\,A_2\mapsto \mu(u^{-1}(A_2))
        \] is a measure on $\mathscr{E}_2$. If $v:E_2\rightarrow E_3$ is another measurable function, then\footnote{The set of measures on a measurable space together with addition form a commutative monoid. The pushforward is then a functor from \cat{Mes} to \cat{CMon}. \textcolor{Aquamarine}{The measures on a measurable space seem to form a so-called $\Sigma$-monoid.}}
        \[
            (v\circ u)_*\mu=v_*(u_*\mu)
        \]
    \end{thm}
    \begin{thm}\label{pushforward of completion of measures}
        Let $(E,\mathscr{E},\mu),(F,\mathscr{F},\nu)$ be measure spaces and $(E,\overline{\mathscr{E}},\overline{\mu}), (F,\overline{\mathscr{F}},\overline{\nu})$ their completions. 
        
        If $u:E\rightarrow F$ is $\mathscr{E}/\mathscr{F}$-measurable with $\{\nu=0\}\subset \{u_*\mu=0\}$, then $u$ is $\overline{\mathscr{E}}/\overline{\mathscr{F}}$-measurable, and
        \[
            \nu=u_*\mu \Longrightarrow \overline{\nu}=u_*\overline{\mu}
        \]
    \end{thm}
    % \begin{proof}
    %     Let $A_2\cup N'_2\in\overline{\mathscr{A}_2}$. $f^{-1}(A_2\cup N'_2)=f^{-1}(A_2)\cup f^{-1}(N'_2)$ with $f^{-1}(A_2)\in\mathscr{A}_1$ and $f^{-1}(N'_2)\subset f^{-1}(N_2)$ for some $N_2\in\mathscr{N}_\nu\subset \mathscr{N}_{f_*\mu}$, i.e., $\mu(f^{-1}(N_2))=0$, i.e., $f^{-1}(N_2)\in\mathscr{N}_\mu$. So $f$ is $\overline{\mathscr{A}_1}/\overline{\mathscr{A}_2}$-measurable. If $\nu=f_*\mu$, then $\overline{\nu}(A_2\cup N'_2)=\nu(A_2)=\mu(f^{-1}A_2)=\overline{\mu}(f^{-1}A_2\cup f^{-1}N_2')$ since $f^{-1}N_2'\subset f^{-1}N_2\in\mathscr{N}_\mu$.
    % \end{proof}
\section{Borel measurability}
    \begin{define}
        Let $X,Y$ be topological spaces. A function $u:X\rightarrow Y$ is \textit{Borel measurable} if it is $\mathscr{B}(X)/\mathscr{B}(Y)$-measurable.
    \end{define}
    \begin{thm}\label{Continuous Maps are Measurable}
        Continuous maps are Borel measurable, i.e., 
        \[\cat{Top}(X,Y)\subset \cat{Mes}(X,Y)\]
        for topological spaces $X,Y$ equipped with Borel $\sigma$-algebras $\mathscr{B}(X),\mathscr{B}(Y)$.
    \end{thm}
    \begin{proof}
        Theorem \ref{measurable functions}.
    \end{proof}
\subsection{Lower semicontinuous functions}
    \begin{thm}
        Let $\varphi:X\rightarrow \,\,]\!-\!\infty,\infty]$ be l.s.c., where $X$ is a topological space. Then $\varphi$ is Borel measurable.
    \end{thm}
\subsection{Functions with at most countably many discontinuities}
    \begin{thm}
        Let $X,Y$ be topological spaces. If $X$ is $T_1$, i.e., singletons are closed in $X$, and $u:X\rightarrow Y$ has at most countably many discontinuities, then $u$ is Borel measurable.
    \end{thm}
    \begin{proof}
        Let $D$ be the set of discontinuities. $u|_{D^c}$ is continuous and hence Borel measurable. Thus, $\{u<r\}=\{u|_{D^c}<r\}\cup \{u|_D<r\}$, where $\{u|_{D^c}<r\}\in\mathscr{B}(D^c)=\mathscr{B}(X)|_{D^c}\subset \mathscr{B}(X)$ since $D$ as an at most countable union of closed sets (singletons in a $T_1$ space) is in $\mathscr{B}(X)$, and $\{u|_D<r\}\subset D$ for the same reason is in $\mathscr{B}(X)$. Thus, $\{u<r\}\in\mathscr{B}(X)$ for all $r\in\mathbb{R}$.
    \end{proof}
    Thus, monotone functions $f:\mathbb{R}\rightarrow\mathbb{R}$ are Borel measurable because of the following:
    \begin{thm}
        Every monotone function $f:\mathbb{R}\rightarrow\mathbb{R}$ has at most countably many discontinuities.
    \end{thm}
\section{Motion Invariance of the Lebesgue Measure}
    \begin{thm}
        Let $v\in\mathbb{R}^p$. The translation $t_v:x\mapsto x+v$ is both Borel measurable and Lebesgue measurable, and
        \[
            (t_v)_*\beta^p=\beta^p,\quad (t_v)_*\lambda^p=\lambda^p
        \]
    \end{thm}
    \begin{thm}
        If $\mu$ is a translation invariant measure on $\mathscr{B}^p$ (resp. $\mathscr{L}^p$) with $\mu([0,1[^p)=1$ or $\mu([0,1]^p)=1$, then $\mu=\beta^p$ (resp. $\mu=\lambda^p$).
    \end{thm}
    \begin{crl}
        If $\mu$ is a translation invariant measure on $\mathscr{B}^p$ (resp. $\mathscr{L}^p$) with $\mu([0,1]^p)=k<\infty$, then $\mu=k\beta^p$ (resp. $\mu=k\lambda^p$).
    \end{crl}
    \begin{thm}
        Every affine automorphism $f$ of $\mathbb{R}^p$ is both Borel and Lebesgue measurable, and
        \[
            f_*\beta^p=|\det f|^{-1}\beta^p,\quad f_*\lambda^p=|\det f|^{-1}\lambda^p
        \]
        In particular, if $A\in\mathscr{B}^p$ (resp. $A\in\mathscr{L}^p$), then $f(A)\in\mathscr{B}^p$ (resp. $f(A)\in\mathscr{L}^p$) and
        \[
            \lambda^p(f(A))=(f^{-1})_*\lambda^p(A)=|\det f|\cdot \lambda^p(A)
        \]
    \end{thm}
    \begin{proof}
        Prove $\lambda^p(f([0,1]^p))=|\det f|$ by verifying multilinearity on $(f(\mathbf{e}_1),\dots, f(\mathbf{e}_p))$ (use continuity of the measure).
    \end{proof}
    \begin{crl}
        Let $M\in\mathbb{R}^{p\times p}$ and $G$ be the Gram matrix of the column vectors.
        \[
            \lambda^p(M([0,1]^p)=|\det M|=(\det G)^{1/2}
        \]
    \end{crl}
    \begin{crl}
        Both $\beta^p$ and $\lambda^p$ are motion invariant.
        % Every motion $f$ in $\mathbb{R}^p$ is both $\mathscr{B}^p/\mathscr{B}^p$- and $\mathscr{L}^p/\mathscr{L}^p$-measurable, and \[f_*\beta^p=\beta^p,\quad f_*\lambda^p=\lambda^p\]
        % and if $A\in\mathscr{B}^p$ (resp. $A\in\mathscr{L}^p$), then $f(A)\in\mathscr{B}^p$ (resp. $f(A)\in\mathscr{L}^p$) and
        % \[
        %     \lambda^p(f(A))=\lambda^p(A)
        % \]
    \end{crl}
    \begin{thm}
        Let $\eta^p$ be the $p$-dimensional Hausdorff outer measure. There exists $k_p\in \,]0,\infty[$ such that
        \[
            (\vol^p)^*=k_p\cdot \eta^p
        \]
    \end{thm}
    \begin{proof}
        The proof only uses the behavior of $(\vol^p)^*$ under translations and homotheties $x\mapsto \alpha x,\,\alpha>0$. Thus, this theorem can be used to prove the motion invariance of the Lebesgue measure as well.
    \end{proof}
\section{Existence of Non-Lebesgue-Measurable Sets}
    \begin{thm}
        Assuming the axiom of choice, every $B\in\mathscr{B}(\mathbb{R}^n)$ with $\lambda^n(B)>0$ contains a non-Lebesgue-measurable set with strictly postive outer Lebesgue measure.
    \end{thm}
    \begin{crl}
        Assuming the axiom of choice, every subset of $\mathbb{R}^n$ with strictly positive outer Lebesgue measure contains a non-Lebesgue-measurable set with strictly postive outer Lebesgue measure.
    \end{crl}
\section{Extended real-valued meaurable functions}
    Let $(E,\mathscr{E})$ be a measurable space and equip $\mathbb{R}$, $\er$ with their Borel $\sigma$-algebras. %$\mathscr{B}(\mathbb{R})$ and $\mathscr{B}(\er)$, resp.
\subsection{Criterion for measurability}
    \begin{thm}\label{criterion for measurability}
        A function $u:E\rightarrow \er$ is measurable if and only if any one of the following
        \[
            \{u>r\},\quad \{u\geqslant r\},\quad \{u<r\},\quad \{u\leqslant r\}
        \]
        is in $\mathscr{E}$ for all $r\in\mathbb{Q}$, and hence for all $r\in\er$.
    \end{thm}
    \begin{proof}
        Theorem \ref{Borel sigma algebra on ER}. Theorem \ref{measurable functions}.
    \end{proof}
    \begin{thm}\label{measurability for countable image}
        A function $u:E\rightarrow \er$ with an at most countable image $\im u$ is measurable if and only if each fiber of $u$ is in $\mathscr{E}$.
    \end{thm}
    \begin{thm}\label{criterion for measurability for subsets of the extended real line}
        For $F\subset \er$ a topological subspace of $\er$ equipped with its Borel $\sigma$-algebra, the inclusion map $\iota:F\rightarrow \er$ is measurable, and a function $u:E\rightarrow F$ is measurable if and only if $\iota\circ u:E\rightarrow \er$ is measurable.
    \end{thm}
    \begin{proof}
        Theorem \ref{measurability w.r.t. codomain restricted to a topological subspace}.
    \end{proof}
\subsection{Lattice structures and sequential closure}
    The set $\er^X$ of functions is partially ordered by
    \[
        f\leqslant g:= f(x)\leqslant g(x)\text{ for all }x\in X
    \]
    Recall the following definition from order theory:
    \begin{define}
        A \textit{lattice} is a poset $P$ where every two-element subset $\{a,b\}\subset P$ has both a join $a\vee b$ and a meet $a\wedge b$.
        A \textit{sublattice} of a lattice is a subset that is both $\vee$-stable and $\wedge$-stable.

        A \textit{complete lattice} is a poset where every subset has both a join and a meet.
    \end{define}
    \begin{thm}
        $\mathbb{R}$ and $\mathbb{R}^X$ are lattices. $\er$ and $\er^X$ are complete lattices.
    \end{thm}
    \begin{thm}\label{lattice}
        If $(u_n)_{n\in\mathbb{N}}$ is a sequence of measurable functions $X\rightarrow \er$, then
        \[
            \sup_{n\in\mathbb{N}}u_n,\,\,\inf_{n\in\mathbb{N}}u_n,\,\,\varliminf_{n\in\mathbb{N}}  u_n,\,\,\varlimsup_{n\in\mathbb{N}}u_n
        \]
        are also measurable. In particular, $\cat{Mes}(X,\er)$ is a sublattice of $\er^X$.
    \end{thm}
    \begin{proof}
        For all $a\in\mathbb{R}$, $\left\{\sup_{n\in\mathbb{N}}u_n>a\right\}=\bigcup_{n\in\mathbb{N}}\{u_n>a\}\in\mathscr{X}$ and $\{-u_n>a\}=\{u_n<-a\}\in\mathscr{X}$. By Theorem \ref{criterion for measurability}, $\sup_{n\in\mathbb{N}}u_n$ and $-u_n$ are measurable. But $\inf_{n\in\mathbb{N}}u_n=\sup_{n\in\mathbb{N}}(-u_n)$.
    \end{proof}
    \begin{thm}\label{sequential closure of measurable functions}
        In the topology of pointwise convergence, $\cat{Mes}(X,\er)$ is sequentially closed in $\er^{X}$.
    \end{thm}
\subsection{Real-valued measurable functions and simple functions}
    \begin{thm}\label{subalgebra and sublattice structure of real measurable functions}
        $\cat{Mes}(X,\mathbb{R})$ is a sub-$\mathbb{R}$-algebra and a sublattice of $\mathbb{R}^X$.
    \end{thm}
    \begin{proof}
        Since addition and multiplication $+,\cdot:\mathbb{R}^2\rightarrow\mathbb{R}$ are continuous and hence Borel-measurable, if $u,v:X\rightarrow\mathbb{R}$ are $\mathscr{X}/\mathscr{B}(\mathbb{R})$-measurable, then $(u,v)$ is $\mathscr{X}/\mathscr{B}(\mathbb{R}^2)$-measurable (by Theorem \ref{Borel sigma algebra of product space}) and $u+v,\,u\cdot v$ as composition of measurable functions, are measurable. If $u,v\in \cat{Mes}(X,\mathbb{R})$, then by Theorem \ref{criterion for measurability for subsets of the extended real line}, $\iota \circ u,\iota\circ v\in\cat{Mes}(X,\er)$, where $\iota:\mathbb{R}\rightarrow\er$ is the inclusion map. By Theorem \ref{lattice}, $\iota\circ(u\vee v)=(\iota\circ u)\vee(\iota\circ v)\in\cat{Mes}(X,\er)$, and hence $u\vee v\in \cat{Mes}(X,\mathbb{R})$.
    \end{proof}
    \begin{define}
        The \textit{indicator function} of a subset $A\subset X$ is
        \[
            \mathbbm{1}_A:X\rightarrow \mathbb{R},\,\, x\mapsto\begin{cases}
                0&\text{if }x\notin A\\
                1&\text{if }x\in A
            \end{cases}
        \] 
    \end{define}
    \begin{thm}
        $\mathbbm{1}_A\in\cat{Mes}(X,\mathbb{R})$ if and only if $A\in\mathscr{X}$.
    \end{thm}
    \begin{proof}
        Theorem \ref{measurability w.r.t. codomain restricted to a topological subspace}. Theorem \ref{measurability for countable image}.
    \end{proof}
    \begin{define}
        A \textit{simple function} is a measurabe function $e:X\rightarrow \mathbb{R}$ such that $e(X)$ is a finite set.
    \end{define}
    \begin{thm}\label{subalgebra and sublattice structure of simple functions}
        Simple functions on $X$ form a sub-$\mathbb{R}$-algebra and a sublattice of $\cat{Mes}(X,\mathbb{R})$.
        %, that is, for    all $f,g\in \mathcal{E}(\mathscr{A};\mathbb{R})$
        %\begin{enumerate}
        %    \item $f\pm g\in \mathcal{E}(\mathscr{A};\mathbb{R})$.
        %    \item $f\cdot g\in \mathcal{E}(\mathscr{A};\mathbb{R})$.
        %    \item $f\vee g,\,f\wedge g\in \mathcal{E}(\mathscr{A};\mathbb{R})$.
        %\end{enumerate}
    \end{thm}
    \begin{proof}
        For example, if $e_1,e_2:X\rightarrow\mathbb{R}$ are simple functions, then $e_1+e_2, e_1\wedge e_2$ are measurable by Theorem \ref{subalgebra and sublattice structure of real measurable functions} and $(e_1+e_2)(X)\subset e_1(X)+e_2(X), (e_1\wedge e_2)(X)\subset e_1(X)\cup e_2(X)$ are finite sets.
    \end{proof}
    \begin{thm}
        A function $e:X\rightarrow \mathbb{R}$ is simple if and only if it is a linear combination of measurable indicator functions.
    \end{thm}
    \begin{proof}
        Sufficiency follows from Theorem \ref{subalgebra and sublattice structure of simple functions} since indicator functions are simple. Necessity. Suppose $e$ is simple. Then $e=\sum_{y\in e(X)}y\mathbbm{1}_{e^{-1}(\{y\})}$ is a finite sum. That $e^{-1}(\{y\})\in\mathscr{X}$ follows from Theorem \ref{measurability w.r.t. codomain restricted to a topological subspace} and Theorem \ref{measurability for countable image}.
    \end{proof}
    In view of Theorem \ref{criterion for measurability for subsets of the extended real line}, the set $\cat{Mes}(E,\mathbb{R})$ can be identified as a subset of $\cat{Mes}(E,\er)$, via the inclusion map $\iota:\mathbb{R}\rightarrow\er$.
\subsection{The sombrero lemma}
    \begin{thm}\label{sombrero}
        Every positive measurable function $u:X\rightarrow \er$ is the pointwise, and uniform if $u$ is bounded, limit of an increasing sequence $(e_n)_{n\in\mathbb{N}}$ of positive simple functions, i.e., $0\leqslant e_n\uparrow u$.
    \end{thm}
    \begin{proof}
        For each $n\in\mathbb{N}$ define
        \[
            A^n_k:=\begin{cases}
                \{{k}/{2^n}\leqslant u<{(k+1)}/{2^n}\}\,\, &\text{for }0\leqslant k\leqslant n2^n-1\\
                \{n\leqslant u\leqslant\infty\}\,\,&\text{for }k=n2^n
         \end{cases}\quad e_n(x):=\sum_{k=0}^{n2^n}\frac{k}{2^n}\mathbbm{1}_{A^n_k}(x)\geqslant 0
        \]
        Then $e_n$ is a simple function for all $n\in\mathbb{N}$, and $|e_n(x)-u(x)|\leqslant 2^{-n}$ for all $x\in \{0\leqslant u<n\}$.
    \end{proof}
\subsection{Positive and negative parts}
    \begin{define}
        The \textit{positive and negative parts} of a function $f:X\rightarrow \er$ are respectively
        \[
            f^+:=f\vee 0,\quad f^-:=-(f\wedge 0)=(-f)^+
        \]
    \end{define}
    \begin{thm}\label{decomposition into positive and negative parts}
        For any function $f:X\rightarrow \er$,
        \[
            f=f^+-f^-,\quad |f|=f^++f^-
        \]
    \end{thm}
    \begin{proof}
        At each $x\in X$, at least one of $f^+(x)$ and $f^-(x)$ is $0$. Thus, $f^+-f^-$ is always defined.
    \end{proof}
    \begin{thm}\label{positive parts preserve order relation while negative parts reverse}
        Let $f,g:X\rightarrow \er$.
        \[
            f(x)\leqslant g(x)\Longleftrightarrow (f^+(x)\leqslant g^+(x))\wedge (f^-(x)\geqslant g^-(x))
        \]
    \end{thm}
    \begin{proof}
        If $f(x)> 0$ then $g(x)\geqslant f(x)> 0$, in which case, $g^+(x)=g(x)\geqslant f(x)=f^+(x)$. 
        
        If $f(x)\leqslant 0$, then $f^+(x)=0$ and $f^+(x)\leqslant g^+(x)$ holds trivially regardless of $g(x)$.
        
        If $g(x)< 0$, then $f(x)\leqslant g(x)<0$, in which case, $f^-(x)=-f(x)\geqslant -g(x)=g^-(x)$.
        
        If $g(x)\geqslant 0$, then $g^-(x)=0$ and $f^-(x)\geqslant g^-(x)$ holds trivially regardless of $f(x)$.
    \end{proof}
    \begin{thm}\label{exsombrero}
        A function $u:X\rightarrow \er$ is measurable if and only if $u^+,u^-$ are measurable, in which case, $|u|$ is also measurable, and there exists a sequence $(e_n)_{n\in\mathbb{N}}$ of simple functions such that
        \begin{tight_enumerate}
            \item $e_n\to u$ pointwise,
            \item $|e_n|\uparrow |u|$ pointwise,
            \item $e_n\cdot u\geqslant 0$ for all $n\in\mathbb{N}$.
        \end{tight_enumerate}
        and the convergence is uniform if $u$ is bounded.

        In particular, the set of simple functions is sequentially dense, and hence dense, in $\cat{Mes}(X,\er)$.
    \end{thm}
    \begin{proof}
        By Theorem \ref{lattice}, $u^+,u^-$ are by definition measurable. $|\ph|:\er\rightarrow \er$ is continuous and hence Borel-measurable. Thus, as a composition, $|u|$ is measurable. If $u^+,u^-$ are measurable, then there exist increasing sequences $(e_n^+)_{n\in\mathbb{N}}$ and $(e_n^-)_{n\in\mathbb{N}}$ of positive simple functions with $e_n^+\uparrow u^+$ and $e_n^-\uparrow u^-$ pointwise. Then $(e_n:=e_n^+-e_n^-)_{n\in\mathbb{N}}$ is a sequence of simple functions that converges pointwise to $u$. Indeed, for $x\in \{u\geqslant 0\}$, $u^-(x)=0$, $e_n^-(x)=0$ for all $n$ and hence $e_n(x)=e_n^+(x)\uparrow u^+(x)=u(x)$; for $x\in \{u< 0\}$, $u^+(x)=0$, $e_n^+(x)=0$ for all $n$ and hence $e_n(x)=-e_n^-(x)\downarrow -u^-(x)=u(x)$.
    \end{proof}
    \begin{thm}
        If $u,v:X\rightarrow \er$ are measurable, then
        \begin{tight_enumerate}
            \item $u+v$ is measurable, provided that $u(x)+v(x)$ is defined for every $x\in X$.
            \item $u\cdot v$ is measurable.
        \end{tight_enumerate} 
    \end{thm}
    \begin{proof}
        By Theorem \ref{exsombrero}, there exist sequences of simple functions $f_n\to u$ and $g_n\to v$ pointwise with $|f_n|\leqslant |u|$ and $|g_n|\leqslant |v|$. By Theorem \ref{subalgebra and sublattice structure of real measurable functions}, $f_n+g_n$, $f_n\cdot g_n$ are measurable. For \textit{1}, if $u(x)+v(x)$ is defined for every $x\in X$, then it can be shown that $f_n(x)+g_n(x)$ always tend to $u(x)+v(x)$ for all $x\in X$. For \textit{2}, if $u(x)=0$, then $f_n(x)=0$ for all $n\in\mathbb{N}$. Thus, $f_n(x)\cdot g_n(x)\to u(x)\cdot v(x)$, since we have defined  $0\cdot \pm\infty=0$. Theorem \ref{sequential closure of measurable functions}.
    \end{proof}
    \begin{thm}
        If $u,v:X\rightarrow \er$ are measurable, then
        \[
            \{u\leqslant v\},\,\{u<v\},\,\{u=v\},\,\{u\neq v\}\in\mathscr{X}
        \]
    \end{thm}
    \begin{proof}
        Let $(e_n)_{n\in\mathbb{N}}$ be a sequence of positive simple functions with $e_n\uparrow u^+$. Then $\{u^+\leqslant v^+\}=\bigcap_{n\in\mathbb{N}}\{e_n\leqslant v^+\}=\bigcap_{n\in\mathbb{N}}\{0\leqslant v^+-e_n\}\in\mathscr{X}$. Similarly, $\{u^-\geqslant v^-\}\in\mathscr{X}$. By Theorem \ref{positive parts preserve order relation while negative parts reverse}, $\{u\leqslant v\}=\{u^+\leqslant v^+\}\cap \{u^-\geqslant v^-\}\in\mathscr{X}$.
    \end{proof}
    % \begin{thm}\label{limit of sum of two sequences}
    %     Let $(a_n),(b_n)\in\er^\mathbb{N}$ be convergent sequences with
    %     \[
    %         \lim_{n\to\infty} a_n,\lim_{n\to\infty} b_n\in\,\,]\!-\!\infty,\infty]\quad\text{or}\quad \lim_{n\to\infty} a_n,\lim_{n\to\infty} b_n\in [-\infty,\infty[
    %     \]
    %     Then $\lim_{n\to\infty}a_n+b_n=\lim_{n\to\infty} a_n+\lim_{n\to\infty} b_n$.
    % \end{thm}
    % \begin{thm}\label{operations under which the set of measurable functions are closed}
    %     Let $u,v\in \mathcal{M}(\mathscr{A};\overline{\mathbb{R}})$.
    %     \begin{tight_enumerate}
    %         \item $u\pm v\in \mathcal{M}(\mathscr{A};\overline{\mathbb{R}})$, whenever defined.
    %         \item $u\cdot v\in \mathcal{M}(\mathscr{A};\overline{\mathbb{R}})$.
    %     \end{tight_enumerate}
    % \end{thm}
    % \begin{proof}
    %     \begin{tight_enumerate}
    %         %\item Write $u=u^+-u^-$ and observe that $u^+,u^-\in \mathcal{M}(\mathscr{A};[0,\infty])$. Apply   sombrero lemma.
    %         %\item By \autoref{exsombrero}, limit laws and \autoref{lattice}.
    %         %\item %$0\in \mathcal{M}(\mathscr{A};\overline{\mathbb{R}})$ and $\mathcal{M}(\mathscr{A};\overline {\mathbb{R}})$ is a sublattice by \autoref{lattice}. 
    %         %Observe that $u=u^{+}-u^{-},\,|u|=u^++u^-$.
    %     \end{tight_enumerate}
    % \end{proof}
\subsection{Monotone class theorem}
    \begin{define}
        A function $f:X\rightarrow \er$ is \textit{bounded} if $\sup_{x\in X}|f(x)|<\infty$.
    \end{define}
    \begin{thm}\label{monotone class theorem}
        Let $M\subset \er^X$ be such that
        \begin{tight_enumerate}
            \item $1\in M$,
            \item $\alpha u+\beta v\in M$ for all $u,v\in M$ bounded and $\alpha,\beta\in\mathbb{R}$,
            \item $u_n\uparrow u\Rightarrow u\in M$ for any increasing sequence $(u_n)_{n\in\mathbb{N}}$ of positive functions from $M$.
        \end{tight_enumerate}
        If $\mathbbm{1}_A\in M$ for all $A$ in a $\cap$-stable generator of $\mathscr{X}$,
        then $M$ contains all positive measurable functions and all bounded measurable functions.
    \end{thm}
    \begin{proof}
        The system $D=\{A\in\mathscr{X}\,|\,\mathbbm{1}_A\in M\}$ is a Dynkin system (Theorem \ref{characterization of Dynkin systems via monotone classes}) that contains a $\cap$-stable generator of the $\sigma$-algebra $\mathscr{X}$. By Theorem \ref{generators}, $\mathscr{X}\subset D$, and hence, $\mathscr{X}=D$. By \textit{2}, $M$ contains all simple functions. By \textit{3} and the sombrero lemma, $M$ contains all positive measurable functions. For $f$ any bounded measurable function, $f^+,f^-$ are bounded positive measurable functions.
    \end{proof}
\subsection{Measurability w.r.t. the \texorpdfstring{$\sigma$}{sigma}-algebra generated by a function}
    \begin{thm}\label{measurability w.r.t. the sigma algebra generated by a function}
        Let $X:E\rightarrow F$ be a function. A function $u:E\rightarrow \er$ is in $\overline{\sigma X}$ if and only if
        \[
            u=f\circ X
        \]
        for some $f\in \overline{\mathscr{F}}$.
    \end{thm}
    \begin{proof}
        Sufficiency is clear. For necessity, consider indicators of sets in $\sigma X$ and use sombrero lemma.
    \end{proof}
    \begin{thm}
        Let $(X_t:E\rightarrow F_t)_{t\in T}$ be a family of measurable functions. Let $X:E\rightarrow \prod_{t\in T}{F_t}$ be the induced measurable function. A function $u:E\rightarrow \er$ is in $\overline{\sigma X}=\overline{\sigma(X_t,t\in T)}$ if and only if there exists a sequence $(t_n)\in T^\mathbb{N}$ and $f\in \overline{\bigotimes_{n\in\mathbb{N}}{\mathscr{F}_{t_n}}}$ such that
        \[
            u=f(X_{t_1},X_{t_2},\dots)
        \]
    \end{thm}
\section{Countably generated measurable spaces}
    Let $X$ be a measurable space.
\subsection{Separating points}
    \begin{thm}
        A $\sigma$-algebra $\mathscr{A}$ on $X$ separates points of $X$ if and only if a generator of $\mathscr{A}$ separates points of $X$.
    \end{thm}
    \begin{thm}
        A measurable space $(X,\mathscr{A})$ is countably generated if and only if there exists a measurable function $f:(X,\mathscr{A})\rightarrow ([0,1],\mathfrak{B})$ such that $f^{-1}\mathfrak{B}[0,1]=\mathscr{A}$.
    \end{thm}
    \begin{thm}
        A measurable space $(X,\mathscr{A})$ is separated and countably generated if and only if there exists $M\subset [0,1]$ such that $(X,\mathscr{A})$ is isomorphic to $(M,\mathfrak{B})$.
    \end{thm}
    \begin{thm}
        There exists a minimal generator of $\mathfrak{B}^1$:
        \[
            \bigcup_{n\in\mathbb{Z}}\{](k-1)\cdot 2^{-n},k\cdot 2^{-n}[\,|\,k\in\mathbb{Z}\}\mathfrak{}
        \]
    \end{thm}
\section{Further Theorems}
    \begin{thm}
        Let $(X,\mathscr{A}),(X',\mathscr{A}')$ be measurable spaces. Let $T:X\rightarrow X'$.
        \[u\in\mathcal{M}(\sigma(T);\overline{\mathbb{R}})\Longleftrightarrow u=w\circ T\text{ for a function } w\in \mathcal{M}(\mathscr{A}';\overline{\mathbb{R}})\]
    \end{thm}
    \begin{thm}
        Let $\mathscr{G}\subset \mathscr{P}(X)$ be closed under finite intersection. Let $V$ be a vector    subspace of $\mathbb{R}^X$ such that 
        \begin{tight_enumerate}
            \item $\mathbbm{1}_X\in V$ and $\mathbbm{1}_G\in V$ for all $G\in\mathscr{G}$.
            \item for every non-negative increasing sequence $(u_n)\in V^\mathbb{N}$, $\sup\limits_{n\in\mathbb {N}}u_n\in \mathbb{R}^X\Longrightarrow \sup\limits_{n\in\mathbb{N}}u_n\in  V$.
        \end{tight_enumerate}
        then $\mathcal{M}(\sigma(\mathscr{G});\mathbb{R})\subset V$.
    \end{thm}
\section{Vector- and Complex-Valued Measurable Functions}
    \begin{thm}
        Let $(X,\mathscr{A})$ be a measurable space. A function $f:X\rightarrow \mathbb{R}^p$ is $\mathscr{A}/\mathfrak{B}^p$-measurable if and only if $f^i:X\rightarrow \mathbb{R}$ is $\mathscr{A}/\mathfrak{B}$-measurable for all $i=1,\dots,p$.
    \end{thm}
    \begin{proof}
        Universal property of product $\sigma$-algebra?
    \end{proof}
    In general, if $Y$ is a normed space, and $f\in\cat{Mes}(\mathscr{A},\mathscr{B}(Y))$, then $\|f\|\in\cat{Mes}(\mathscr{A})$ since $\|\|$ is continuous and hence measurable.
    
    The homeomorphism between $\mathbb{C}$ and $\mathbb{R}^2$
    \[\psi:\mathbb{C}\rightarrow \mathbb{R}^2,\,\,z\mapsto (\RE{z},\IM{z})\quad\text{and}\quad\psi^{-1}:\mathbb{R}^2\rightarrow \mathbb{C},\,\,(x,y)\mapsto x+iy\]
    entails that $\mathscr{B}(\mathbb{C})=\psi^{-1}[\mathscr{B}(\mathbb{R}^2)]$.%$C\in\mathscr{B}(\mathbb{C})\Longleftrightarrow \psi[C]\in\mathscr{B}(\mathbb{R}^2)$.
    \begin{crl}
        \begin{align*}
            f\in\mathcal{M}(\mathscr{A};\mathbb{C})&\Longleftrightarrow \RE{f},\IM{f}\in\mathcal{M}(\mathscr{A};\mathbb{R})\\
            &\Longleftrightarrow (\RE{f})^+,(\IM{f})^+,(\RE{f})^-,(\IM{f})^-\in\mathcal{M}(\mathscr{A};\mathbb{R}_+)
        \end{align*}
    \end{crl}
    \begin{thm}
        If $f,g:(X,\mathscr{A})\rightarrow (\mathbb{C},\mathfrak{B})$ are measurable and $\alpha,\beta\in\mathbb{C}$, then
        \[
            \alpha f+\beta g,\,\,f\cdot g,\,\,|f| 
        \]
        are measurable.
    \end{thm}
\section{Random variables}
    Let $(\Omega,\mathscr{H},\mathbb{P})$ be a probability space.
    \begin{define}
        A \textit{random variable on $(\Omega,\mathscr{H},\mathbb{P})$} is a morphism $X:(\Omega,\mathscr{H})\rightarrow (E,\mathscr{E})$ in \cat{Mes}; the pushforward measure $X_*\mathbb{P}$ is called the \textit{distribution} of $X$; the distribution of a stochastic process $(\Omega,\mathscr{H})\rightarrow (E,\mathscr{E})^T$ is called its \textit{probability law}.

        A \textit{simple random variable} is a random variable that takes only finitely many values. A \textit{discrete random variable} is a random variable that takes at most countably many values.
    \end{define}
    \begin{thm}
        The distribution of any random variable is a probability measure on its target space.
    \end{thm}
\subsection{Joint distributions}
    \begin{define}
        The \textit{joint distribution} of a family $(X_t:(\Omega,\mathscr{H})\rightarrow (E_t,\mathscr{E}_t))_{t\in T}$ of random variables on $(\Omega,\mathscr{H},\mathbb{P})$ is the distribution of the induced random variable $X:(\Omega,\mathscr{H})\rightarrow (\prod_{t\in T}E_t,\bigotimes_{t\in T}\mathscr{E}_t),\,\,\omega\mapsto (X_t(\omega))_{t\in T}$.
    \end{define}
    \begin{thm}
        If the joint distribution of two random variables $X:(\Omega,\mathscr{H})\rightarrow (E,\mathscr{E})$ and $Y:(\Omega,\mathscr{H})\rightarrow (F,\mathscr{F})$ has the form of a product measure $\mu\times\nu$ of two probability measures on $\mathscr{E}$ and $\mathscr{F}$ respectively, then $\mu$ is necessarily the distribution of $X$ and $\nu$ necessarily the distribution of $Y$. 
    \end{thm}
    \begin{proof}
        $\mu(A)=(\mu\times \nu)(A\times F)=(X,Y)_*\mathbb{P}(A\times F)=(X,Y)_*\mathbb{P}(\pi_1^{-1}(A))=X_*\mathbb{P}(A)$.
    \end{proof}
\section{Standard Borel spaces}
    \begin{thm}
        The product of standard Borel spaces is standard Borel.
    \end{thm}
\chapter{The Lebesgue Integral}
    Let $(E,\mathscr{E},\mu)$ be a measure space.
\section{Integration of positive measurable functions}
    Let $S_+$ be the set of positive simple functions on $(E,\mathscr{E})$.
\subsection{Definition via simple functions}
    \begin{define}
        Let
        \[
            I_\mu:S_+\rightarrow [0,\infty],\,\, f\mapsto \sum_{y\,\in\,\im f} y\cdot \mu(f^{-1}\{y\})
        \]
    \end{define}
    \begin{thm}\label{properties of I_mu}
        For $f,g\in S_+$, $\alpha\in\mathbb{R}_+$, $A\in\mathscr{E}$,
        \begin{gather*}
            I_{\mu}(\mathbbm{1}_A)=\mu(A)\\
            I_{\mu}(\alpha f)=\alpha\cdot I_{\mu}(f)\\
            I_{\mu}(f+g)=I_{\mu}(f)+I_{\mu}(g)\\
            f\leqslant g\Longrightarrow I_{\mu}(f)\leqslant I_{\mu}(g)
        \end{gather*}
    \end{thm}
    \begin{proof}
        1. For 1 to hold it is necessary to define $0\cdot  \infty=0$, since $I_\mu(\mathbbm{1}_A)=1\cdot \mu(A)+0\cdot \mu(A^c)$.
    
        3. $I_\mu(f+g)=\sum_{(y,y')\in\im f\times \im g}(y+y')\mu(f^{-1}\{y\}\cap g^{-1}\{y'\})=\sum_y(y\cdot \sum_{y'}\mu(f^{-1}\{y\}\cap g^{-1}\{y'\}))+\sum_{y'}(y'\cdot \sum_{y}\mu(f^{-1}\{y\}\cap g^{-1}\{y'\}))=\sum_y(y\cdot \mu(f^{-1}\{y\}))+\sum_{y'}(y'\cdot \mu(g^{-1}\{y'\}))$.

        4. Since $g=f+(g-f)$ with $g-f\in S_+$, by 3, $I_\mu(g)=I_\mu(f)+I_\mu(g-f)\geqslant I_\mu(f)$.
    \end{proof}
    The proof of Theorem \ref{properties of I_mu} uses only finite additivity, i.e., Theorem \ref{properties of I_mu} holds for $\mu$ any content.
    \begin{define}\label{definition of integral of positive measurable functions}
        The \textit{$\mu$-integral} of a positive $\mathscr{E}/\mathscr{B}(\er_+)$-measurable function $u:E\rightarrow \er_+$ is
        \[
            \int_\mu u:=% \int u(x)\mu(\odif x):=
            \sup_{S_+ \ni f\leqslant u}I_{\mu}(f)% \in [0,\infty]
        \]
    \end{define}
    \begin{thm}
        For $f\in S_+$,
        \[
            \int_\mu f=I_\mu (f)
        \]
    \end{thm}
    % \begin{proof}
    %     This follows from the monotonicity of $I_\mu$ (Theorem \ref{properties of I_mu}).
    % \end{proof}
\subsection{Isotonicity}
    \begin{thm}\label{monotonicity of integral of positive measurable functions}
        Let $u,v\in\overline{M}_+$.
        \[
            u\leqslant v\Longrightarrow \int_\mu u\leqslant \int_\mu v
        \]
    \end{thm}
    \begin{proof}
        If $u\leqslant v$, then $\{f\in S_+\,|\,f\leqslant u\}\subset \{f\in S_+\,|\,f\leqslant v\}$. By Definition \ref{definition of integral of positive measurable functions}, $\int_\mu u\leqslant \int_\mu v$.
    \end{proof}
\subsection{Beppo Levi's theorem}
    \begin{thm}\label{Beppo Levi's Theorem}
        Let $(u_n)_{n\in\mathbb{N}}$ be an increasing sequence of positive measurable functions $X\rightarrow \er_+$.
        \[
            u_n\uparrow u\Longrightarrow \int_\mu u_n\uparrow \int_\mu u
        \]
        where the convergence $u_n\uparrow u$ is meant, as usual, pointwise.
    \end{thm}
    \begin{proof}
        By Theorem \ref{sequential closure of measurable functions}, $u$ is positive measurable and $\int_\mu u$ is defined. By isotonicity of the integral, $\sup_{n}\int_\mu u_n\leqslant \int_\mu u$. It remains to show that if $S_+ \ni f\leqslant u=\sup_{n}u_n$, then $I_\mu(f)\leqslant \sup_{n}\int_\mu u_n$. Let $0<\alpha<1$. Let $A_k=\{\alpha f\leqslant u_k\}$. Since $\alpha f\cdot \mathbbm{1}_{A_k}\leqslant u_k$, by Definition \ref{definition of integral of positive measurable functions}, $I_\mu(\alpha f\cdot \mathbbm{1}_{A_k})\leqslant \int_\mu u_k \leqslant  \sup_{n}\int_\mu u_n$. Now $\alpha I_\mu(f)=I_\mu(\alpha f)=I_\mu(\alpha f\cdot \mathbbm{1}_{A_k})+I_\mu(\alpha f\cdot \mathbbm{1}_{A_k^c})\leqslant \sup_{n}\int_\mu u_n+\alpha y_{\text{max}}\mu(A_k^c)$, where $y_\text{max}=\max\im f$. As $k\to \infty$, $A_k^c\downarrow \varnothing$. % for if $x\in\bigcap_k A_k^c$ then $u_k(x)<\alpha f(x)$ for all $k$ and $u(x)=\sup_k u_k(x)\leqslant \alpha f(x)<f(x)\leqslant u(x)$, contradiction. 
        By continuity from above $\mu(A_k^c)\downarrow 0$. Thus, $\alpha I_\mu(f)\leqslant \sup_{n}\int_\mu u_n$, and $I_\mu(f)=\sup\,\{\alpha I_\mu(f)\,|\,0<\alpha<1\}\leqslant \sup_{n}\int_\mu u_n$.
    \end{proof}
\subsection{Vanishing integral \texorpdfstring{$\Longleftrightarrow$}{<->} vanishes a.e.}
    \begin{thm}\label{null integral}
        Let $u:X\rightarrow \er_+$ be a positive measurable function.
        \[
            \int_\mu u=0\Longleftrightarrow \mu(\{u>0\})=0
        \]
    \end{thm}
    \begin{proof}
        Assume $\int_\mu u=0$. By the sombrero lemma, $(S_+)^\mathbb{N}\ni (f_n)\uparrow u$. By isotonicity of the integral, $I_\mu(f_n)=0$ and hence $\mu(\{f_n>0\})=0$. Since $\{f_n>0\}\uparrow \{u>0\}$, by continuity from below, $\mu(\{u>0\})=0$. Conversely, let $N=\{u>0\}$. Then $n\cdot \mathbbm{1}_N\uparrow \infty\cdot \mathbbm{1}_N\geqslant u$ and $I_\mu(n\cdot \mathbbm{1}_N)\uparrow \int_\mu\infty\cdot \mathbbm{1}_N\geqslant \int_\mu u$, but $I_\mu(n\cdot \mathbbm{1}_N)=n\cdot \mu(N)=0$ for all $n$.
    \end{proof}
    \begin{thm}
        Let $u:X\rightarrow \er_+$ be a strictly positive measurable function. If $\mu$ is non-trivial, then
        \[
            \int_\mu u>0
        \]
    \end{thm}
    \begin{proof}
        Assume $u>0$. If $\int_\mu u=0$, then $\mu(X)=\mu(\{u>0\})=0$ by Theorem \ref{null integral}.
    \end{proof}
\subsection{Positive homogeneity and \texorpdfstring{$\sigma$}{sigma}-additivity}
    \begin{thm}\label{properties of integral of positive measurable functions}
        For $u,v:X\rightarrow \er_+$ positive measurable and $\alpha\in [0,\infty]$ holds the following:
        \begin{gather*}
            \int_\mu \alpha u=\alpha \int_\mu u\\
            \int_\mu u+v=\int_\mu u+\int_\mu v
        \end{gather*}
    \end{thm}
    \begin{proof}
        Let $\alpha=\infty$. The proof of Theorem \ref{null integral} shows that $n\cdot \mu(\{u>0\})\uparrow \int_\mu\infty\cdot u$, i.e., $\int_\mu \infty\cdot u=\infty\cdot \mu(\{u>0\})$, which by Theorem \ref{null integral} is the same as $\infty\cdot \int_\mu u$.
    \end{proof}
    \begin{thm}\label{integral of a series}
        For $(u_n)_{n\in\mathbb{N}}$ a sequence of positive measurable functions $X\rightarrow \er_+$, the series
        \[
            \sum_{n=1}^\infty u_n
        \]
        is positive measurable, and
        \[
            \int_\mu \sum_{n=1}^\infty u_n=\sum_{n=1}^\infty \int_\mu u_n
        \]
    \end{thm}
    % \begin{proof}
    %     By Beppo Levi's theorem and additivity of the integral.
    % \end{proof}
\subsection{Insensitivity}
    \begin{thm}\label{integrals of a.e. equal positive measurable functions}
        Let $u,v:X\rightarrow [0,\infty]$ be measurable.
        \[
            u=v\,\,\,\mu\text{-a.e.}\Longrightarrow \int_\mu u=\int_\mu v
        \]
    \end{thm}
    \begin{proof}
        By $u=v$ $\mu$-a.e. is meant $u|_F=v|_F$ for a $\mu$-full set $F$. By Theorem \ref{properties of integral of positive measurable functions}, $\int u\odif\mu=\int u\one_F\odif\mu+\int u\one_{F^c}\odif\mu$. By Theorem \ref{null integral}, $\int u\one_{F^c}\odif\mu=0$. Thus, $\int u\odif\mu=\int u\one_F\odif\mu=\int v\one_F\odif\mu=\int v\odif\mu$.
    \end{proof}
\subsection{Measure density functions}
    \begin{thm}
        Let $u:X\rightarrow \er_+$ be a positive measurable function. The set function
        \[
            u\odot \mu:\mathscr{X}\rightarrow [0,\infty],\,\,A\mapsto \int_\mu u\cdot \mathbbm{1}_A 
        \]
        is a measure on $\mathscr{X}$, and
        \[
            \{\mu=0\}\subset\{u\odot\mu=0\}
        \]
    \end{thm}
    \begin{proof}
        By Theorem \ref{integral of a series} to show $\sigma$-additivity or by Beppo Levi's theorem to use Theorem \ref{characterization of premeasures by monotone convergence}.
    \end{proof}
    \begin{thm}
        Let $u,v:X\rightarrow \er_+$ be positive measurable functions.
        \[
            \int_{u\odot \mu} v=\int_\mu u\cdot v
        \]
        In particular, $v\odot (u\odot \mu)=(v\cdot u)\odot \mu$.
    \end{thm}
    \begin{proof}
        By linearity, the equality clearly holds for $v$ any simple function. It remains to apply the sombrero lemma and Beppo Levi's theorem.
    \end{proof}
\subsection{Fatou's lemma}
    \begin{thm}
        Let $(u_n)_{n\in\mathbb{N}}$ be a sequence of positive measurable functions $X\rightarrow \er_+$.
        \[
            \int_\mu \varliminf_{n\to\infty} u_n \leqslant \varliminf_{n\to\infty} \int_\mu u_n
        \]
    \end{thm}
    \begin{proof}
        
    \end{proof}
\subsection{Characterization of measures via integrals}
    \begin{thm}\label{Characterization of measures via integrals}
        Let $L:\cat{Mes}(X,\er_+)\rightarrow \er_+$ be a function. If
        \begin{tight_enumerate}
            \item $L(0)=0$,
            \item $L(\alpha u+\beta v)=\alpha L(f)+\beta L(g)$ for all $u,v\in\cat{Mes}(X,\er_+)$ and $\alpha,\beta\in\mathbb{R}_+$,
            \item $u_n\uparrow u\Rightarrow L(u_n)\uparrow L(u)$ for all increasing sequence $u_\bullet\in\cat{Mes}(X,\er_+)^\mathbb{N}$.
        \end{tight_enumerate}
        then there exists a unique measure $\mu$ on $\mathscr{X}$ such that
        \[
            L=\int_\mu
        \]
    \end{thm}
    \begin{proof}
        If $\mu$ is a such a measure, then it is uniquely determined: $\mu(A)=\int_\mu\mathbbm{1}_A=L(\mathbbm{1}_A)$ for all $A\in\mathscr{X}$. But $\mu:A\mapsto L(\mathbbm{1}_A)$ is indeed a measure (Theorem \ref{characterization of premeasures by monotone convergence}) that satisfies $L=\int_\mu$ for all measurable indicator functions, and by linearity, for all positive simple functions, and by the sombrero lemma and the monotone convergence theorem, for all positive measurable functions.
    \end{proof}
\subsection{Integration over sub-\texorpdfstring{$\sigma$}{sigma}-algebras}
    \begin{thm}
        If $\mathscr{E}_0$ is a sub-$\sigma$-algebra of $\mathscr{E}$, then for $u:E\rightarrow \er_+$ any positive $\mathscr{E}_0/\mathscr{B}$-measurable function, $u$ is also $\mathscr{E}/\mathscr{B}(\er_+)$-measurable, and
        \[
            \int_{\mu|_{\mathscr{E}_0}}u=\int_{\mu} u
        \]
    \end{thm}
    \begin{proof}
        Apply Theorem \ref{Characterization of measures via integrals} for $L:\cat{Mes}(E_0,\er_+)\rightarrow \er_+,\,\,v\mapsto \int_\mu v$ to obtain a measure $\nu$ on $\mathscr{E}_0$ such that $\nu(A_0)=\int_\nu \mathbbm{1}_{A_0}=L(\mathbbm{1}_{A_0})=\int_\mu\mathbbm{1}_{A_0}=\mu(A_0)$ for all $A_0\in\mathscr{E}_0$, i.e., $\nu=\mu|_{\mathscr{E}_0}$.
    \end{proof}
\subsection{The general transformation theorem}
    \begin{thm}\label{the general transformation theorem}
        Let $\begin{tikzcd}
            (E,\mathscr{E}) \ar[r, "f"] & (F,\mathscr{F}) \ar[r, "u"] & \er_+
        \end{tikzcd}$ be a diagram in \cat{Mes}.
        \[\int_\mu f^* u=\int_{f_*\mu} u\]
        where $f^* u=u\circ f$ is the pullback of $u$ by $f$.
    \end{thm}
    \begin{proof}
        Apply Theorem \ref{Characterization of measures via integrals} for $L:\cat{Mes}(F,\er_+)\rightarrow \er_+,\,\,u\mapsto \int_\mu f^*u$ to obtain a measure $\nu$ on $\mathscr{F}$ such that $\nu(A)=\int_\nu \mathbbm{1}_{A}=L(\mathbbm{1}_{A})=\int_\mu f^*\mathbbm{1}_{A}=\mu(\{f\in A\})$ for all $A\in\mathscr{F}$, i.e., $\nu=f_*\mu$.
    \end{proof}
\subsection{Integration over measurable subsets}
    \begin{thm}
        Let $A\in\mathscr{X}$ and $u:X\rightarrow \er_+$. Equip $A$ with the trace $\sigma$-algebra $\mathscr{X}|_A$. Then
        \[
            u\cdot\mathbbm{1}_A\in\cat{Mes}(X,\er_+)\Longleftrightarrow u|_A\in\cat{Mes}(A,\er_+)
        \]
        in which case,
        \[
            \int_\mu u\cdot\mathbbm{1}_A=\int_{\mu|_{\mathscr{X}|_A}} u|_A
        \]
        Moreover, for all $u\in \cat{Mes}(X,\er_+)$,
        \[
            \int_\mu u\cdot\mathbbm{1}_A=\int_{\mu_{\cap A}}u
        \]
    \end{thm}
    \begin{proof}
        For $1$, first $\Sigma X|_A\subset \Sigma X$ so $\mu|_{\Sigma X|_A}$ is a well-defined measure on $\Sigma X|_A$. Suppose $u\geqslant 0$ and thus $\tilde{u}\geqslant 0$. Let $(f_n)$ be a sequence of positive simple functions on $X$ such that $f_n\uparrow \tilde{u}$. Then $f_n|_A$ is a sequence of positive simple functions on $A$ with $f_n|_A\uparrow u$, since $\widetilde{f_n|_A}=f_n$. Thus, $\int \tilde{u}\odif\mu=\lim I_\mu(f_n)=\lim I_{\mu|_{\Sigma X|_A}}(f_n|_A)=\int u\odif\mu_{\Sigma X|_A}$ since $(f_n|_A)^{-1}\{y\}=(f_n)^{-1}\{y\}$ for $y>0$, that is, $I_\mu(f_n)=I_{\mu|_{\Sigma X|_A}}(f_n|_A)$. It remains to point out $\widetilde{u^\pm}=\tilde{u}^\pm$.

        Apply Theorem \ref{Characterization of measures via integrals} for $L:\cat{Mes}(X,\er_+)\rightarrow \er_+,\,\,u\mapsto \int_\mu u\cdot\mathbbm{1}_A$ to obtain a measure $\nu$ on $\mathscr{X}$ such that $\nu(A')=\int_\nu \mathbbm{1}_{A'}=L(\mathbbm{1}_{A'})=\int_\mu\mathbbm{1}_{A'}\cdot\mathbbm{1}_{A}=\mu(A'\cap A)$ for all $A'\in\mathscr{X}$, i.e., $\nu=\mu_{\cap A}$.
    \end{proof}
\subsection{Markov's inequality}
    \begin{thm}\label{Markov's inequality}
        Let $u:E\rightarrow\er_+$ be positive measurable. For $c>0$, %$A\in\mathscr{E}$,
        % \[
        %     \mu(\{|u|\geqslant c\}\cap A)\leqslant\frac{1}{c}\int_\mu |u|\cdot\mathbbm{1}_A
        % \]
        % and if $A=X$, in particular,
        \[
            \mu(\{u\geqslant c\})\leqslant\frac{1}{c}\int_\mu u
        \]
    \end{thm}
    \begin{thm}\label{Markov's inequality, generalized}
        Let $u:E\rightarrow \er$ be measurable. For $f:\er\rightarrow\er_+$ increasing and $c\in\mathbb{R}$ such that $f(c)>0$,
        \[
            \mu(\{u\geqslant c\})\leqslant\frac{1}{f(c)}\int_\mu f\circ u
        \]
    \end{thm}
    \begin{proof}
        $\{u\geqslant c\}\subset \{f\circ u\geqslant f(c)\}$ since $f$ is increasing. Thus, $\mu(\{u\geqslant c\})\leqslant \mu(\{f\circ u\geqslant f(c)\})\leqslant f(c)^{-1}\int_\mu f\circ u$ by Theorem \ref{Markov's inequality}.
    \end{proof}
\subsection{Integration w.r.t. a series of measures}
    \begin{thm}\label{integration w.r.t. a series of measures}
        Let $(\mu_n)_{n\in\mathbb{N}}$ be a sequence of measures on $(E,\mathscr{E})$. For $u\in\overline{\mathscr{E}}_+$,
        \[
            \int_{\sum_{n=1}^\infty \mu_n} u = \sum_{n=1}^\infty \int_{\mu_n} u
        \]
    \end{thm}
    \begin{proof}
        Write $\mu=\sum_{n=1}^{\infty}\mu_n$. For every positive
        simple measurable function $s$, the definition of the integral gives
        \[
        \int_\mu s=\sum_{n=1}^{\infty}\int_{\mu_n}s.
        \]
        Choose nonnegative simple measurable functions $s_k\uparrow u$.
        By the monotone convergence theorem,
        \[
        \begin{aligned}
        \int_\mu u
        &=\sup_k\int_\mu s_k\\
        &=\sup_k\sup_m\sum_{n=1}^{m}\int_{\mu_n}s_k\\
        &=\sup_m\sup_k\sum_{n=1}^{m}\int_{\mu_n}s_k\\
        &=\sup_m\sum_{n=1}^{m}\int_{\mu_n}u\\
        &=\sum_{n=1}^{\infty}\int_{\mu_n}u.
        \end{aligned}
        \]
        Here the two suprema commute because both equal the supremum
        over all pairs $(k,m)$; the penultimate equality follows from
        monotone convergence and passage of increasing limits through
        finite sums.
    \end{proof}
\section{Integrability}
    \begin{define}
        A function $u:E\rightarrow \er$ is \textit{$\mu$-quasi-integrable} if it is measurable and
        \[
            \int_\mu u^+<\infty\quad\vee\quad \int_\mu u^-<\infty
        \]
        in which case the \textit{$\mu$-integral} of $u$ is 
        \[
            \int_\mu u:=\int_\mu u^+-\int_\mu u^-
        \]
    \end{define}
    \begin{define}
        A function $u:E\rightarrow \er$ is \textit{$\mu$-integrable} if it is measurable and
        \[
            \int_\mu u^+<\infty\quad\wedge\quad \int_\mu u^-<\infty
        \]
    \end{define}
    \begin{thm}
        A measurable function $u:E\rightarrow \er$ is \textit{$\mu$-integrable} if and only if
        \[
            \int_\mu |u|<\infty
        \]
    \end{thm}
    \begin{thm}\label{integrable functions are a.e. real-valued}
        A measurable function $u:X\rightarrow \er$ is $\mu$-integrable if and only if $\{|u|=\infty\}$ is $\mu$-null and $u\cdot \mathbbm{1}_{\{|u|=\infty\}^c}:X\rightarrow \mathbb{R}$ is $\mu$-integrable. 
        % , in which case
        % \[
        %     \int u\odif \mu=\int u\cdot \mathbbm{1}_{\{|u|=\infty\}^c}\odif \mu
        % \]
    \end{thm}
    \begin{proof}
        Let $A=\{|u|=\infty\}$. If $u$ is integrable, then $\int |u\cdot \mathbbm{1}_{A^c}|\odif \mu\leqslant \int |u|\odif \mu<\infty$. By Theorem \ref{properties of integral of positive measurable functions} \& \ref{monotonicity of integral of positive measurable functions}, $\infty\cdot \mathbbm{1}_A\leqslant |u|\Longrightarrow \infty\cdot \mu(A)=\int \infty\cdot \mathbbm{1}_A\odif \mu\leqslant \int |u|\odif \mu<\infty\Longrightarrow \mu(A)=0$. Conversely, $|u|=|u\cdot \mathbbm{1}_{A^c}|+\infty\cdot \mathbbm{1}_A$. Thus, by Theorem \ref{properties of integral of positive measurable functions}, $\int |u|\odif \mu=\int |u\cdot \mathbbm{1}_{A^c}|\odif \mu+\infty\cdot \mu(A)=\int |u\cdot \mathbbm{1}_{A^c}|\odif \mu<\infty$, so $u$ is integrable. % Theorem \ref{integral of almost everywhere equal measurable functions}.
    \end{proof}
\subsection{Insensitivity}
    \begin{thm}\label{integral of almost everywhere equal measurable functions}
        Let $u,v:E\rightarrow \er$ be measurable functions. 
        
        If $u=v$ $\mu$-a.e., then $u$ is $\mu$-quasi-integrable if and only if $v$ is $\mu$-quasi-integrable, in which case
        \[
            \int_\mu u=\int_\mu v
        \]
    \end{thm}
    \begin{proof}
        If $u=v$ $\mu$-a.e., then $u^\pm=v^\pm$ $\mu$-a.e., and the claim follows from Theorem \ref{integrals of a.e. equal positive measurable functions}.
    \end{proof}
\subsection{Isotonicity}
    \begin{thm}\label{monotonicity of integral of quasi-integrable functions}
        Let $u,v:E\rightarrow \er$ be $\mu$-quasi-integrable.
        \[
            u\leqslant v\,\,\,\mu\text{-a.e.}\Longrightarrow \int_\mu u\leqslant \int_\mu v
        \]
    \end{thm}
    \begin{proof}
        If there is a $\mu$-full set $F$ such that $u(x)\leqslant v(x)$, which is, by Theorem \ref{positive parts preserve order relation while negative parts reverse}, equivalent to $u^+(x)\leqslant v^+(x)$ and $u^-(x)\geqslant v^-(x)$, for all $x\in F$, then by Theorem \ref{integrals of a.e. equal positive measurable functions} \& \ref{monotonicity of integral of positive measurable functions}, $\int_\mu u^+=\int_\mu u^+\cdot\mathbbm{1}_F\leqslant \int_\mu v^+\cdot\mathbbm{1}_F=\int_\mu v^+$ and similarly $\int_\mu u^-\geqslant \int_\mu v^-$.
        % Note that for $a,b,c,d\in\er$, if $a\leqslant b$ and $c\leqslant d$, then $a+c\leqslant b+d$ provided both $a+c$ and $b+d$ are defined. Indeed, if $a+c=-\infty$ or $b+d=\infty$ then the equality $a+c\leqslant b+d$ is trivially true. If $a+c=\infty$ then either $a=\infty$ or $c=\infty$, in which case, $b\geqslant a=\infty$ and $d\geqslant c=\infty$.
    \end{proof}
\subsection{Linearity}
    \begin{thm}\label{homogeneity}
        Let $u:E\rightarrow \er$ be $\mu$-quasi-integrable. Let $\alpha\in\mathbb{R}$.
        \[
            \int_\mu\alpha u=\alpha\int_\mu u
        \]
    \end{thm}
    \begin{thm}\label{additivity}
        Let $u,v:E\rightarrow \er$ be $\mu$-quasi-integrable. If $u+v:E\rightarrow \er$ is defined and
        \[
            \int_\mu u+\int_\mu v
        \]
        is defined, then $u+v$ is $\mu$-quasi-integrable and
        \[
            \int_\mu u+v=\int_\mu u+\int_\mu v
        \]
    \end{thm}
    \begin{proof}
        $({u}+{v})^+-({u}+{v})^-={u}+{v}=({u}^+-{u}^-)+({v}^+-{v}^-)=({u}^++{v}^+)-({u}^-+{v}^-)$. Thus, $({u}+{v})^++({u}^-+{v}^-)=({u}^++{v}^+)+({u}+{v})^-$. By repeated application of Theorem \ref{properties of integral of positive measurable functions}, 
        \begin{equation}\label{equation for linearity}
            \int ({u}+{v})^++\int {u}^-+\int {v}^-=\int{u}^++\int{v}^++\int({u}+{v})^-
        \end{equation}
        If $\int u+\int v$ is defined, then either $\int u^+,\int v^+<\infty$ or $\int u^-,\int v^-<\infty$. For reasons of symmetry, assume $\int u^+,\int v^+<\infty$. Then $\int (u+v)^+<\infty$ (so $u+v$ is quasi-integrable) since $(u+v)^+\leqslant u^++v^+$ and $(u+v)^-\leqslant u^-+v^-$. If either $\int u^-=\infty$ or $\int v^-=\infty$, then in \eqref{equation for linearity} necessrily $\int (u+v)^-=\infty$ and hence $\int u+v=-\infty=\int u +\int v$; otherwise, all terms in \eqref{equation for linearity} are finite.
    \end{proof}
\subsection{\texorpdfstring{$L^1$}{L1} as a Riesz space}
    \begin{define}
        A partial function $u$ from $E$ is $L^1_\mathbb{R}$ if the domain of $u$ contains a $\mu$-full subset $F$ on which $u$ is $\mathbb{R}$-valued and agrees with a measurable function $\tilde{u}:E\rightarrow\mathbb{R}$ that is $\mu$-integrable;
        % \[
        %     \tilde{u}:E\rightarrow\mathbb{R},\,\,x\mapsto \begin{cases}
        %         u(x)&\text{for }x\in F\\
        %         0&\text{for }x\notin F
        %     \end{cases}
        % \]
        furthermore, the $\mu$-integral of $u$ is defined\footnote{
            If $\tilde{u}_1,\tilde{u}_2:E\rightarrow \mathbb{R}$ are both measurable functions that agree with $u$ on $\mu$-full sets $F_1,F_2$ resp., then $F_1\cap F_2$ is $\mu$-full and $\tilde{u}_1(x)=u(x)=\tilde{u}_2(x)$ for all $x\in F_1\cap F_2$, that is, $\tilde{u}_1=\tilde{u}_2$ $\mu$-a.e. By Theorem \ref{integral of almost everywhere equal measurable functions}, $\tilde{u}_1$ is $\mu$-quasi-integrable if and only if $\tilde{u}_2$ is $\mu$-quasi-integrable, and in which case, $\int_\mu\tilde{u}_1=\int_\mu\tilde{u}_2$. Thus, the definition is unambiguous.
        } as that of $\tilde{u}$:
        \[
            \int_\mu u:=\int_{\mu}\tilde{u}
        \]
        
        The quotient, of the set of $L^1_\mathbb{R}$ functions on $(E,\mathscr{E},\mu)$, by almost everywhere equality, is denoted by $L^1_\mathbb{R}(E,\mathscr{E},\mu)$, or simply $L^1_\mathbb{R}$ when the underlying measure space is clear from context.
    \end{define}
    \begin{thm}\label{linearity of the integral on L^1}
        $L^1_\mathbb{R}$ is a Riesz space on which
        \[
            \int_\mu:L^1_\mathbb{R}\rightarrow \mathbb{R}
        \]
        is a positive linear form.
    \end{thm}
    \begin{proof}
        By definition, for $u,v\in L^1_\mathbb{R}$ there exist integrable measurable functions $\tilde{u}\in u,\tilde{v}\in v$. Then $\alpha u,u+v,u\vee v, u\wedge v$ are defined as the equivalence classes of $\alpha\tilde{u},\tilde{u}+\tilde{v}, \tilde{u}\vee\tilde{v},\tilde{u}\wedge \tilde{v}$ resp., whose integrability follows from $\int |\alpha u|=|\alpha|\int |u| <\infty$ and $\int |u+v|, \int |u\vee v|, \int |u\wedge v|\leqslant \int (|u|+|v|)=\int |u|+\int |v|<\infty$. It is clear that this definition of $\alpha u,u+v,u\vee v, u\wedge v$ does not depend on the choice of the integrable measurable functions $\tilde{u}\in u$, $\tilde{v}\in v$. Thus, $L^1_\mathbb{R}$ is a Riesz space. Theorem \ref{homogeneity} \& \ref{additivity}.
    \end{proof}
    \begin{thm}\label{triangle inequality 1}
        For all $u\in L^1$,
        \[
            \left|\int_\mu u\right|\leqslant \int_\mu |u|
        \]
    \end{thm}
    \begin{proof}
        $|\int  u|=\pm\int u=\int \pm u\leqslant \int |u|$ by Theorem \ref{monotonicity of integral of quasi-integrable functions}.
    \end{proof}
\subsection{Characterization of quasi-integrability via integrability}
    \begin{thm}\label{characetrization of quasi-integrable functions}
        % Let $u:X\rightarrow \er$ be a measurable function.
        % \begin{align*}
        %     \int u\odif \mu\in\,\,]\!-\!\infty,\infty]&\Longleftrightarrow\int u^-\odif \mu<\infty \\
        %     &\Longleftrightarrow u\text{ is a majorant of some $\mu$-integrable function }X\rightarrow \er
        % \end{align*}
        % \begin{align*}
        %     \int u\odif \mu\in[-\infty,\infty[&\Longleftrightarrow \int u^+\odif \mu<\infty \\
        %     &\Longleftrightarrow u\text{ is a minorant of some $\mu$-integrable function }X\rightarrow \er
        % \end{align*}
        A measurable function $u:X\rightarrow \er$ is $\mu$-quasi-integrable if and only if it is a majorant or a minorant of some $\mu$-integrable function.

        If there exists $\mu$-integrable functions $v,w$ such that $v\leqslant u\leqslant w$, then $u$ is $\mu$-integrable.
    \end{thm}
    % \begin{thm}
    %     Let 
    %     \begin{gather*}
    %         \widetilde{M}:=\{(u,A_1,A_2)\in M\times \Sigma X\times \Sigma X\,|\,A_1,A_2\subset \{u=0\},\, A_1\cap A_2=\varnothing\}\\
    %         \widetilde{\mathfrak{L}^1}:=\{(u,N_1,N_2)\in \mathfrak{L}^1\times \mathscr{N}_\mu\times \mathscr{N}_\mu\,|\,N_1,N_2\subset \{u=0\},\, N_1\cap N_2=\varnothing\}
    %     \end{gather*}
        
    %     The function
    %     \[
    %         \overline{M}\rightarrow \widetilde{M},\quad u\mapsto (u\cdot \mathbbm{1}_{\{u=\infty\}^c\cap \{u=-\infty\}^c},\{u=\infty\},\{u=-\infty\})
    %     \]
    %     is an isomorphism in \cat{Set}, whose inverse is given by
    %     \[
    %         \widetilde{M}\rightarrow \overline{M},\quad (u,A_1,A_2)\mapsto u+\infty\mathbbm{1}_{A_1}+(-\infty)\mathbbm{1}_{A_2}
    %     \]

    %     Moreover, under this isomorphism,
    %     \[
    %         \{\mu\text{-integrable functions in }\overline{M}\}\cong \widetilde{\mathfrak{L}^1}
    %     \]
    % \end{thm}
\subsection{Monotone convergence theorem}
    \begin{thm}
        If $(u_n)_{n\in\mathbb{N}}$ is an increasing (resp. decreasing) sequence of measurable functions $X\rightarrow \er$ minorized (resp. majorized) by a $\mu$-integrable measurable function, then $u_n$ is $\mu$-quasi-integrable for all $n\in\mathbb{N}$ and so is $\lim_{n\to\infty}u_n$, and
        \[
            \int_\mu \lim_{n\to\infty}u_n=\lim_{n\to\infty} \int_\mu u_n
        \]
    \end{thm}
    \begin{proof}
        Let $u=\lim_{n\to\infty}u_n$. Quasi-integrability of $u_n$ and $u$ follows directly from Theorem \ref{characetrization of quasi-integrable functions}. By Theorem \ref{positive parts preserve order relation while negative parts reverse}, $u_n^+\uparrow u^+$ and $u_n^-\downarrow u^-$. Since $(u_n)$ is minorized by some integrable function, $\int u_n^-<\infty$ and $\int u^-<\infty$. By Theorem \ref{integrable functions are a.e. real-valued}, $u_n^-,u^-$ is a.e. real-valued and $u_n^-,u^-\in L^1_\mathbb{R}$ for all $n$.\footnote{
            Here $L^1_\mathbb{R}$ should be understood as the set of $\mu$-almost measurable $\mu$-a.e.\! real-valued $\mu$-a.e.\! defined functions on $E$ that are $\mu$-integrable, without taking the quotient.
        } Moreover, $u_1^--u_n^-$ is a.e. defined, positive and increases to $u_1^--u^-$. By Beppo Levi's theorem, $\int u_n^+\uparrow \int u^+$ and $\int u_1^--\int u_n^-=\int u_1^--u_n^-\uparrow \int u_1^--u^-=\int u_1^--\int u^-$. Again by finiteness of $\int u_1^-$, I have $-\int u_n^-\uparrow -\int u^-$. Thus, $\int u_n\uparrow \int u$. If $(u_n)$ is decreasing and majorized by an integrable function $w$, then $(-u_n)$ is increasing and minorized by $-w$.
    \end{proof}
    % As a corollary to the preceding theorem,
    % \begin{thm}
    %     For $(u_n)_{n\in\mathbb{N}}$ a monotone sequence of $\mu$-integrable measurable functions $X\rightarrow\er$,
    %     \[
    %         \int_\mu \lim_{n\to\infty}u_n=\lim_{n\to\infty} \int_\mu u_n
    %     \]
    % \end{thm}
    % \begin{proof}
    %     % Alternatively, $u=\lim_{n\to\infty}u_n$ is a function $X\rightarrow \er$, which is measurable and quasi-integrable (Theorem \ref{characetrization of quasi-integrable functions}). Assume $u_n$ is increasing, so $v_n=u_n-u_1\geqslant 0$ and $\int u-u_1\odif \mu=\int \lim_{n\to\infty}v_n\odif \mu=\lim_{n\to\infty}\int v_n\odif \mu=\lim_{n\to\infty}\int u_n\odif \mu-\int u_1\odif \mu$. It then suffices to show that $\int u-u_1\odif \mu+\int u_1\odif \mu=\int u\odif \mu$. If $u$ is integrable, then the equality follows from linearity of the integral (by identifying $u$ with $u\mathbbm{1}_{\{|u|=\infty\}^c}$). If $\int u\odif \mu=\infty$, then $\int u-u_1\odif \mu=\infty$ since otherwise $N=\{|u-u_1|=\infty\}\in\mathscr{N}_\mu$ and $\int u\odif \mu=\int u\mathbbm{1}_{N^c}\odif\mu=\int (u-u_1)\mathbbm{1}_{N^c}\odif\mu+\int u_1\mathbbm{1}_{N^c}\odif\mu<\infty$. 
    % \end{proof}
\subsection{A partial converse to isotonicity and insensitivity}
    \begin{thm}\label{partial converse to isotonicity and insensitivity for integrable functions}
        Let $u,v:E\rightarrow \er$ be $\mu$-integrable. If for all $A\in\mathscr{E}$,
        \[
            \int_\mu u\cdot\mathbbm{1}_A\leqslant \int_\mu v\cdot\mathbbm{1}_A
        \]
        then $u\leqslant v$ $\mu$-a.e. Thus, if for all $A\in\mathscr{E}$,
        \[
            \int_\mu u\cdot\mathbbm{1}_A= \int_\mu v\cdot\mathbbm{1}_A
        \]
        then $u=v$ $\mu$-a.e.
    \end{thm}
    \begin{proof}
        Let $A=\{u>v\}$ and $A_n=\{u>v+1/n\}$. Then $\int u\cdot\mathbbm{1}_{A_n}\geqslant \int (v+1/n)\cdot\mathbbm{1}_{A_n}=\int v\cdot\mathbbm{1}_{A_n}+\mu(A_n)/n\geqslant \int u\cdot\mathbbm{1}_{A_n}+\mu(A_n)/n$. Thus, $\mu(A_n)=0$. But $A_n\uparrow A$.
    \end{proof}
    \begin{thm}
        Let $\mu$ be $\sigma$-finite. Let $u,v:E\rightarrow \er$ be $\mu$-quasi-integrable. If for all $A\in\mathscr{E}$,
        \[
            \int_\mu u\cdot\mathbbm{1}_A\leqslant \int_\mu v\cdot\mathbbm{1}_A
        \]
        then $u\leqslant v$ $\mu$-a.e. Thus, if for all $A\in\mathscr{E}$,
        \[
            \int_\mu u\cdot\mathbbm{1}_A= \int_\mu v\cdot\mathbbm{1}_A
        \]
        then $u=v$ $\mu$-a.e.
    \end{thm}
    \begin{proof}
        Assume $\int u^-<\infty$. Then $\int v^-<\infty$. Let $B_n\uparrow X$ with $\mu(B_n)<\infty$. Set $A_n:=B_n\cap \{v\leqslant n\}\in\mathscr{E}$. Then $A_n\uparrow \{v<\infty\}$ with $\mu(A_n)<\infty$. Thus, $\int v^+\cdot\mathbbm{1}_{A_n}\leqslant n\cdot\mu(A_n)<\infty$ and $v\cdot \mathbbm{1}_{A_n}$ is integrable. Then $u\cdot \mathbbm{1}_{A_n}$ is also integrable. By Theorem \ref{partial converse to isotonicity and insensitivity for integrable functions}, $u\cdot \mathbbm{1}_{A_n}\leqslant v\cdot\mathbbm{1}_{A_n}$ a.e. Then $u\cdot \mathbbm{1}_{\{v<\infty\}}\leqslant v\cdot\mathbbm{1}_{\{v<\infty\}}$ a.e. But $v|_{\{v<\infty\}^c}=\infty$. Thus, $u\leqslant v$ a.e.

        The case when $\int v^+<\infty$ is similar. If $\int u^-=\infty$ and $\int v^+=\infty$, then both $u\cdot\mathbbm{1}_{\{u>v\}}$ and $v\cdot\mathbbm{1}_{\{u>v\}}$ are integrable. By Theorem \ref{partial converse to isotonicity and insensitivity for integrable functions}, $u\cdot\mathbbm{1}_{\{u>v\}}\leqslant v\cdot\mathbbm{1}_{\{u>v\}}$ a.e. Thus, $\{u>v\}$ is null.
    \end{proof}
\subsection{The sine integral}
    \[\lim_{T\to\infty}\int_0^T \frac{\sin\xi}{\xi}\odif \xi=\frac{\pi}{2}\]
    \begin{proof}
        
    \end{proof}
\subsection{The general transformation theorem for quasi-integrable functions}
    \begin{thm}\label{the general transformation theorem for quasi-integrable functions}
        Let $\begin{tikzcd}
            (E,\mathscr{E}) \ar[r, "f"] & (F,\mathscr{F}) \ar[r, "u"] & \er
        \end{tikzcd}$ be a diagram in \cat{Mes}. $u$ is $f_*\mu$-quasi-integrable if and only if $f^*u$ is $\mu$-quasi-integrable, in which case,
        \[\int_\mu f^* u=\int_{f_*\mu} u\]
        where $f^* u=u\circ f$ is the pullback of $u$ by $f$.
    \end{thm}
    \begin{proof}
        Apply Theorem \ref{the general transformation theorem}: $\int_\mu(f^*u)^\pm=\int_\mu f^*u^\pm=\int_{f_*\mu}u^\pm$.
    \end{proof}
\subsection{**Continuous functions with compact support}
\subsection{*Integration of complex-valued functions}
    %\[\mathcal{L}^p_{\mathbb{C}}(\mu):=\{f\in\mathcal{M}(\mathscr{A};\mathbb{C})\,|\,|f|\in\mathcal{L}^p_{\mathbb{R}}(\mu)\}\]
    \begin{lem}
        The following are equivalent:
        \begin{tight_enumerate}
            \item $f\in \mathcal{L}^1_{\mathbb{C}}(\mu)$.
            \item $|f|\in \mathcal{L}^1_{\mathbb{R}}(\mu)$.
            \item $\RE{f}\in\mathcal{L}^1_{\mathbb{R}}(\mu)\wedge\IM{f}\in\mathcal{L}^1_{\mathbb{R}}(\mu)$.
        \end{tight_enumerate}
        in which case
        \[
        \int f\odif \mu:=\int \RE{f}\odif \mu+i\int\IM{f}\odif \mu
    \]
    \end{lem}
    \begin{lem}
        Let $f,g\in \mathcal{L}^1_{\mathbb{C}}(\mu)$. Let $\alpha,\beta\in\mathbb{C}$.
        \begin{tight_enumerate}
            \item $\displaystyle \RE\int f\odif\mu=\int \RE f\odif\mu,\quad \IM\int f\odif\mu=\int \IM f\odif\mu,\quad \overline{\int f\odif \mu}=\int \overline{f}\odif \mu$.
            \item $\alpha f+\beta g\in \mathcal{L}^1_{\mathbb{C}}(\mu)$ and $\displaystyle \int \alpha f+\beta g\odif \mu=\alpha\int f\odif \mu+\beta\int g\odif \mu$.
            \item $\displaystyle \left|\int f\odif \mu\right|\leqslant\int |f|\odif \mu$.
        \end{tight_enumerate}
    \end{lem}
    \begin{proof}
        $|\int f\odif \mu|=\lambda\cdot \int f\odif\mu$ for some $\lambda\in\mathbb{C}$ with $|\lambda|=1$. Thus, $|\int f\odif \mu|=\int \lambda f\odif\mu=\RE \int \lambda f\odif\mu=\int \RE \lambda f\odif\mu\leqslant \int |\lambda f|\odif \mu=\int |f|\odif \mu$.
    \end{proof}
    \begin{thm}\label{linearity of the integral on L^1_C}
        $\mathfrak{L}_\mathbb{C}^1$ is a vector space on which
        \[
            \int \ph\odif \mu:\mathfrak{L}^1\rightarrow \mathbb{C}
        \]
        is a (continuous w.r.t. seminorm $\|\|_1$) and positive linear form.
    \end{thm}
\subsection{Integration on a measurable subset}
    \begin{thm}
        Let $A\in\Sigma X$ and $u:A\rightarrow \er$. Equip $A$ with the trace $\sigma$-algebra $\Sigma X|_A$.

        Let $\tilde{u}:X\rightarrow \er,\,\,x\mapsto (x\in A\,?\,u(x):0)$.
        \[
            \tilde{u}\in\cat{Mes}(X,\er)\Longleftrightarrow u\in\cat{Mes}(A,\er)
        \]
        in which case,
        \[
            \int \tilde{u}\odif\mu=\int u\odif\mu|_{\Sigma X|_A}
        \]
    \end{thm}
    \begin{proof}
        For $2$, first $\Sigma X|_A\subset \Sigma X$ so $\mu|_{\Sigma X|_A}$ is a well-defined measure on $\Sigma X|_A$. Suppose $u\geqslant 0$ and thus $\tilde{u}\geqslant 0$. Let $(f_n)$ be a sequence of positive simple functions on $X$ such that $f_n\uparrow \tilde{u}$. Then $f_n|_A$ is a sequence of positive simple functions on $A$ with $f_n|_A\uparrow u$, since $\widetilde{f_n|_A}=f_n$. Thus, $\int \tilde{u}\odif\mu=\lim I_\mu(f_n)=\lim I_{\mu|_{\Sigma X|_A}}(f_n|_A)=\int u\odif\mu_{\Sigma X|_A}$ since $(f_n|_A)^{-1}\{y\}=(f_n)^{-1}\{y\}$ for $y>0$, that is, $I_\mu(f_n)=I_{\mu|_{\Sigma X|_A}}(f_n|_A)$. It remains to point out $\widetilde{u^\pm}=\tilde{u}^\pm$.
    \end{proof}
\subsection{Integration of vector-valued functions}
    Let $\mathbb{K}\in\{\mathbb{R},\mathbb{C}\}$ and $n\in\mathbb{N}$. Let $\mathbb{K}^n$ be equipped with the Borel $\sigma$-algebra $\mathscr{B}(\mathbb{K}^n)$.
    \begin{define}
        A function $u:E\rightarrow\mathbb{K}^n$ is \textit{$\mu$-integrable} if it is measurable and $u_i=\pi_i\circ u:E\rightarrow\mathbb{K}$ is $\mu$-integrable for all $i=1,\dots,n$, in which case, the \textit{$\mu$-integral} of $u$ is
        \[
            \int_\mu u:=\left(\int_\mu u_1,\dots,\int_\mu u_n\right)\in\mathbb{K}^n
        \]
    \end{define}
    \begin{thm}\label{commutation of integral and linear map}
        Let $u:E\rightarrow\mathbb{K}^n$ be $\mu$-integrable and $A:\mathbb{K}^n\rightarrow\mathbb{K}^m$ a $\mathbb{K}$-linear map. Then $A\circ u$ is $\mu$-integrable and
        \[
            \int_\mu A\circ u=A\int_\mu u
        \]
    \end{thm}
    \begin{proof}
        By definition, it suffices to show for $m=1$. In this case, $A$ is a $\mathbb{K}$-linear form, which can be written as $x\mapsto a_1x_1+\cdots+a_nx_n$, where $a_1,\dots,a_n\in\mathbb{K}$. Then $A\circ u=a_1 u_1+\cdots+a_nu_n\in L_\mathbb{K}^1$ and $\int A\circ u=a_1\int u_1+\cdots+a_n\int u_n=A\int u$.
    \end{proof}
    % \begin{thm}
    %     Let $V$ be an $n$-dimensional vector space over $\mathbb{K}$. Let $\alpha,\beta:\mathbb{K}^n\rightarrow V$ be $\mathbb{K}$-\cat{Vect} isomorphisms induced by two bases in $V$. Let $u:E\rightarrow V$ a function. Then $\alpha^{-1}\circ u$ is $\mu$-integrable if and only if $\beta^{-1}\circ u$ is $\mu$-integrable, in which case,
    %     \[
    %         \alpha\left(\int_\mu \alpha^{-1}\circ u\right)=\beta\left(\int_\mu \beta^{-1}\circ u\right)
    %     \]
    % \end{thm}
    % \begin{proof}
    %     By Theorem \ref{commutation of integral and linear map}, if $\alpha^{-1}\circ u$ is integrable, then $\beta^{-1}\circ u=(\beta^{-1}\circ\alpha)\circ(\alpha^{-1}\circ u)$ is integrable, and conversely by symmetry, and in which case, $\int \beta^{-1}\circ u=(\beta^{-1}\circ\alpha)\int \alpha^{-1}\circ u$, i.e., $\beta(\int \beta^{-1}\circ u)=\alpha(\int \alpha^{-1}\circ u)$.
    % \end{proof}
    % Thus, to integrate a function valued in a finite-dimensional vector space $V$, one can choose an arbitrary basis in $V$ and integrate coordinate-wise. The resulting integral (as a vector in $V$) does not depend on the choice of the basis.
\subsection{Minkowski's Inequality for Integrals}
    \begin{thm}
        Be $(X,\mathscr{A},\mu)$ and $(Y,\mathscr{B},\nu)$ $\sigma$-finite measure spaces. Let $u\in\mathcal{M}(\mathscr{A}\otimes\mathscr{B};\overline{\mathbb{R}})$.
        \[\left(\int\left(\int|u(x,y)|\nu(\odif y)\right)^p\mu(\odif x)\right)^{1/p}\leqslant \int\left(\int |u(x,y)|\mu(\odif x)\right)^{1/p}\nu(\odif y)\]
        holds for all $p\in [1,\infty)$, with equality for $p=1$.
    \end{thm}
\subsection{*Relation between Beta and Gamma Functions}
\subsection{*A More General Version of Green's Theorem}
\subsection{*Cauchy's Integral Theorem}
\subsection{*Pompeiu's Formula}
\section{Convergence theorems and integral calculus}
\subsection{Dominated convergence}
    \begin{thm}\label{dominated convergence theorem}
        Let $(u_n)\in\cat{Mes}(X,\mathbb{K})^\mathbb{N}$. If
        \begin{tight_enumerate}
            \item there exists $w\in\mathfrak{L}_\mathbb{R}^1(\mu)$ such that $|u_n|\leqslant w\,\,\mu$-a.e. for all $n\in\mathbb{N}$,
            \item there exists $u\in\cat{Mes}(X,\mathbb{K})$ such that $u_n\to u$ $\mu$-a.e.,
        \end{tight_enumerate}
        then $u_n\in\mathfrak{L}^1_\mathbb{K}(\mu)$ for all $n\in\mathbb{N}$, $u\in\mathfrak{L}^1_\mathbb{K}(\mu)$, and the following holds:
        \begin{gather*}
            \lim_{n\to\infty}\int |u_n-u|\odif\mu=0 \\
            \lim_{n\to\infty}\int u_n\odif\mu=\int u\odif\mu
        \end{gather*}
    \end{thm}
    % Remark. The theorem still holds for $\mathbb{K}=\er$ by using the usual identification of integrable $\er$-valued functions with integrable $\mathbb{R}$-valued functions, since $N=\{|w|=\infty\}\cup \{|u|=\infty\} \cup \bigcup{n=1}^\infty\{|u_n|=\infty\}$ has to be $\mu$-null (as a consequence of the existence of a $\mu$-integrable positive function $w\in\mathfrak{L}^1_{\er})$.
    \begin{proof}
        Since a countable union of $\mu$-null sets is $\mu$-null, there exists $N\in\mathscr{N}_\mu$ such that $|u_n\one_{N^c}|\leqslant w \one_{N^c}\leqslant w$ for all $n\in\mathbb{N}$ and $(u_n\one_{N^c})\to u\one_{N^c}$. Thus, $u_n\one_{N^c}$ and hence $u_n$ are integrable. Since $|\sph|:\mathbb{K}\rightarrow\mathbb{R}$ is continuous, $|u\one_{N^c}|=\lim_{n\to\infty}|u_n\one_{N^c}|\leqslant w$. Thus, $u\one_{N^c}$ and hence $u$ are integrable. Since $|u_n-u|\one_{N^c}\leqslant 2w$ for all $n\in\mathbb{N}$, applying Fatou's theorem to the sequence $(2w-|u_n-u|\one_{N^c})\in\mathfrak{L}^1_{\mathbb{R}_+}(\mu)^\mathbb{N}$ yields $\int 2w\odif \mu\leqslant \varliminf_{n\to\infty} \int (2w-|u_n-u|\one_{N^c})\odif \mu=\int 2w\odif \mu- \varlimsup_{n\to\infty} \int |u_n-u|\one_{N^c}\odif \mu$. So $0\leqslant \varliminf_{n\to\infty} \int |u_n-u|\one_{N^c}\odif \mu\leqslant \varlimsup_{n\to\infty} \int |u_n-u|\one_{N^c}\odif \mu=0$. Theorem \ref{triangle inequality 1}.
    \end{proof}
    % \begin{remark}
    %     If $(u_n(x))$ converges in $\mathbb{K}$ for $\mu$-almost all $x\in X$, then there exists $N\in\mathscr{N}_\mu$ such that $u=\lim_{n\to\infty}u_n\one_{N^c}$ is a function $X\rightarrow \mathbb{K}$. Since $\cat{Mes}(X,\mathbb{K})$ is sequentially closed, $u$ is a measurable function such that $u_n\to u$ $\mu$-a.e..
    % \end{remark}
    Alternative presentation: \textbf{Actually not good since that the limit is taken in the topology of pointwise convergence is not obvious:}
    \begin{thm}
        Let $(u_n)\in\cat{Mes}(X,\mathbb{K})^\mathbb{N}$. If
        \begin{tight_enumerate}
            \item there exists a $\mu$-integrable $w\in L^1_\mathbb{R}(\mu)$ such that $|u_n|\leqslant w\,\,\mu$-a.e. for all $n\in\mathbb{N}$,
            \item $(u_n)$ converges $\mu$-a.e.,
        \end{tight_enumerate}
        then $u_n\in L^1_\mathbb{K}(\mu)$ for all $n\in\mathbb{N}$, $\lim_{n\to\infty} u_n\in {L}^1_\mathbb{K}(\mu)$, and the following holds:
        \begin{gather*}
            \lim_{n\to\infty}\int \left|u_n-\lim_{n\to\infty} u_n\right|\odif\mu=0 \\
            \lim_{n\to\infty}\int u_n\odif\mu=\int \lim_{n\to\infty} u_n\odif\mu
        \end{gather*}
    \end{thm}
    \begin{crl}
        Let $u:X\rightarrow \mathbb{K}$ be $\mu$-integrable over $A\in\Sigma X$. If
        \[
            A=\bigcup_{n=1}^\infty A_n
        \]
        for $(A_n)\in\Sigma X^\mathbb{N}$ with $A_i\cap A_j\in\mathscr{N}_\mu$ for all $i,j\in\mathbb{N}$ distinct, then
        \[
            \int _Au\odif \mu=\sum_{n=1}^\infty\int_{A_n}u\odif\mu
        \]
    \end{crl}
    \begin{proof}
        
    \end{proof}
\subsection{Parameter-dependent integrals}
    \begin{thm}\label{continuity of parameter-dependent integral}
        Let $T$ be a first-countable space and $f:T\times X\rightarrow \mathbb{K}$ a function. Let $t_0\in T$. If
        \begin{tight_enumerate}
            \item $f(t,\ph)\in L^1_{\mathbb{K}}$ for all $t\in T$.
            \item $f(\ph,x):T\rightarrow \mathbb{K}$ is continuous at $t_0$ for $\mu$-almost all $x\in X$.
            \item there exists a neighborhood $U$ of $t_0$ and a $\mu$-integrable function $g\in L^1_\mathbb{R}$ such that $|f(t,x)|\leqslant g(x)$ for all $t\in U$ and $\mu$-almost all $x\in X$.
        \end{tight_enumerate}
        then the function
        \[
            F:T\rightarrow \mathbb{K},\,\,t\mapsto \int f(t,\bullet )\odif \mu
        \]
        is continuous at $t_0$, and also the function
        \[
            \Phi:T\rightarrow \mathfrak{L}^1_{\mathbb{K}},\,\,t\mapsto f(t,\bullet)
        \]
        is continuous at $t_0$.
    \end{thm}
    \begin{proof}
        Let $(t_n)$ be a sequence in $U$ tending to $t_0$, so $|f(t_n,\bullet)|\leqslant g$ $\mu$-a.e.. Since $f(\bullet,x)$ is continuous at $t_0$ for almost all $x\in X$, $f(t_n,\bullet)\to f(t_0,\bullet)$ $\mu$-a.e.. By dominated convergence, $F(t_n)\to F(t_0)$. By first-countability (hereditary) of $U$, $F(t)$ is continuous at $t_0$ in $U$, and hence in $T$.
    \end{proof}
    \begin{thm}\label{differentiability of parameter-dependent integral}
        Let $]a,b[\,\,\subset\mathbb{R}$ and $f:\,\,]a,b[\,\,\times X\rightarrow\mathbb{R}$ a function. Let $t_0\in\,\,]a,b[$. If
        \begin{tight_enumerate}
            \item $f(t,\bullet)\in L^1_\mathbb{R}$ for all $t\in \,\,]a,b[$,
            \item the partial derivative $\displaystyle \pdv{f}{t}$ exists at $(t_0,x)$ for all $x\in X$,
            \item there exist a neighborhood $U$ of $t_0$ and $g\in L^1_\mathbb{R}$ such that for all $t\in U\setminus\{t_0\}$ and $\mu$-a.a. $x\in X$
            \[
                \left|\frac{f(t,x)-f(t_0,x)}{t-t_0}\right|\leqslant g(x)
            \]
        \end{tight_enumerate}
        then the function
        \[F:\,\,]a,b[\,\,\rightarrow \mathbb{R},\,\,t\mapsto \int f(t,\bullet)\odif\mu\]
        is differentiable at $t_0$, the function $\displaystyle \pdv{f}{t}(t_0,\bullet)\in L^1_\mathbb{R}$ and
        \[
            F'(t_0)=\odv{}{t}_{t_0}\int f(t,\bullet)\odif\mu=\int \pdv{f}{t}(t_0,\bullet)\odif\mu
        \]
    \end{thm}
    \begin{proof}
        Let $(t_n)$ be a sequence in $U\setminus \{t_0\}$ tending to $t_0$, so $u_n=(f(t_n,\bullet)-f(t_0,\bullet))/(t_n-t_0)\leqslant g$ $\mu$-a.e.. Since the partial derivative exists at $(t_0,x)$ for all $x$, $u_n\to\pdv{f}{t}(t_0,\bullet)$ (pointwise).
    \end{proof}
\subsection{**Holomorphic Dependence of the Integral on a Complex Parameter}
\subsection{**Scheff\'{e}'s Lemma}
    $||a|-|b||\leqslant |a-b|$.
\subsection{The Lebesgue measure}
    The following subsections concern functions $f:X\rightarrow \mathbb{K}$ with $X\in \mathscr{L}^p$, and $X$ is equipped with the trace $\sigma$-algebra $\mathscr{L}^p|_X$. As usual, the codomain is equipped with the Borel $\sigma$-algebra $\mathscr{B}\mathbb{K}$.
    \begin{thm}
        For $X\in \mathscr{L}^p$,
        \[
            \mathscr{L}^p|_X=\{A\in \mathscr{L}^p\,|\,A\subset X\}
        \]
    \end{thm}
    \begin{thm}
        $\cat{Mes}(X,\mathbb{K})$ is a $\mathbb{K}$-algebra. $f\in \cat{Mes}(X,\mathbb{K})\Longrightarrow |f|\in \cat{Mes}(X,\mathbb{K})$. $\cat{Mes}(X,\mathbb{R})$ is a lattice. 
    \end{thm}
    \begin{thm}
        If $f:X\rightarrow \mathbb{K}$ is continuous, then $f\in\cat{Mes}((X,\mathscr{B}^p|_X),\mathbb{K})\subset\cat{Mes}((X,\mathscr{L}^p|_X),\mathbb{K})$.
    \end{thm}
    \begin{proof}
        A continuous map is $\mathscr{B}A/\mathscr{B}\mathbb{K}$-measurable, i.e., $\mathscr{B}^p|_A/\mathscr{B}\mathbb{K}$-measurable, but $\mathscr{B}^p|_A\subset \mathscr{L}^p|_A$.
    \end{proof}
    \begin{thm}
        If $f:I\rightarrow \mathbb{R}$ is monotone, where $I\subset \mathbb{R}$ is an interval, then $f\in\cat{Mes}(I,\mathbb{K})$.
    \end{thm}
    \begin{proof}
        $f^{-1}([y,\infty[)$ is a subinterval of $I$. Recall Theorem \ref{measurable functions}.
    \end{proof}
    \begin{define}
        Let $X\in\mathscr{L}^p$. A function $f:X\rightarrow\mathbb{K}$ is \textit{integrable} if it is $\lambda^p|_{\mathscr{L}^p|_X}$-integrable.
    \end{define}
\subsection{The fundamental theorem of calculus}
    \begin{thm}\label{antiderivative of a continuous function}
        If $f:\mathbb{R}\rightarrow \mathbb{K}$ is continuous, then for any fixed $a\in\mathbb{R}$ the function
        \[
            F_a:\mathbb{R}\rightarrow \mathbb{K},\,\,t\mapsto \int_a^t f:=\begin{cases}
                \displaystyle\int_{\lambda}f\cdot\mathbbm{1}_{[a,t[} & \text{for }t\geqslant a\\
                -\displaystyle\int_{\lambda}f\cdot\mathbbm{1}_{[t,a[} & \text{for }t< a
            \end{cases}
        \]
        is differentiable with $F_a'=f$.
    \end{thm}
    \begin{proof}
        Since $f$ is continuous, it is integrable over every compact
        interval, so $F_a$ is well-defined.
        Fix $t\in\mathbb{R}$. For every $h\neq0$, additivity of
        the integral gives
        \[
            \frac{F_a(t+h)-F_a(t)}{h}-f(t)
            =\frac1h\int_t^{t+h}\bigl(f(x)-f(t)\bigr)\,dx.
        \]
        Consequently,
        \[
            \left|
                \frac{F_a(t+h)-F_a(t)}{h}-f(t)
            \right|
            \leqslant
            \sup_{|x-t|\leqslant |h|}|f(x)-f(t)|
            \to 0
            \qquad \text{as }h\to0
        \]
        by continuity of $f$ at $t$ (Cauchy criterion).
    \end{proof}
    \begin{thm}
        Let $-\infty<a<b<\infty$ and $f:\,\,]a,b[\,\,\rightarrow\mathbb{K}$ a differentiable function.
        
        If $f'$ is bounded, then $f'\in L^1(]a,b[)$, the limits
        \[
            f(a):=\lim_{x\downarrow a}f(x),\qquad f(b):=\lim_{x\uparrow b}f(x)
        \]
        exist in $\mathbb{K}$, and
        \[
            \int f'=f(b)-f(a)
        \]
    \end{thm}
    \begin{proof}
        Choose $M<\infty$ such that $|f'(x)|\leqslant M$ for
        every $x\in \,\,]a,b[$. Then by finite increment lemma, $|f(y)-f(x)|\leqslant M|y-x|$ for all $x,y\in\,\,]a,b[$. By completeness of $\mathbb{K}$ and Cauchy's criterion, the endpoint limits $f(a),f(b)$ exist in $\mathbb{K}$. Define extensions to $\mathbb{R}$ by
        \[
            \overline{f}(x):=
            \begin{cases}
                f(a),&x\leqslant a,\\
                f(x),&a< x< b,\\
                f(b),&x\geqslant b,
            \end{cases}
            \qquad
            \overline{f'}(x):=
            \begin{cases}
                f'(x),&a<x<b\\
                0,&x\notin \,\,]a,b[
            \end{cases}
        \]
        Then $\overline f$ is $M$-Lipschitz on $\mathbb{R}$ and $\overline f'(x)=\overline{f'}(x)$ for $x\in \mathbb{R}\setminus\{a,b\}$.
    
        For $n\geqslant1$, define the continuous function
        \[
            g_n(x):=
            n\left(\overline f\left(x+\frac1n\right)
            -\overline f(x)\right),
            \qquad x\in\mathbb{R}.
        \]
        Then, for $x\in\mathbb{R}\setminus\{a,b\}$,
        \[
            g_n(x)\to\overline{f'}(x)\qquad \text{as }n\to\infty
        \]
        Thus $\overline{f'}$ is measurable. Moreover, the Lipschitz
        bound gives $|g_n(x)|\leqslant M$ for every $x\in\mathbb{R}$.
        Since $]a,b[$ has finite measure, by dominated convergence theorem, $g_n\mathbbm{1}_{]a,b[},\overline{f'}\mathbbm{1}_{]a,b[}\in L^1$ and
        \[
            \int f'
            =\int_a^b\overline{f'}
            =\lim_{n\to\infty}\int_a^b g_n
        \]
    
        By translation invariance of Lebesgue measure, and Theorem \ref{antiderivative of a continuous function},
        \begin{align*}
            \int_a^b g_n&=n\left(
                \int_{a+1/n}^{b+1/n}\overline f
                -\int_a^b\overline f
            \right)=n\left(\int_b^{b+1/n}\overline f-\int_a^{a+1/n}\overline f\right)\overset{n\to\infty}{\longrightarrow} f(b)-f(a)
        \end{align*}
    \end{proof}
\subsection{Improper integrals}
    \begin{thm}
        Let $I\subset\mathbb{R}$ be an open or half-open interval. If $f:I\rightarrow \mathbb{K}$ is $\lambda$-integrable over every compact subinterval of $I$, then $f$ is $\lambda$-integrable if and only if for some $[a_n,b_n]\uparrow I$
        \[
            \lim_{n\to\infty}\int_{a_n}^{b_n} |f|\odif\lambda <\infty
        \]
        in which case,
        \[
            \int f\odif\lambda=\lim_{n\to\infty}\int_{a_n}^{b_n} f\odif\lambda
        \]
    \end{thm}
    \begin{proof}
        Since $|f|\cdot \one_{[a_n,b_n]}\uparrow |f|$, by Theorem \ref{Beppo Levi's Theorem}, $\int |f|\odif\lambda=\lim_{n\to\infty}\int |f|\cdot\one_{[a_n,b_n]}\odif\lambda$. Since $|f\cdot \one_{[a_n,b_n]}|\leqslant |f|$ and $f\cdot \one_{[a_n,b_n]}\to f$, if $f$ is $\lambda$-integrable, then by Theorem \ref{dominated convergence theorem}, $\int f\odif\lambda=\lim_{n\to\infty}\int f\cdot \one_{[a_n,b_n]}\odif\lambda$. % Alternatively, $f^+\cdot\one_{[a_n,b_n]}\uparrow f^+$ and $f^-\cdot\one_{[a_n,b_n]}\uparrow f^-$. Thus, $\int f^+\cdot\one_{[a_n,b_n]}\odif\lambda \uparrow \int f^+\odif\lambda$ and $\int f^-\cdot\one_{[a_n,b_n]}\odif\lambda \uparrow \int f^-\odif\lambda$ (note that for an arbitrary sequence $(u_n)$, generally $(u_n^+)$ and $(u_n^-)$ are not increasing), that is, $\int f\odif\lambda=\lim_{n\to\infty}\int f^+\cdot\one_{[a_n,b_n]}\odif\lambda-\lim_{n\to\infty}\int f^-\cdot\one_{[a_n,b_n]}\odif\lambda$.
    \end{proof}
\subsection{*Mean Value Theorems}
    Let $f,g:[a,b]\rightarrow \mathbb{R}$, where $[a,b]\subset\mathbb{R}$.
    \begin{thm}
        If $f$ is $\lambda$-integrable, $f\geqslant 0$ and $g$ is continuous, then $f\cdot g$ is $\lambda$-integrable and there exists $\xi\in [a,b]$ such that
        \[
            \int _a^b(f\cdot g)\odif\lambda = g(\xi)\int _a^bf\odif\lambda
        \]
    \end{thm}
    \begin{proof}
        $g([a,b])$ is connected and compact, thus a closed interval $[c,d]$. In particular, $|f\cdot g|\leqslant d\cdot |f|$ is integrable. Since $f\geqslant 0$, if $\int _a^bf\odif\lambda=0$, then $f$ and $f\cdot g$ vanish $\lambda$-a.e.; otherwise, $\int _a^b(f\cdot g)\odif\lambda/\int _a^bf\odif\lambda \in [c, d]=g([a,b])$.
    \end{proof}
    \begin{thm}\label{Bonnet's second mean value theorem}
        If $f$ is $\lambda$-integrable and $g$ is monotone, then $f\cdot g$ is $\lambda$-integrable and there exists $\xi\in[a,b]$ such that
        \[
            \int _a^b(f\cdot g)\odif\lambda = g(a)\int _a^\xi f\odif\lambda+g(b)\int_\xi^b f\odif\lambda
        \]
    \end{thm}
    \begin{proof}
        
    \end{proof}
    \begin{crl}
        If $f$ is $\lambda$-integrable, $g$ is decreasing and $g\geqslant 0$, then $f\cdot g$ is $\lambda$-integrable and there exists $\xi\in[a,b]$ such that
        \[
            \int _a^b(f\cdot g)\odif\lambda = g(a)\int _a^\xi f\odif\lambda
        \]
    \end{crl}
    \begin{proof}
        Apply Theorem \ref{Bonnet's second mean value theorem} to $\tilde{g}=g\cdot\one_{[a,b[}$.
    \end{proof}
\subsection{*Eulerian Integrals}
    Let $f_\alpha:(0,\infty)\rightarrow \mathbb{R},\,\,x\mapsto x^\alpha,\,\,\alpha\in\mathbb{R}$.
    \begin{enumerate}
        \item $\exp\circ\alpha\ln|_{(0,1)}\in\mathcal{L}^1_\mathbb{R}((0,1))\Longleftrightarrow \alpha>-1$.
        \item $\exp\circ\alpha\ln|_{[1,\infty)}\in\mathcal{L}^1_\mathbb{R}([1,\infty))\Longleftrightarrow \alpha<-1$.
    \end{enumerate}
    \par Let $f:(0,\infty)\rightarrow\mathbb{R},\,x\mapsto x^{\alpha}e^{-\beta x}$. Then $f\in\mathcal{L}^1((0,\infty);\mathbb{R})$ if $\alpha>-1$ and $\beta>0$.
    \begin{proof}
        Using the exponential series we find for all $N\in\mathbb{N}$ and $x>0$
        \[f(x)=x^{\alpha}e^{-\beta x}\leqslant x^\alpha\mathbbm{1}_{(0,1)}(x)+\frac{N!}{\beta^N}x^{\alpha-N}\mathbbm{1}_{[1,\infty)}(x)\]
        $\in\mathcal{L}^1((0,\infty);\mathbb{R})$ for large enough $N$.
    \end{proof}
    \[B(\alpha,\beta):=\int_{(0,1)}x^{\alpha -1}(1-x)^{\beta -1}\lambda^1(\odif x),\quad\Gamma(t):=\int_{(0,\infty)} x^{t-1}e^{-x}\lambda^1(\odif x)\]
    \begin{enumerate}
        \item $\Gamma$ is well-defined for all $t>0$.
        \item $\Gamma$ is continuous.
        \item $\Gamma$ is infinitely differentiable.
        \item $t\Gamma(t)=\Gamma(t+1)$, in particular, $\Gamma(n+1)=n!$.
    \end{enumerate}
    \begin{thm}
        The Beta function:
        \begin{enumerate}
        \item $B(\alpha,\beta)$ is defined for $\alpha>0,\,\,\beta>0.$
        \item $B(\alpha,\beta)=B(\beta,\alpha).$
        \item $\displaystyle B(\alpha,\beta)=\frac{\alpha -1}{\alpha+\beta -1}B(\alpha -1,\beta)=
        \frac{\beta -1}{\alpha+\beta -1}B(\alpha,\beta -1).$
        \item $\displaystyle B(\alpha, n)=\frac{(n-1)!}{\alpha(\alpha +1)\cdots (\alpha +n-1)};\,\,\,\,
        \displaystyle B(m, n)=\frac{(m-1)!(n-1)!}{(m +n-1)!}.$
        \item $\displaystyle B(\alpha,\beta)=\int_{(0,\infty)}\frac{y^{\alpha -1}}{(1+y)^{\alpha+\beta}}\lambda^1(\odif y).$
    \end{enumerate}
    \end{thm}
    \begin{thm}
        The Gamma function:
        \begin{enumerate}
        \item $\Gamma(t)$ is defined for $t>0.$
        \item $\Gamma(t)$ is infinitely differentiable and
        \[
            \Gamma^{(n)}(t)=\int_{(0,\infty)}\ln^nx\cdot x^{t-1} e^{-x}\lambda^1(\odif x)
        \]
        \item $\Gamma(t+1)=t\Gamma(t);\,\,\,\,(\Gamma(1)=1)\Longrightarrow(\Gamma(n+1)=n!).$ 
        \item $\ln\circ\Gamma$ is convex.
        \item $\displaystyle \Gamma(t)=\lim_{n\to\infty}n^t\cdot\frac{(n-1)!}{t(t+1)\cdots(t+n-1)}.$
        \item $\displaystyle \Gamma(t)\cdot\Gamma(1-t)=\frac{\pi}{\sin{\pi t}}\,\,\,\,(0<t<1).$
        \item $\displaystyle 2\int_{(0,\infty)}e^{-u^2}\lambda^1(\odif u)=\Gamma(\frac{1}{2})=\sqrt{\pi}.$
        \item $\displaystyle B(\alpha,\beta)=\frac{\Gamma(\alpha)\cdot\Gamma(\beta)}{\Gamma(\alpha+\beta)}.$
    \end{enumerate}
    \end{thm}
    \begin{enumerate}
        \item $\displaystyle
        \int_{0}^{\pi/2}\sin^{\alpha-1}\varphi\cos^{\beta-1}\varphi\odif{\varphi}=\frac{1}{2}B(\frac{\alpha}{2},\frac{\beta}{2});$
        \item $\displaystyle
        \int_{0}^{\pi/2}\sin^{\alpha-1}\varphi\odif{\varphi}=\int_{0}^{\pi/2}\cos^{\alpha-1}\varphi\odif{\varphi}=
        \frac{\sqrt{\pi}}{2}\frac{\Gamma(\frac{\alpha}{2})}{\Gamma(\frac{\alpha+1}{2})};$
        \item $\displaystyle
        V_n(r)=\frac{\pi^{n/2}}{\Gamma(\frac{n+2}{2})}r^n,$ where $V_n(r)$ is the volume of an $n$-dimensional ball of radius $r$.
    \end{enumerate}
\section{Expectation}
    \begin{define}
        The \textit{expectation} of a random variable $X:(\Omega,\mathscr{H})\rightarrow\er$ is its $\mathbb{P}$-integral.
    \end{define}
    \begin{thm}
        Let $X:(\Omega,\mathscr{H})\rightarrow (E,\mathscr{E})$ be a random variable. If $\mu:\mathscr{E}\rightarrow [0,\infty]$ is the distribution $X_*\mathbb{P}$ of $X$, then for all $\mathscr{E}/\mathscr{B}(\er)$-measurable function $u:E\rightarrow \er$,
        \begin{equation}\label{expectation}
            \mathbb{E}(u\circ X)=\int_{\mu}u
        \end{equation}
        Conversely, if \eqref{expectation} holds for $u=\mathbbm{1}_A$ for all $A$ in a $\cap$-stable generator of $\mathscr{E}$, then $\mu$ is necessarily the distribution $X_*\mathbb{P}$ of $X$.
    \end{thm}
    \begin{thm}
        Let $X:(\Omega,\mathscr{H})\rightarrow \er$ be a random variable.
        \[
            \mathbb{E} X=\int_{\mu(\odif x)}x
        \]
        where $\mu:\mathscr{B}(\er)\rightarrow [0,\infty]$ is the distribution $X_*\mathbb{P}$ of $X$.

        In particular, identically distributed random variables have the same expecation.
    \end{thm}
\subsection{Expectation of random variables valued in a convex set}
    \begin{thm}\label{expectation of random variables valued in a convex set}
        Let $X$ be an integrable $\mathbb{R}^r$-valued random variable.

        If $X$ takes values in a convex set $D$ almost surely, then $\mathbb{E}X\in D$.
    \end{thm}
    \begin{proof}
        Suppose $\mathbb{E}X\notin \cl D$. By the strict separation theorem of Hahn-Banach, there exists a continuous linear functional $f\in(\mathbb{R}^r)^\star$ and $\alpha\in\mathbb{R}$ such that $f(\mathbb{E}X)<\alpha<f(x)$ for $x\in \cl D$. Thus, $\mathbb{E}(f(X))\geqslant \alpha>f(\mathbb{E}X)=\mathbb{E}(f(X))$, contradiction. Therefore, $\mathbb{E}X\in \cl D$.

        Suppose $\mathbb{E}X\in \cl D\setminus D$. There exists a supporting hyperplane $\{f=\alpha\}$ at the boundary point $\mathbb{E}X$ such that $f(x)\geqslant \alpha$ for $x\in D$ and $f(X)\geqslant \alpha$ almost surely. But $\mathbb{E}(f(X))=f(\mathbb{E}X)=\alpha$. Thus, $f(X)=\alpha$ a.s., i.e., $X\in\{f=\alpha\}\cap D=:D'$ a.s. $D'$ is clearly convex with dimension at most $r-1$. Applying the result of the first paragraph of the proof, $\mathbb{E}X\in \cl D'$. Thus, $\mathbb{E}X\in \cl D'\setminus D'$. Repeating the process leads to a contradiction.
    \end{proof}
\subsection{Jensen's inequality}
    % \begin{thm}\label{Jensen's inequality}
    %     Let $X$ be an integrable $\mathbb{R}^r$-valued random variable and $\varphi:\mathbb{R}^r\rightarrow\mathbb{R}$ a convex function. Then $\varphi(X)$ is quasi-integrable and
    %     \[
    %         \varphi(\mathbb{E}X)\leqslant \mathbb{E}(\varphi(X))
    %     \]
    % \end{thm}
    % \begin{proof}
    %     Recall from convex analysis that $\varphi$ is necessarily continuous.
    % \end{proof}
    \begin{thm}\label{Jensen's inequality, general form}
        Let $X$ be an integrable $\mathbb{R}^r$-valued random variable and $\varphi:\mathbb{R}^r\rightarrow \,\,]\!-\!\infty,\infty]$ a l.s.c. convex function. Then $\varphi(X)$ is quasi-integrable and
        \[
            \varphi(\mathbb{E}X)\leqslant \mathbb{E}(\varphi(X))
        \]
    \end{thm}
    \begin{proof}
        Since l.s.c. functions are Borel measurable, $\varphi(X)$ is indeed measurable. If $\varphi$ is identically $\infty$, then the inequality holds trivially. Assume $\varphi$ is not identically $\infty$. Let $\varphi^*$ be the convex conjugate of $\varphi$. By Fenchel-Moreau theorem, $\varphi(x)=\sup_{f\in(\mathbb{R}^r)^\star}{\{f(x)-\varphi^*(f)\}}$ for all $x\in \mathbb{R}^r$. Since $X$ is integrable, by Theorem \ref{commutation of integral and linear map}, $f(X)$ is integrable for all $f\in(\mathbb{R}^r)^\star$. But $\varphi^*$ is not identically $\infty$, i.e., there exists $f_0\in (\mathbb{R}^r)^\star$ such that $\varphi^*(f_0)<\infty$. Thus, $\varphi(X)$ is minorized by the integrable function $f_0(X)-\varphi^*(f_0)$, and is hence quasi-integrable. Thus, by Theorem \ref{monotonicity of integral of quasi-integrable functions} and Theorem \ref{commutation of integral and linear map}, $\mathbb{E}(\varphi(X))\geqslant \mathbb{E}(f(X)-\varphi^*(f))=f(\mathbb{E}X)-\varphi^*(f)$ for all $f\in (\mathbb{R}^r)^\star$, where $f(X)-\varphi^*(f)$ is quasi-integrable since it either majorizes or minorizes the integrable function $f(X)$. Thus, $\mathbb{E}(\varphi(X))\geqslant \sup_{f\in (\mathbb{R}^r)^\star}{\{f(\mathbb{E}X)-\varphi^*(f)\}}=\varphi(\mathbb{E}X)$.
    \end{proof}
    \begin{thm}\label{Jensen's inequality}
        Let $X$ be an integrable $\mathbb{R}^r$-valued random variable and $\varphi:\mathbb{R}^r\rightarrow \mathbb{R}$ a convex function. Then $\varphi(X)$ is quasi-integrable and
        \[
            \varphi(\mathbb{E}X)\leqslant \mathbb{E}(\varphi(X))
        \]
    \end{thm}
    \begin{proof}
        $\varphi$ is necessarily continuous so Theorem \ref{Jensen's inequality, general form} applies.
    \end{proof}
    \begin{thm}\label{Jensen's inequality, for continuous convex functions on a convex set}
        Let $X$ be an integrable $\mathbb{R}^r$-valued random variable that takes values in a convex set $D$. Let $\varphi:D\rightarrow \mathbb{R}$ be a continuous convex function. Then $\varphi(X)$ is quasi-integrable, $\mathbb{E}X\in D$ and
        \[
            \varphi(\mathbb{E}X)\leqslant \mathbb{E}(\varphi(X))
        \]
    \end{thm}
    \begin{proof}
        $\varphi$ can be extended to a convex l.s.c. function $\widetilde{\varphi}:\mathbb{R}^r\rightarrow \,\,]\!-\!\infty,\infty]$. Theorem \ref{expectation of random variables valued in a convex set}.
    \end{proof}
\subsection{The AM-GM inequality}
    \begin{thm}\label{AM-GM}
        Let $\alpha_1,\dots,\alpha_n\in[0,1]$ be such that $\alpha_1+\cdots+\alpha_n=1$. For $x_1,\dots,x_n\geqslant 0$,
        \[
            \prod_{i=1}^n x_i^{\alpha_i}\leqslant \sum_{i=1}^n\alpha_ix_i
        \]
    \end{thm}
    \begin{proof}
        It suffices to prove for $x_1,\dots,x_n>0$. Let $\{1,\dots,n\}$ be equipped with its discrete $\sigma$-algebra and $\mathbb{P}(\{k\})=\alpha_k$ for $k=1,\dots,n$. Then $\mathbb{P}$ is a probability measure and $\ln x_\bullet$ is a $\mathbb{R}$-valued random variable on $\{1,\dots,n\}$ that is clearly $\mathbb{P}$-integrable. Jensen's inequality (Theorem \ref{Jensen's inequality}) then applies with $r=1$, $X=\ln x_\bullet$ and $\varphi=\exp$.
    \end{proof}
\subsection{Young's inequality}
    \begin{thm}\label{Young's inequality}
        Let $1<p,q<\infty$ be such that $p^{-1}+q^{-1}=1$. For all $x,y\geqslant 0$,
        \[
            xy\leqslant \frac{x^p}{p}+\frac{y^q}{q}
        \]
        Equality occurs if and only if $y=x^{p-1}$.
    \end{thm}
    \begin{proof}
        $xy=(x^p)^{\frac{1}{p}}(x^q)^{\frac{1}{q}}\leqslant \frac{1}{p}x^p+\frac{1}{q}x^q$ by Theorem \ref{AM-GM}.
    \end{proof}
\subsection{Expectation of functions of a random variable}
    \begin{thm}
        Let $X$ be a random variable on a probability space $(\Omega,\mathscr{H},\mathbb{P})$ taking values in a measurable space $(E,\mathscr{E})$. Let $\mu$ be the distribution of $X$ on $(E,\mathscr{E})$.

        A measurable function $u:E\rightarrow \er$ is $\mu$-quasi-integrable if and only if $u\circ X$ is $\mathbb{P}$-quasi-integrable, in which case,
        \[
            \mathbb{E}(u\circ X)=\int_\mu u
        \]
    \end{thm}
\subsection{Uniqueness theorem, the integral version}
    \begin{thm}
        Let $\mu,\nu$ be probability measures on $(E,\mathscr{E})$. Then $\mu=\nu$ if and only if
        \[
            \int_\mu u=\int_\nu u
        \]
        \begin{tight_enumerate}
            \item for $u:E\rightarrow\er_+$ that is positive measurable, or equivalently,
            \item for $u=\mathbbm{1}_A$ for all $A$ in some $\cap$-stable generator of $\mathscr{E}$, or equivalently,
            \item for $u:E\rightarrow\mathbb{R}$ positive bounded continuous, if $E$ is a metrizable topological space with $\mathscr{E}=\mathscr{B}(E)$.
        \end{tight_enumerate}
    \end{thm}
    \begin{proof}
        The first two seem to work for arbitrary measures.
    
        For a closed set $F$, the function $u_n=(1-n\cdot d(x,F))^+$ is positive bounded continuous and $u_n\downarrow \mathbbm{1}_F$. By the monotone convergence theorem (since $u_n\leqslant 1$ and $\mu,\nu$ are probability measures), $\int_\mu \mathbbm{1}_F=\int_\nu \mathbbm{1}_F$. But closed sets form a $\cap$-stable generator of $\mathscr{E}=\mathscr{B}(E)$.
    \end{proof}
\subsection{Existence of L1 convergent stochastic process}
    \begin{thm}
        Let $(\mathscr{H}_n)_{n\in\mathbb{N}}$ be a filtration of $\mathscr{H}$. For every bounded random variable $V\in \overline{\mathscr{H}_\vee}$, there exists a stochastic process $X:\Omega\rightarrow \mathbb{R}^\mathbb{N}$ adapted to $(\mathscr{H}_n)_{n\in\mathbb{N}}$ such that $X_n$ is bounded and $\|X_n-V\|_1\to 0$ as $n\to\infty$.
    \end{thm}
\section{\texorpdfstring{$L^p$}{Lp}-spaces, introduction}
    \begin{define}
        For $u:E\rightarrow \mathbb{C}$ a measurable function, define
        % \[
        %     \|u\|_p:=
        %     \begin{cases}
        %             \displaystyle\left(\int_\mu |u|^p\right)^{1/p}&\text{for }0< p<\infty\\
        %             \inf\,\{c>0\,|\,\mu(\{|u|\geqslant c\})=0\}&\text{for }p=\infty
        %     \end{cases}
        % \]
        \[
            \|u\|_p:=\left(\int_\mu |u|^p\right)^{1/p}
        \]
        for $0<p<\infty$, and
        \[
            \|u\|_\infty:=\inf{\{c\in\er_+\,|\,\mu(\{|u|> c\})=0\}}
        \]
    \end{define}
    % \begin{remark}
    %     For $p\in[1,\infty[$, the functions $\mathbb{Y}\rightarrow [0,\infty],\,\, x\mapsto |x|$ and $[0,\infty]\rightarrow [0,\infty],\,\,x\mapsto x^p$ are continuous and hence Borel-measurable. Thus, $|u|^p:X\rightarrow [0,\infty]$ is measurable.
    % \end{remark}
    \begin{thm}\label{reformulation of essential supremum}
        For $u:E\rightarrow \mathbb{C}$ any measurable function, $|u|\leqslant \|u\|_\infty$ a.e.
        
        Furthermore, if $|u|\leqslant c\in\er_+$ a.e., then $\|u\|_\infty\leqslant c$.
    \end{thm}
    \begin{proof}
        $\{|u|>\|u\|_\infty\}=\bigcup_{n=1}^\infty \{|u|>\|u\|_\infty+n^{-1}\}$, as a countable union of null sets, is null.
    \end{proof}
\subsection{H\"{o}lder's inequality}
        \begin{thm}
            Let $p,q,r\in [1,\infty]$ be such that $p^{-1}+q^{-1}=r^{-1}$, where $\infty^{-1}:=0$. 
            
            For $u,v:E\rightarrow\mathbb{C}$ measurable, 
            \[\|u\cdot v\|_r\leqslant \|u\|_p\cdot \|v\|_q\]
            If $u\in\mathcal{L}_{\mathbb{R}}^p(\mu)$ and $v\in\mathcal{L}_{\mathbb{R}}^q(\mu)$, then equality occurs if and only if
        \begin{tight_enumerate}
            \item $\displaystyle \frac{|u(x)|^p}{\left\|u\right\|_p^p}=\frac{|v(x)|^q}{\left\|v\right\|_q^q}\quad \mu\text{-a.e.}\quad\vee\quad u=0\quad \mu\text{-a.e.}\quad\vee\quad v=0\quad \mu\text{-a.e.}\quad$ for $p,q\in(0,\infty)$.
            \item $v|_{\{u\neq 0\}}(x)=\left\|v\right\|_\infty\quad \mu\text{-a.e.}\quad$ for $p=1,q=\infty$.
        \end{tight_enumerate}
        \end{thm}
        \begin{proof}
            If $r=\infty$, then necessarily $p=q=\infty$, and $|u\cdot v|=|u|\cdot |v|\leqslant \|u\|_\infty\cdot \|v\|_\infty$ a.e. By Theorem \ref{reformulation of essential supremum}, $\|u\cdot v\|_\infty\leqslant \|u\|_\infty\cdot \|v\|_\infty$. Assume $r<\infty$. If $p=\infty$, then necessarily $q=r$ and by Theorem \ref{reformulation of essential supremum} and Theorem \ref{properties of integral of positive measurable functions}, $\|u\cdot v\|_r=(\int |u|^r|v|^r)^{1/r}\leqslant (\int \|u\|_\infty^r|v|^r)^{1/r}=(\|u\|_\infty^r\int |v|^r)^{1/r}=\|u\|_\infty\cdot \|v\|_r=\|u\|_\infty\cdot \|v\|_q$. The case when $q=\infty$ is analogous. 
            
            Now assume $p,q,r<\infty$. For $0<b\leqslant 1$, the function $\varphi:\mathbb{R}_+^2\rightarrow\mathbb{R},\,\, (x,y)\mapsto x^by^{1-b}$ is continuous and concave. By Theorem \ref{Jensen's inequality, for continuous convex functions on a convex set},
        \end{proof}
        \begin{crl}
            Let $(p_n)\in ]1,\infty[^N$ such that $\sum_{n=1}^Np_n^{-1}=1$. Then $\forall (u_n)\in\mathcal{M}(\mathscr{A};\mathbb{R})^N$. 
            \[\|u_1\cdot u_2\cdot\cdots\cdot u_N\|_1\leqslant \|u_1\|_{p_1}\cdot \|u_2\|_{p_2}\cdot\cdots\cdot \|u_N\|_{p_N}\]
        \end{crl}
        \begin{thm}
            If $u\in L^p_\mathbb{K}(\beta^r)$ and $f\in C^\infty_c(\mathbb{R}^r,\mathbb{K})$, then $u\cdot f\in L^p_\mathbb{K}(\beta^r)$.
        \end{thm}
\subsection{Minkowski's inequality}
            % \begin{thm}
            %         Let $u,v\in\mathcal{L}^p(\mu;\mathbb{R}), p\in[1,\infty]$. Then
            %     \[u+v\in\mathcal{L}^p(\mu;\mathbb{R})\quad\wedge\quad\left\|u+v\right\|_p\leqslant \left\|u\right\|_p+\left\|v\right\|_p\]
            % \end{thm}
        \begin{thm}
            Let $u,v\in\mathcal{M}(\mathscr{A};\mathbb{R})$. Let $p\in[1,\infty]$. Then
        \[\left\|u+v\right\|_p\leqslant \left\|u\right\|_p+\left\|v\right\|_p\]
        \end{thm}
\subsection{Application to the sequence space}
    Let $\gamma$ be the counting measure on $(\mathbb{N},\mathscr{P}(\mathbb{N}))$. Then
    \[
        \ell^p(\mathbb{N}):=\mathcal{L}_{\mathbb{R}}^p(\gamma)=\begin{cases}
            \left\{(a_n)\in\mathbb{R}^\mathbb{N}\,\big|\, \sum\limits_{n=1}^\infty |a_n|^p<\infty\right\}\quad &p\in[1,\infty)\\
            \{\text{bounded sequences in }\mathbb{R}^\mathbb{N}\}\quad &p=\infty
        \end{cases}
    \]
    \par H\"{o}lder's inequality now reads
    \[
        \sum_{n=1}^\infty |a_nb_n|\leqslant\begin{cases}
        \displaystyle\left(\sum_{n=1}^\infty|a_n|^p\right)^{1/p}\left(\sum_{n=1}^\infty|b_n|^q\right)^{1/q}\quad &\text{if }\displaystyle\frac{1}{p}+\frac{1}{q}=1, p,q\in(1,\infty)\\
        \displaystyle\sup_{n\in\mathbb{N}}|b_n|\sum_{n=1}^\infty|a_n|\quad &\text{if }p=1,q=\infty
    \end{cases}
    \]
    \par Minkowski's inequality now reads
    \[  
        \begin{cases}
        \displaystyle\left(\sum_{n=1}^\infty |a_n\pm b_n|^p\right)^{1/p}\leqslant\left(\sum_{n=1}^\infty |a_n|^p\right)^{1/p}+\left(\sum_{n=1}^\infty |b_n|^p\right)^{1/p}\quad &\text{if }p\in[1,\infty)\\
        \displaystyle\sup_{n\in\mathbb{N}}|a_n\pm b_n|\leqslant\sup_{n\in\mathbb{N}}|a_n|+\sup_{n\in\mathbb{N}}|b_n|\quad &\text{if }p=1,q=\infty
        \end{cases}
    \]
\subsection{\texorpdfstring{$L^p$}{Lp} spaces}
    \begin{thm}
        Let $u:X\rightarrow \er$ be a measurable function. For any $p\in [1,\infty]$,
        \[
            u=0\,\,\, \mu\text{-a.e.}\Longleftrightarrow \|u\|_p=0
        \]
    \end{thm}
    \begin{proof}
        For $1\leqslant p<\infty$, by Theorem \ref{null integral}, $\|u\|_p=0\Leftrightarrow \int |u|^p\odif\mu=0\Leftrightarrow \mu(\{|u|^p>0\})=0\Leftrightarrow \mu(\{|u|>0\})\Leftrightarrow u=0\,\,\mu\text{-a.e.}$. For $p=\infty$, 
    \end{proof}
    \begin{thm}\label{u=v almost everywhere}
        Let $u,v:X\rightarrow \er$ be measurable functions.
        \[
            u=v\,\,\,\mu\text{-a.e.}\Longrightarrow \|u\|_p=\|v\|_p \text{ for all }p\in [1,\infty]
        \]
    \end{thm}
    \begin{proof}
        If $F$ is a $\mu$-full set, then $\|u\|_p^p=\int |u|^p\odif\mu=\int |u\one_F|^p\odif\mu=\|u\one_F\|_p^p$ for $1\leqslant p<\infty$.
    \end{proof}
    \begin{define}
        Let $L^p_{\mathbb{Y}}(\mu)$ be the quotient set of $\mathfrak{L}^p_{\mathbb{Y}}(\mu)$ by the equivalence relation of $\mu$-a.e. equality. That is, functions and $\mu$-almost functions that agree on a $\mu$-full set are identified. By Theorem \ref{u=v almost everywhere}, the $p$-norm of $u\in L^p_{\mathbb{Y}}(\mu)$ can be defined as the $p$-norm of any function (or $\mu$-almost function) in $\mathfrak{L}^p_{\mathbb{Y}}(\mu)$ that belongs to the class $u$. By Theorem \ref{integrals of a.e. equal positive measurable functions}, the $\mu$-integral of $u\in L^1_\mathbb{Y}(\mu)$ can be defined in a similar vein. Moreover, by Theorem \ref{integrable functions are a.e. real-valued} or Theorem \ref{integrable functions are a.e. real valued}, $L^p_\mathbb{R}(\mu)\cong L^p_{\er}(\mu)$ in \cat{Set}. Indeed, for every $u\in L^p_{\er}(\mu)$ there exists a unique $v\in L^p_\mathbb{R}(\mu)$ such that $v\subset u$.
    \end{define}
    Let $\mathbb{K}\in\{\mathbb{R},\mathbb{C}\}$.
    \begin{thm}
        $L_\mathbb{K}^p(X,\mu)$ is a vector space over $\mathbb{K}$, on which $u\mapsto \|u\|_p$ is a norm.
    \end{thm}
    \begin{thm}
        The $\mu$-integral
        \[
            \int \ph\odif \mu:L^1_\mathbb{K}(\mu)\rightarrow \mathbb{K}
        \]
        is a continuous positive linear functional.
    \end{thm}
    \begin{proof}
        $|\int u\odif\mu|/\|u\|_1\leqslant (\int |u|\odif\mu)/\|u\|_1=1$.
    \end{proof}
\subsection{\texorpdfstring{$L^p$}{Lp} spaces for a finite measure space}
    \begin{thm}\label{antitonicity of Lp spaces}
        If $\mu$ is a finite measure, then for $1\leqslant p\leqslant q\leqslant \infty$,
        \[
            L^p\supset L^q
        \]
    \end{thm}
    \begin{proof}
        H\"{o}lder's inequality.
    \end{proof}
    The reverse inclusion does not hold in general. On the probability space $([0,1],\mathscr{B}([0,1]),\lambda)$, the function $x^{-1/2}\in L^1\setminus L^2$.
\subsection{Lp functions are a.e. real-valued}
    \begin{thm}\label{integrable functions are a.e. real valued}
        Let $u\in\mathcal{L}^p_{\overline{\mathbb{R}}}(\mu)$, $p\in[1,\infty]$. Then
        \[
            \{|u|= \infty\}\in\mathscr{N}_\mu\quad\wedge\quad  u\cdot\mathbbm{1}_{\{|u|= \infty\}^c}\in\mathcal{L}^p_{\mathbb{R}}(\mu)
        \]
    \end{thm}
    \begin{proof}
        It suffices to prove for $p=1$ since $u\in\mathcal{L}^p_{\overline{\mathbb{R}}}(\mu)\Longrightarrow |u|^p\in\mathcal{L}^1_{\overline{\mathbb{R}}_+}(\mu)$ for $p\in[1,\infty)$.
        \[\mu(\{|u|= \infty\})=\mu\left(\bigcap_{n\in\mathbb{N}}\{|u|\geqslant n\}\right)=\lim_{n\to\infty}\mu(\{|u|\geqslant n\})\leqslant\lim_{n\to\infty}\left(\frac{1}{n}\int|u|\odif \mu\right)=0\]
        and note that $\{|u|^p=\infty\}=\{|u|=\infty\}$.
        \par Then use Corollary \ref{u=v almost everywhere}.2 to derive $u\cdot\mathbbm{1}_{\{|u|= \infty\}^c}\in\mathcal{L}^p(\mu;\mathbb{R})$.
    \end{proof}
\section{Uniform integrability}
    Let $(E,\mathscr{E},\mu)$ be a measure space.
    \begin{define}
        A family $F\subset\cat{Mes}((E,\mathscr{E}),\mathbb{R})$ is \textit{uniformly integrable} if 
    \end{define}
\section{Moments}
    Let $(\Omega,\mathscr{H},\mathbb{P})$ be a probability space. In this section, $L^p$ is short for $L^p_\mathbb{R}(\Omega,\mathscr{H},\mathbb{P})$, $1\leqslant p\leqslant \infty$.
    \begin{define}
        For $X$ an $\er$-valued random variable, $\mathbb{E}(X^p)$ is called the \textit{$p$th moment} of $X$, provided that $X^p$ is quasi-integrable;
        $\mathbb{E}((X-\mathbb{E}X)^p)$ is called the $p$th \textit{centered moment} of $X$, provided that $X\in L^1$ and $(X-\mathbb{E}X)^p$ is quasi-integrable.
    \end{define}
    \begin{define}
        The \textit{zero-mean} subspace of $L^p$ is the kernel $N^p$ of the expectation $\mathbb{E}$ as a linear operator $L^p\rightarrow \mathbb{R}$. Explicitly,
        \[
            N^p:=\{X\in L^p\,|\,\mathbb{E}X=0\}
        \]
        % In particular, \textit{zero-mean} subspace of $L^2$ inherits the inner product of $L^2$.
    \end{define}
    \begin{thm}
        The map
        \[
            \pi:L^p\rightarrow N^p,\,\,X\mapsto X-\mathbb{E}X
        \]
        is a linear projection of $L^p$ onto $N^p$ and induces an equivalence relation on $L^p$:
        \[
            X\sim Y:=(X-Y\text{ is a constant})
        \]
    \end{thm}
\subsection{Variances and standard deviation}
    \begin{define}
        The \textit{variance} of $X\in L^1$ is the second centered moment 
        \[
           \Var X:=\mathbb{E}((X-\mathbb{E}X)^2)=\|X-\mathbb{E}X\|_2^2
        \]
        and its \textit{standard deviation} is 
        \[
            \sigma_{X}:=\sqrt{\Var X}
        \]
        i.e., the $L^2$-norm of the zero-mean projection $X-\mathbb{E}X$.
    \end{define}
    \begin{thm}\label{properties of variances and standard deviation}
        Let $X,Y\in L^1$.
        \begin{tight_enumerate}
            \item $\Var{\alpha X}=\alpha^2\Var X$ and $\sigma_{\alpha X}=|\alpha|\sigma_{X}$ for $\alpha\in\mathbb{R}$.
            \item $\sigma_{X+Y}\leqslant \sigma_{X}+\sigma_{Y}$.
        \end{tight_enumerate}
    \end{thm}
    \begin{proof}
        2. Triangle inequality of norms.
    \end{proof}
    Let $\alpha,\beta>0$ be such that $\alpha+\beta =1$. Then by Theorem \ref{properties of variances and standard deviation}, $\sigma_{\alpha X+\beta Y}$ is generally smaller than $\alpha \sigma_X+\beta\sigma_Y$. This corresponds to the financial fact that combining assets yields a smaller volitility than the weighted average of the individual volitilities of each asset.
    \begin{thm}
        Let $X\in L^1$.
        \[
            \Var X=\mathbb{E}(X^2)-(\mathbb{E}X)^2
        \]

        Thus, $X$ has finite variance if and only if $X\in L^2$.
    \end{thm}
\subsection{Chebyshev's inequality}
    \begin{thm}
        Let $X\in L^1$.
        \[
            \mathbb{P}(\{|X-\mathbb{E}X|\geqslant\varepsilon\})\leqslant \frac{1}{\varepsilon^2}\Var{X}
        \]
    \end{thm}
    \begin{proof}
        Apply Markov's inequality to $(X-\mathbb{E}X)^2$.
    \end{proof}
\subsection{Covariances}
    \begin{define}
        The \textit{covariance} of two random variables $X,Y\in L^2$ is
        \[
            \Cov(X,Y):=\mathbb{E}((X-\mathbb{E}X)(Y-\mathbb{E}Y))=\braket{X-\mathbb{E}X\,}{\,Y-\mathbb{E}Y}
        \]
        i.e., the inner product of their projections onto the zero-mean subspace of $L^2$.
    \end{define}
    \begin{thm}\label{properties_of_covariance}
        Let $X,Y\in L^2$.
        \begin{tight_enumerate}
            \item $\Cov(X,X)=\Var X$.
            \item $\Cov(\alpha X,\beta Y)=\alpha\beta\Cov(X,Y)$ for $\alpha,\beta\in\mathbb{R}$.
            \item $\Cov(X,Y)=\mathbb{E}(XY)-\mathbb{E}X\cdot \mathbb{E}Y$.
            \item $|\Cov(X,Y)|\leqslant \sqrt{\Var X}\sqrt{\Var Y}$,
            \item $\Var(X+Y)=\Var X+\Var Y+2\Cov(X,Y)$.
        \end{tight_enumerate}
    \end{thm}
    \begin{proof}
        3. Cauchy-Schwartz inequality. 4. $\Var(X+Y)=\|X-\mathbb{E}X+Y-\mathbb{E}Y\|_2^2$.
    \end{proof}
    \begin{define}
        $X,Y\in L^2$ are said to be \textit{uncorrelated} if $\Cov(X,Y)=0$.

        The \textit{correlation} of $X,Y\in L^2$ that have non-vanishing variance is
        \[
            \Corr(X,Y):=\frac{\Cov(X,Y)}{\sqrt{\Var{X}\Var{Y}}}\in [-1,1]
        \]
    \end{define}
    \begin{thm}
        If $X_1,\dots,X_n\in L^2$ are pairwise uncorrelated, then
        \[
            \Var(X_1+\cdots+X_n)=\Var X_1+\cdots+\Var X_n
        \]
    \end{thm}
\section{Density Function of a Measure}
    \begin{thm}
        Let $f\in \mathcal{M}(\mathscr{A};\overline{\mathbb{R}}_+)$. The set function
        \[
            f\mu:\mathscr{A}\rightarrow \overline{\mathbb{R}}_+,\,\,A\mapsto \int_A f\odif \mu
        \]
        is a measure on $(X,\mathscr{A})$. It is called the measure with density function $f$ w.r.t. $\mu$.
    \end{thm}
    \begin{thm}
        \begin{tight_enumerate}
            \item If $u,w\in\mathcal{L}_{\mathbb{R}}^1(\mu)$ and $\displaystyle \int_A u\odif \mu=\int_A w\odif \mu$ for all $A\in\mathscr{A}$, then $u=w$ $\mu$-a.e.
            \item Let $f,g\in\mathcal{M}(\mathscr{A};\overline{\mathbb{R}}_+)$. Under the additional assumption that $\mu$ is $\sigma$-finite
            \[f\mu=g\mu\Longrightarrow f=g\quad\mu\text{-a.e.}\]
        \end{tight_enumerate}
    \end{thm}
    % \begin{proof}
    %     Let $A=\{f=\infty\}\cap\{g=\infty\}$. Since $\int_{A^c} |f-g|\odif\mu=\int_{A^c\cap\{f\geqslant g\}}(f-g)\odif \mu+\int_{A^c\cap\{f<g\}}(g-f)\odif \mu=0+0=0$, we have $\mu(\{f\neq g\})=\mu(A^c\cap\{f\neq g\})=0$.
    % \end{proof}
    \begin{thm}
        Let $f\in \mathcal{M}(\mathscr{A};\overline{\mathbb{R}}_+)$.
        \[u\in \mathcal{L}(f\mu;\overline{\mathbb{R}})\Longrightarrow u\cdot f\in \mathcal{L}(\mu;\overline{\mathbb{R}})\wedge\int u\odif{f\mu}=\int u\cdot f\odif \mu\]
        % in the sense that either both sides are undefined, or both sides are defined and the equality holds.
    \end{thm}
\subsection{The Radon-Nikod\'{y}m Theorem}
    Be $\mu,\nu$ two measures on the measurable space $(X,\mathscr{A})$. 
    \[\nu\ll\mu:=\mathscr{N}_\mu\subset\mathscr{N}_{\nu}\]
    \begin{thm}
        If $\mu$ is $\sigma$-finite, then the following are equivalent
        \begin{tight_enumerate}
            \item $\nu:A\mapsto \displaystyle \int_A f\odif\mu$ for some a.e. unique $f\in\mathcal{M}(\mathscr{A};\mathbb{R}_+)$.
            \item $\nu\ll\mu$.
        \end{tight_enumerate}
        The a.e. unique function $f$ is called the Radon-Nikod\'{y}m derivative of $\nu$ w.r.t. $\mu$ and denoted by $\displaystyle \frac{\odif\nu}{\odif\mu}$.
    \end{thm}
    %$\nu$ is called absolutely continuous w.r.t. $\mu$ and write $\nu\ll\mu$ if $\mathscr{N}_\mu\subset\mathscr{N}_{\nu}$.
\subsection{Lebesgue Decomposition Theorem}
    The two measures $\mu,\nu$ are called mutually singular and denoted $\mu\perp\nu$ if $\exists N\in\mathscr{A}$
    \[N\in\mathscr{N}_\mu\quad\wedge\quad N^c\in\mathscr{N}_\nu\]
    \begin{thm}
        If $\mu,\nu$ are $\sigma$-finite, then there exists a unique (up to null sets) decomposition \[\nu=\nu^\ll+\nu^\perp\]where $\nu^\ll\ll\mu$ and $\nu^\perp\perp\mu$.
    \end{thm}
\chapter{Transition Kernels and Fubini-Tonelli Theorem}
\section{Transition kernels}
    \begin{define}
        A \textit{transition kernel} from $E$ to $F$ is a function $\kappa:E\times\mathscr{F}\rightarrow [0,\infty]$ such that
        \begin{tight_enumerate}
            \item $\kappa(\ph,B):E\rightarrow [0,\infty]$ is measurable for all $B\in\mathscr{F}$.
            \item $\kappa(x,\ph):\mathscr{F}\rightarrow [0,\infty]$ is a measure for all $x\in E$.
        \end{tight_enumerate}
    \end{define}
    \begin{thm}
        For $\nu:\mathscr{F}\rightarrow [0,\infty]$ a measure and $u:E\rightarrow \er_+$ measurable,
        \[
            \kappa:E\times\mathscr{F}\rightarrow [0,\infty],\,\,(x,B)\mapsto u(x)\cdot \nu(B)
        \]
        is a transition kernel from $E$ to $F$.
    \end{thm}
    \begin{thm}
        If $(\kappa_n)_{n\in\mathbb{N}}$ is a sequence of transition kernels from $(E,\mathscr{E})$ to $(F,\mathscr{F})$, then
        \[
            \sum_{n=1}^\infty \kappa_n
        \]
        is also a transition kernel from $(E,\mathscr{E})$ to $(F,\mathscr{F})$.
    \end{thm}
    \begin{thm}
        Let $\kappa:E\times\mathscr{F}\rightarrow [0,\infty]$ be a transition kernel from $E$ to $F$. Then
        \begin{tight_enumerate}
            \item for every $u\in\overline{\mathscr{F}}_+$,
            \[
                \bracket{\kappa}{u}:E\rightarrow \er_+,\,\,x\mapsto \int_{\kappa(x,\ph)}u
            \]
            is measurable, i.e., in $\overline{\mathscr{E}}_+$.
            \item for every measure $\mu$ on $(E,\mathscr{E})$,
            \[
                \bracket{\mu}{\kappa}:\mathscr{F}\rightarrow [0,\infty],\,\,B\mapsto \int_\mu \kappa(\ph,B)
            \]
            is a measure on $(F,\mathscr{F})$.
            \item for every $u\in \overline{\mathscr{F}}_+$ and every measure $\mu$ on $(E,\mathscr{E})$,
            \begin{equation}\label{eq:associativity_of_kernel_operators}
                \int_{\bracket{\mu}{\kappa}}u=\int_{\mu}\bracket{\kappa}{u}=\int_{\mu}\int_{\kappa(x,\ph)}u\odif y\odif x
            \end{equation}
        \end{tight_enumerate}
    \end{thm}
    Let $u\in \mathscr{E}/\mathscr{F}$ and $\kappa(x,B):=\mathbbm{1}_{\{u\in B\}}$. Then \eqref{eq:associativity_of_kernel_operators} states the general transformation theorem under pushforward of measures.
\subsection{Characterization}
    \begin{thm}
        A \textit{transition kernel} from $(E,\mathscr{E})$ to $(F,\mathscr{F})$ is equivalently an operator $\kappa:\overline{\mathscr{F}}_+\rightarrow \overline{\mathscr{E}}_+$ satisfying
        \begin{tight_enumerate}
            \item $\kappa(0)=0$,
            \item $\kappa(\alpha u+\beta v)=\alpha \kappa(u)+ \beta \kappa(v)$ for all $u,v\in \overline{\mathscr{F}}_+$ and $\alpha,\beta\in\mathbb{R}_+$,
            \item $\kappa(u_n)\uparrow\kappa(u)$ whenever $u_n\uparrow u$ in $\overline{\mathscr{F}}_+$.
        \end{tight_enumerate}
    \end{thm}
    \begin{proof}
        Indeed, $\kappa(x,B)=\kappa(\mathbbm{1}_B)(x)$.
    \end{proof}
    \begin{thm}
        \[
            (\kappa_1\circ\kappa_2)(x,C)=\kappa_1(\kappa_2(\mathbbm{1}_C))(x)=\int_{\kappa_1(x,\ph)}\kappa_2(\mathbbm{1}_C)=\int_{\kappa_1(x,\odif y)}\kappa_2(y,C)
        \]
    \end{thm}
\subsection{Classification}
    \begin{define}
        A transition kernel $\kappa:\overline{\mathscr{F}}_+\rightarrow \overline{\mathscr{E}}_+$ from $(E,\mathscr{E})$ to $(F,\mathscr{F})$ is called 
        \begin{tight_enumerate}
            \item a \textit{Markov kernel} if $\kappa(1)=1$.
            \item a \textit{sub-Markov kernel} if $\kappa(1)\leqslant 1$.
            \item a \textit{finite kernel} if $\kappa(1)(x)<\infty$ for all $x\in E$.
            \item a \textit{bounded kernel} if $\kappa(1)$ is bounded.
            \item a \textit{$\sigma$-finite kernel} if $\kappa(x,\ph)$ is a $\sigma$-finite measure for all $x\in E$.
            \item a \textit{$\sigma$-bounded kernel} if there exists a countable measurable partition $F=F_1\sqcup F_2\sqcup \cdot$ such that $\kappa(\mathbbm{1}_{F_n})$ is bounded for all $n\in\mathbb{N}$.
            \item a \textit{s-finite kernel} if $\kappa=\sum_{n=1}^\infty\kappa_n$ for a sequence $(\kappa_n)$ of finite kernels.
            \item a \textit{s-bounded kernel} if $\kappa=\sum_{n=1}^\infty\kappa_n$ for a sequence $(\kappa_n)$ of bounded kernels.
        \end{tight_enumerate}
    \end{define}
    \begin{thm}
        Let $\iota:\overline{\mathscr{E}}_+\rightarrow \overline{\mathscr{E}}_+$ be the identity kernel. Then
        \[
            \iota(x,A)=\iota(\mathbbm{1}_A)(x)=\mathbbm{1}_A(x)=\delta_x(A)
        \]
    \end{thm}
    It should be obvious that
    \begin{thm}
        If $\kappa:\overline{\mathscr{E}}_+\rightarrow \overline{\mathscr{E}}_+$ is a Markov kernel on $(E,\mathscr{E})$, then so are
        \[
            \kappa^0:=\iota,\quad\kappa^1:=\kappa,\quad\kappa^{n+1}=\kappa\circ\kappa^n,\quad n\geqslant 0
        \]
    \end{thm}
    \begin{thm}
        If $\kappa$ is a finite kernel from $(E,\mathscr{E})$ to $(F,\mathscr{F})$, then
        \[
            \hat{\kappa}:(x,B)\mapsto \frac{\kappa(x,B)}{h(x)}
        \]
        is a sub-Markov kernel, where
        \[
            h:x\mapsto 
            % \begin{cases}
            %     \kappa(x,F) &\text{if }\kappa(x,F)>0\\
            %     1 &\text{otherwise}
            % \end{cases} =
            \kappa(x,F)\mathbbm{1}_{\{\kappa(\ph,F)>0\}}+\mathbbm{1}_{\{\kappa(\ph,F)=0\}}
        \]
        is measurable and strictly positive.
    \end{thm}
\subsection{Extension to measurable functions on the product space}
    \begin{thm}\label{Extension to measurable functions on the product space}
        Let $\kappa$ be a s-finite kernel from $(E,\mathscr{E})$ to $(F,\mathscr{F})$. Then for every $u\in \overline{\mathscr{E}\otimes\mathscr{F}}_+$, the function
        \[
            \bar{\kappa}(u):E\rightarrow\er_+,\,\,x\mapsto \int_{\kappa(x,\odif y)}u(x,y)
        \]
        is in $\overline{\mathscr{E}}_+$. Moreover, the operator $\bar{\kappa}:\overline{\mathscr{E}\otimes\mathscr{F}}_+\rightarrow \overline{\mathscr{E}}_+$ is linear and continuous under increasing limits, in the sense that
        \begin{tight_enumerate}
            \item $\bar{\kappa}(\alpha u+\beta v)=\alpha\bar{\kappa}(u)+\beta\bar{\kappa}(v)$,
            \item $\bar{\kappa}(u_n)\uparrow \bar{\kappa}(u)$ for $u_n\uparrow u$ in $\overline{\mathscr{E}\otimes\mathscr{F}}_+$.
        \end{tight_enumerate}
    \end{thm}
    \begin{proof}
        Measurability of $y\mapsto u(x,y)$ for any fixed $x$ follows from Theorem \ref{sections are measurable}.
        By Theorem \ref{integration w.r.t. a series of measures}, it suffices to show for $\kappa$ a finite kernel.
        For $C\in\mathscr{E}\otimes\mathscr{F}$, write
        $C_x=\{y\in F:(x,y)\in C\}\in\mathscr{F}=\{\mathbbm{1}_{C}(x,\ph)=1\}$, and define
        \[
        \mathscr{D}
        =\left\{C\in\mathscr{E}\otimes\mathscr{F}:
        x\mapsto\kappa(x,C_x)\text{ is measurable}\right\}.
        \]
        For every measurable rectangle $A\times B$,
        \[
        \kappa(x,(A\times B)_x)
        =\mathbbm{1}_A(x)\kappa(x,B),
        \]
        so $\mathscr{D}$ contains all measurable rectangles, including
        $E\times F$.
        
        If $C\in\mathscr{D}$, then
        \[
        \kappa(x,(C^c)_x)
        =\kappa(x,F)-\kappa(x,C_x)
        \]
        is measurable, since both terms are finite measurable functions.
        Thus $C^c\in\mathscr{D}$.
        If $(C_n)_{n\geqslant1}$ is a sequence of pairwise disjoint
        members of $\mathscr{D}$, then
        \[
        \kappa\left(x,\left(\bigsqcup_{n=1}^{\infty}C_n\right)_x\right)
        =\sum_{n=1}^{\infty}\kappa(x,(C_n)_x)
        \]
        is measurable. Hence $\mathscr{D}$ is a Dynkin system.
        Since measurable rectangles are $\cap$-stable, Theorem \ref{generators} gives $\mathscr{D}=\mathscr{E}\otimes\mathscr{F}$.
        
        Consequently, $\bar{\kappa}(u)$ is measurable when $u$ is
        an indicator, and hence when $u$ is a positive simple
        function, and hence for all $u$ positive measurable functions via monotone convergence theorem.
    \end{proof}
    The assumption of $s$-finiteness cannot be omitted. Let
    \[
        (E,\mathscr{E})=(\mathbb{R},\mathscr{B}(\mathbb{R})),
        \qquad
        (F,\mathscr{F})=(\mathbb{R},\mathscr{P}(\mathbb{R})),
    \]
    and let $\nu$ be counting measure on $\mathscr{P}(\mathbb{R})$. Define
    \[
        \kappa(x,B)=\nu(B),
        \qquad x\in E,\quad B\in\mathscr{F}.
    \]
    Then $\kappa$ is a kernel, since $\kappa(x,\cdot)=\nu$ is a measure
    for every $x\in E$, and $x\mapsto\kappa(x,B)$ is constant, hence
    $\mathscr{E}$-measurable, for every $B\in\mathscr{F}$.
    
    Choose a non-Borel set $A\subseteq\mathbb{R}$ and define
    \[
        u(x,y)=\mathbbm{1}_{D}(x,y)\mathbbm{1}_{A}(y),
        \qquad
        D=\{(x,y)\in\mathbb{R}^2:x=y\}.
    \]
    The diagonal $D$ is closed, so
    \[
        D\in\mathscr{B}(\mathbb{R}^2)
        =\mathscr{B}(\mathbb{R})\otimes\mathscr{B}(\mathbb{R})
        \subseteq\mathscr{E}\otimes\mathscr{F}.
    \]
    Moreover, $\mathbb{R}\times A\in\mathscr{E}\otimes\mathscr{F}$.
    Thus $u=\mathbbm{1}_{D\cap(\mathbb{R}\times A)}$ belongs to
    $\overline{\mathscr{E}\otimes\mathscr{F}}_+$. However,
    \[
        x\mapsto \int_{\kappa(x,\odif y)}u(x,y)
        =\int_{\nu(\odif y)}
            \mathbbm{1}_{\{x\}}(y)\mathbbm{1}_{A}(y)
        =\mathbbm{1}_{A}(x),
    \]
    which is not $\mathscr{E}$-measurable.
    
    Finally, $\kappa$ is not $s$-finite. Otherwise, fixing any $x\in E$
    would express $\nu$ as a countable sum of finite measures. Each
    finite measure assigns strictly positive mass to at most countably
    many singletons, so such a sum does as well. This contradicts
    $\nu(\{y\})=1$ for every $y\in\mathbb{R}$.
\subsection{Measures on the product space}
    \begin{thm}\label{measures on the product space}
        Let $\mu$ be a measure on $(E,\mathscr{E})$, and let $\kappa$
        be an s-finite transition kernel from $(E,\mathscr{E})$
        to $(F,\mathscr{F})$. Then there exists a unique measure $\pi=:\mu\times \kappa$ on
        $(E\times F,\mathscr{E}\otimes\mathscr{F})$ such that
        \begin{equation}\label{kernel product measure}
        \int_\pi u
        =\int_{\mu(\odif x)}\int_{\kappa(x,\odif y)}u(x,y)
        \qquad
        \forall u\in\overline{\mathscr{E}\otimes\mathscr{F}}_+
        \end{equation}
        Moreover, if $\mu$ is $\sigma$-finite and $\kappa$ is
        $\sigma$-bounded, then $\pi$ is $\sigma$-finite and is
        the unique measure on that product space satisfying
        \begin{equation}\label{kernel rectangular measure}
        \pi(A\times B)
        =\int_{\mu(\odif x)}\mathbbm{1}_A(x)\cdot \kappa(x,B)\qquad \forall A\in\mathscr{E}\quad\forall B\in\mathscr{F}
        \end{equation}
    \end{thm}
    \begin{proof}
        If $u_n$ is an increasing sequence in $\overline{\mathscr{E}\otimes\mathscr{F}}_+$, then $(x\mapsto \int_{\kappa(x,\odif y)}u_n(x,y))_{n\in\mathbb{N}}$ is an increasing sequence of measurable functions by Theorem \ref{Extension to measurable functions on the product space} ($s$-finiteness of the kernel is required here). Thus, repeated application of the monotone convergence theorem shows that \eqref{kernel product measure} defines a measure on the product space via Theorem \ref{Characterization of measures via integrals}.

        To prove the second, start by observing that $\pi$ satisfies
        \eqref{kernel rectangular measure}.
        Supposing that $\mu$ is $\sigma$-finite and $\kappa$ is $\sigma$-bounded,
        there remains to show that $\pi$ is $\sigma$-finite and is the only measure
        satisfying \eqref{kernel rectangular measure}.
        To that end, let $\hat{\pi}$ be another measure satisfying
        \eqref{kernel rectangular measure}.
        Since $\mu$ is $\sigma$-finite, there is a measurable partition $(E_m)$
        of $E$ such that $\mu(E_m)<\infty$ for each $m$.
        Since $\kappa$ is $\sigma$-bounded, there is a measurable partition $(F_n)$
        of $F$ such that $x\mapsto\kappa(x,F_n)$ is bounded for each $n$.
        Note that the measurable rectangles $E_m\times F_n$ form a partition
        of $E\times F$ and that, by the formula
        \eqref{kernel rectangular measure}
        for $\pi$ and $\hat{\pi}$,
        \[
            \pi(E_m\times F_n)=\hat{\pi}(E_m\times F_n)<\infty
        \]
        for each $m$ and $n$. Use Theorem \ref{uniqueness theorem}.
    \end{proof}
    \begin{thm}
        Let $\mu$ be a measure on $(E,\mathscr{E})$ and $\nu$ an s-finite measure on $(F,\mathscr{F})$. Then
        \[
            (\mu\times\nu)(A\times B)=\mu(A)\cdot \nu(B)
        \]
    \end{thm}
    \begin{thm}\label{the unique measure on the product of n spaces}
        Let $\mu$ be a measure on $(E_1,\mathscr{E}_1)$. For $i=1,\dots,n-1$, let $\kappa_i$ be
        an s-finite transition kernel from $(E_1\times\cdots\times E_i,\mathscr{E}_1\otimes\cdots\otimes\mathscr{E}_i)$ to $(E_{i+1},\mathscr{E}_{i+1})$. Then there exists a unique measure $\pi$ on the product space
        $(E_1\times\cdots\times E_n,\mathscr{E}_1\otimes\cdots\otimes \mathscr{E}_n)$ such that for all $u\in\overline{\mathscr{E}_1\otimes\cdots\otimes \mathscr{E}_n}_+$
        \begin{equation}\label{eq: the unique measure on the product of n spaces}
            \int_{\pi} u
            =\int_{\mu(\odif{x_1})}\int_{\kappa_1(x_1,\odif{x_2})}\cdots \int_{\kappa_{n-1}(x_1,\dots,x_{n-1},\odif{x_n})}u(x_1,\dots,x_{n-1},x_n)
        \end{equation}
    \end{thm}
    \begin{proof}
        Existence follows from Theorem \ref{measures on the product space} by induction. Uniqueness is immediate by setting $u=\mathbbm{1}_A$, $A\in \mathscr{E}_1\otimes\cdots\otimes \mathscr{E}_n$ in \eqref{eq: the unique measure on the product of n spaces}.
    \end{proof}
\subsection{Transition densities}
    Ex6.30 of chapter I of GTM261.
\section{Fubini-Tonelli theorem}
    \begin{thm}\label{Fubini-Tonelli}
        Let $\mu$ and $\nu$ be $s$-finite measures on $(E,\mathscr{E})$ and $(F,\mathscr{F})$ respectively. Then

        Then there exists a unique measure $\pi=\mu\otimes \nu$ on the product space $(E\times F,\mathscr{E}\otimes \mathscr{F})$ such that 
        \begin{equation}\label{Fubini's measure formula}
            \int_\pi u=\int_{\mu(\odif x)}\int_{\nu(\odif y)} u(x,y)=\int_{\nu(\odif y)}\int_{\mu(\odif x)} u(x,y)\qquad \forall u\in \overline{\mathscr{E}\otimes\mathscr{F}}_+
        \end{equation}
        Moreover, $\pi$ is $s$-finite; if $\mu$ and $\nu$ are $\sigma$-finite, then $\pi$ is $\sigma$-finite, and is the unique measure on that product space satisfying
        \begin{equation}\label{rectangular measure}
            \pi(A\times B)=\mu(A)\cdot \nu(B)\qquad \forall A\in\mathscr{E}\quad\forall B\in\mathscr{F}
        \end{equation}
        and may be denoted by $\mu\times \nu$.

        If $u\in L^1_\mathbb{K}(\pi)$, then
            \begin{gather*}
                u(\bullet,y)\in L^1_\mathbb{K}(\mu)\text{ for $\nu$-a.a. }y\in F\\
                u(x,\bullet)\in L^1_\mathbb{K}(\nu)\text{ for $\mu$-a.a. }x\in E\\
                y\mapsto \int_{\mu(\odif x)} u(x,y)\text{ is a measurable $\nu$-almost function }F\rightarrow \mathbb{K}\text{ and is in }L^1_\mathbb{K}(\nu)\\
                x\mapsto \int_{\mu(\odif y)} u(x,y)\text{ is a measurable $\mu$-almost function }E\rightarrow \mathbb{K}\text{ and is in }L^1_\mathbb{K}(\mu)\\
                \int_\pi u=\int_{\mu(\odif x)}\int_{\nu(\odif y)} u(x,y)=\int_{\nu(\odif y)}\int_{\mu(\odif x)} u(x,y)
            \end{gather*}
    \end{thm}
    \begin{proof}
        Let $h:F\times E\rightarrow E\times F,\,\,(y,x)\mapsto (x,y)$, which is clearly measurable. 
        By Theorem \ref{measures on the product space}, there exists a unique measure 
        $\pi$ on $(E\times F,\mathscr{E}\otimes \mathscr{F})$ satisfying 
        $\int_\pi u=\int_{\mu}\int_{\nu} u$ for all $u\in \overline{\mathscr{E}\otimes\mathscr{F}}_+$,
        as well as a unique measure $\tilde{\pi}$ on $(F\times E,\mathscr{F}\otimes \mathscr{E})$ satisfying $\int_{\tilde{\pi}} v=\int_{\nu}\int_{\mu} v$ for all $v\in \overline{\mathscr{F}\otimes\mathscr{E}}_+$. 
        In particular, $\pi(A\times B)=\mu(A)\cdot \nu(B)=\tilde{\pi}(B\times A)=h_*\tilde{\pi}(A\times B)$ for all $A\in\mathscr{E}, B\in\mathscr{F}$.  
        If $\mu$, $\nu$ are $\sigma$-finite, then Theorem \ref{measures on the product space} applies with the $\sigma$-bounded kernel $\kappa(x,B):=\nu(B)$ to give the second part of the theorem, from which follows $\pi=h_*\tilde{\pi}$ and hence $\int_{\mu}\int_{\nu} u=\int_{h_*\tilde{\pi}} u=\int_{\tilde{\pi}} u\circ h=\int_{\nu}\int_{\mu} u\circ h$. That is, \eqref{Fubini's measure formula} holds for $\mu,\nu$ $\sigma$-finite, and in particular, for $\mu,\nu$ finite. By Theorem \ref{integration w.r.t. a series of measures}, if $\mu=\sum_i\mu_i$ and $\nu=\sum_j\nu_j$, then $\int_\pi u=\int_{\mu}\int_{\nu} u=\sum_i\sum_j\int_{\mu_i}\int_{\nu_j}u=\sum_i\sum_j\int_{\nu_j}\int_{\mu_i}u=\sum_j\sum_i\int_{\nu_j}\int_{\mu_i}u=\int_{\nu}\int_{\mu} u$, from which it is clear that $\pi=\sum_{i,j}\mu_i\times\nu_j$ is $s$-finite, where $\mu_i\times\nu_j$ is the unique measure defined by $\int_{\mu_i\times\nu_j}u=\int_{\mu_i}\int_{\nu_j}u$.
    \end{proof}
    Without the assumption of $s$-finiteness, the theorem may not hold. Let $E=F=[0,1]$ and $\mathscr{E}=\mathscr{F}=\mathscr{B}([0,1])$. Let $\lambda$ be the Lebesgue measure on $\mathscr{E}$ and $\nu$ the counting measure on $\mathscr{F}$.
    Define
    \[
        D=\{(x,y)\in [0,1]^2:x=y\}.
    \]
    Then
    \[
        \int_{\lambda(\odif x)}\int_{\nu(\odif y)} \mathbbm{1}_{D}(x,y)=1\neq 0=\int_{\nu(\odif y)} \int_{\lambda(\odif x)}\mathbbm{1}_{D}(x,y)
    \]
    despite that $x\mapsto \int_{\nu(\odif y)} \mathbbm{1}_{D}(x,y)\equiv 1$ and $y\mapsto \int_{\lambda(\odif x)}\mathbbm{1}_{D}(x,y)\equiv 0$ are measurable.
    \begin{thm}
        If $n>d\geqslant 1$, then
        \[(\mathbb{R}^n,\mathscr{B}(\mathbb{R}^n),\lambda^n)=\left(\mathbb{R}^d\times \mathbb{R}^{n-d},\mathscr{B}(\mathbb{R}^d)\otimes \mathscr{B}(\mathbb{R}^{n-d}),\lambda^d\times \lambda^{n-d}\right)\]
    \end{thm}
    \begin{proof}
        Let $\lambda^d\times \lambda^{n-d}$ be defined by \eqref{Fubini's measure formula}. Since the Lebesgue measures are $\sigma$-finite, it is the unique measure on $\mathscr{B}(\mathbb{R}^d)\otimes \mathscr{B}(\mathbb{R}^{n-d})=\mathscr{B}(\mathbb{R}^n)$ satisfying \eqref{rectangular measure}. But $\lambda^n$ also satisfies it.
    \end{proof}
    \begin{thm}\label{Fubini-Tonelli generalized to n measures}
        Let $\mu_1,\dots,\mu_n$ be $s$-finite measures on $(E_1,\mathscr{E}_1),\dots,(E_n,\mathscr{E}_n)$ respectively. Then there exists a unique measure $\pi$ on $(E_1\times \cdots\times E_n,\mathscr{E}_1\otimes\cdots\otimes\mathscr{E}_n)$ such that
        \[
            \int_\pi u=\int_{\mu_1(\odif{x_1})}\int_{\mu_2(\odif{x_2})}\cdots \int_{\mu_n(\odif{x_n})} u(x_1,x_2,\dots,x_n)\qquad \forall u\in\overline{\mathscr{E}_1\otimes\cdots\otimes\mathscr{E}_n}_+
        \]
        and the integrals on the right side can be performed in any order desired.

        Moreover, $\pi$ is s-finite; if $\mu_1,\dots,\mu_n$ are $\sigma$-finite, then $\pi$ is $\sigma$-finite, and is the unique measure on $(E_1\times \cdots\times E_n,\mathscr{E}_1\otimes\cdots\otimes\mathscr{E}_n)$ satisfying
        \begin{equation}\label{rectangle formula for n measures}
            \pi(A_1\times\cdots\times A_n)=\mu_1(A_1)\cdots\mu_n(A_n)\qquad\forall A_1\in\mathscr{E}_1,\dots,\forall A_n\in\mathscr{E}_n
        \end{equation}
    \end{thm}
    \begin{proof}
        Existence and uniqueness of $\pi$ follows from Theorem \ref{the unique measure on the product of n spaces}. 
        Since every permutation is a composition of adjacent transpositions, it suffices to justify swapping two adjacent integrations, say $\int_{\mu_i}$ and $\int_{\mu_{i+1}}$. Set $w(x_1,\dots,x_i,x_{i+1})=\int_{\mu_{i+2}(\odif x_{i+2})}\cdots\int_{\mu_n(\odif x_n)}u(x_1,\dots,x_n)$, which is in $\mathscr{E}_1\otimes\cdots\otimes\mathscr{E}_i\otimes\mathscr{E}_{i+1}$ by repeated application of Theorem \ref{Extension to measurable functions on the product space}. Then by Theorem \ref{Fubini-Tonelli}, 
        \[
            \int_{\mu_i(\odif x_i)}\int_{\mu_{i+1}(\odif x_{i+1})}w(x_1,\dots,x_{i-1},x_{i},x_{i+1})=\int_{\mu_{i+1}(\odif x_{i+1})}\int_{\mu_i(\odif x_i)}w(x_1,\dots,x_{i-1},x_{i},x_{i+1})
        \]
        as a measurable function of $x_1,\dots,x_{i+1}$. 
        
        The uniqueness of measures satisfying \eqref{rectangle formula for n measures} is a routine application of Theorem \ref{uniqueness theorem}.
    \end{proof}
\subsection{*For Complete Sigma-Finite measures}
\subsection{The tail formula}
    \begin{thm}\label{thm:tail_formula}
        Let $\mu$ be an $s$-finite measure on $(E,\mathscr{E})$ and $u\in\mf{E}_+$.
        \begin{equation}\label{eq:tail_formula}
            \int_\mu u=\int_0^\infty \mu(\{u>t\})\odif t=\int_0^\infty \mu(\{u\geqslant t\})\odif t
        \end{equation}
        If $\varphi:\mathbb{R}_+\rightarrow\mathbb{R}_+$ with $\varphi(0)=0$ is increasing and continuously differentiable, then
        \begin{equation}\label{eq:tail_formula_generalized}
            \int_\mu \varphi\circ u=\int_0^\infty \varphi'(t)\mu(\{u> t\})\odif t=\int_0^\infty \varphi'(t)\mu(\{u\geqslant t\})\odif t
        \end{equation}
        where $\varphi(\infty):=\sup_{x\in\mathbb{R}_+}\varphi(x)$.
    \end{thm}
    \begin{proof}
        For \eqref{eq:tail_formula}, by Fubini-Tonelli theorem and Theorem \ref{measurability on a product space},
        \[
            \int_\mu u=\int_{\mu(\odif x)}\int_{\mathbb{R}_+}\mathbbm{1}_{[0,u(x)[}(t)\odif t=\int_{\mu(\odif x)}\int_{\mathbb{R}_+}\mathbbm{1}_{\{u>t\}}(x)\odif t=\int_{\mathbb{R}_+} \int_{\mu(\odif x)}\mathbbm{1}_{\{u>t\}}(x)\odif t=\int_{\mathbb{R}_+}\mu(\{u>t\})\odif t
        \]
        and the equalities remain valid if $>$ is replaced with $\geqslant$ and $\mathbbm{1}_{[0,u(x)[}$ with $\mathbbm{1}_{[0,u(x)]}$.

        For \eqref{eq:tail_formula_generalized}, by the fundamental theorem of calculus, and the assumption that $\varphi(0)=0$,
        \begin{equation}\label{eq:tail_formula_generalized_proof}
            \varphi(y)=\int_0^{y}\varphi'(t)\odif t=\int_{\mathbb{R}_+} \varphi'(t)\mathbbm{1}_{[0,y[}(t)\odif t=\int_{\mathbb{R}_+} \varphi'(t)\mathbbm{1}_{[0,y]}(t)\odif t
        \end{equation}
        for all $y\in\mathbb{R}_+$. Since $\varphi$ is increasing, $\varphi'$ is positive. By monotone convergence theorem, \eqref{eq:tail_formula_generalized_proof} holds for $y=\infty$ with $\varphi(\infty):=\sup_{x\in\mathbb{R}_+}\varphi(x)$. It then suffices to appply Fubini-Tonelli theorem.

        An alternative proof for \eqref{eq:tail_formula_generalized}. By \eqref{eq:tail_formula}, $\int_\mu \varphi\circ u=\int_0^\infty \mu(\{\varphi \circ u\geqslant t\})\odif t$. On the other hand, $\varphi'(t)\mu(\{u\geqslant t\})=\varphi'(t)\mu(\{\varphi\circ u\geqslant \varphi(t)\})$. Indeed, by isotonicity, $\{u\geqslant t\}\subset\{\varphi\circ u\geqslant \varphi(t)\}$ (with $\varphi(\infty):=\sup_{x\in\mathbb{R}_+}\varphi(x)$). If $x\in \{\varphi\circ u\geqslant \varphi(t)\}\setminus \{u\geqslant t\}$, then $u(x)<t<\infty$ and $\varphi(u(x))\leqslant \varphi(t)$, i.e., $\varphi(u(x))=\varphi(t)$. But then, by isotonicity, $\varphi(y)=\varphi(t)$ for all $y\in [u(x),t]$, and $\varphi'(t)=0$. Finally, apply transformation theorem: $\int_{\mathbb{R}_+} \varphi'(t)\mu(\{\varphi\circ u\geqslant \varphi(t)\})\odif t=\int_{\mathbb{R}_+} \mu(\{\varphi\circ u\geqslant \varphi(t)\}) \odif\varphi(t)=\int_{\mathbb{R}_+}\mu(\{\varphi\circ u\geqslant t\})\odif t$.
    \end{proof}
    \begin{thm}\label{tail formula for Lp norm}
        Let $\mu$ be an $s$-finite measure on $(E,\mathscr{E})$ and $u\in\mf{E}$. For $p>0$,
        \[\|u\|_p^p=\int_\mu |u|^p=\int_0^\infty pt^{p-1}\mu\{|u|> t\}\odif t=\int_0^\infty pt^{p-1}\mu\{|u|\geqslant t\}\odif t
        \]
    \end{thm}
    \begin{proof}
        Let $\varphi:\mathbb{R}_+\rightarrow\mathbb{R}_+,\,\,x\mapsto x^p$ in Theorem \ref{thm:tail_formula}.
    \end{proof}
\subsection{Moments of positive random variables}
    \begin{thm}
        Let $X$ be a $\er_+$-valued random variable. For $p>0$,
        \[
            \mathbb{E}(X^p)=\int_0^\infty px^{p-1}\mathbb{P}(\{X>x\})\odif x
        \]
        In particular, for $X$ taking values in $\overline{\mathbb{Z}}_+=\mathbb{N}\cup\{0,\infty\}$,
        \[
            \mathbb{E}X=\sum_{n=0}^\infty \mathbb{P}(\{X>n\}),\quad \mathbb{E}(X^2)=2\sum_{n=0}^\infty n\cdot \mathbb{P}(\{X>n\})+\mathbb{E}X
        \]
    \end{thm}
    \begin{proof}
        Theorem \ref{tail formula for Lp norm}.
    \end{proof}
\subsection{The error function}
    \[\int_{\mathbb{R}_+} e^{-t^2}\odif \lambda^1=\frac{\sqrt{\pi}}{2}\]
    \begin{proof}
        Let $I=\displaystyle \int_{\mathbb{R}_+} e^{-t^2}\odif \lambda^1$.
        \begin{align*}
            I^2&=\int_{\mathbb{R}_+} \left(\int_{\mathbb{R}_+} e^{-x^2}\lambda^1(\odif x)\right) e^{-y^2}\lambda^1(\odif y)\\
            &=\int_{\mathbb{R}_+} \int_{\mathbb{R}_+} e^{-(x^2+y^2)}\lambda^1(\odif x)\lambda^1(\odif y)
        \end{align*}
    \end{proof}
    \begin{define}
        The \textit{error function} is
        \[
            \operatorname{erf} z:=\frac{2}{\sqrt{\pi}}\int_0^z e^{-t^2}\odif t
        \]
    \end{define}
\section{Independence}
    Let $(\Omega,\mathscr{H},\mathbb{P})$ be a probability space.
    \begin{define}
        A finite family $\{\mathscr{H}_1,\dots,\mathscr{H}_n\}$ of sub-$\sigma$-algebras of $\mathscr{H}$ is an \textit{independency} if
        \[
            \mathbb{E}(V_1\cdots V_n)=\mathbb{E}V_1\cdots \mathbb{E}V_n
        \]
        for all positive random variables $V_1\in \overline{\mathscr{H}_1}_+,\dots,V_n\in \overline{\mathscr{H}_n}_+$.

        An indexed family of sub-$\sigma$-algebras of $\mathscr{H}$ is an \textit{independency} if every finite subfamily is an independency.
    \end{define}
    \begin{thm}\label{subfamily of an independency is independent}
        If $(\mathscr{H}_i)_{i\in I}$ is an independency, then so is $(\mathscr{H}_j)_{j\in J}$ for any $J\subset I$.
    \end{thm}
    \begin{thm}\label{sub-sigma-algebra of an independency is independent}
        If $(\mathscr{H}_i)_{i\in I}$ is an independency, then so is $(\mathscr{E}_i)_{i\in I}$ for $\mathscr{E}_i$ a sub-$\sigma$-algebra of $\mathscr{H}_i$ for each $i\in I$.
    \end{thm}
\subsection{Independence of generated sub-\texorpdfstring{$\sigma$}{sigma}-algebras}
    \begin{thm}\label{independence of generated sub sigma algebras}
        Let $\mathscr{F}_1,\dots,\mathscr{F}_n$ be sub-$\sigma$-algebras of $\mathscr{H}$, $n\geqslant 2$, with $\cap$-stable generators $\mathscr{C}_1,\dots,\mathscr{C}_n$ respectively.
        Then $\mathscr{F}_1,\dots,\mathscr{F}_n$ are independent if and only if
        \begin{equation}\label{eq:independence-generators}
            \mathbb{P}\left(\bigcap_{i=1}^n A_i\right)
            =
            \prod_{i=1}^n\mathbb{P}(A_i)
            \qquad
            \forall A_i\in\overline{\mathscr{C}}_i:=\mathscr{C}_i\cup\{\Omega\},\quad
            i=1,\dots,n.
        \end{equation}
    \end{thm}
    \begin{proof}
        Necessity follows immediately from
        $\overline{\mathscr{C}}_i\subset\mathscr{F}_i$.
    
        Conversely, suppose that \eqref{eq:independence-generators}
        holds. We extend this identity one coordinate at a time.
        More precisely, we prove by induction on $k\in\{0,\dots,n\}$
        that it holds whenever
        \[
            A_i\in\mathscr{F}_i\quad (i\leqslant k),
            \qquad
            A_i\in\overline{\mathscr{C}}_i\quad (i>k).
        \]
        The case $k=0$ is the assumption.
    
        Suppose that the assertion holds for $k-1$, where
        $1\leqslant k\leqslant n$. Fix
        \[
            A_i\in\mathscr{F}_i\quad (i<k),
            \qquad
            A_i\in\overline{\mathscr{C}}_i\quad (i>k),
        \]
        and put
        \[
            B:=\bigcap_{i\neq k}A_i,
            \qquad
            c:=\prod_{i\neq k}\mathbb{P}(A_i).
        \]
        Taking $A_k=\Omega$ in the induction hypothesis gives
        $\mathbb{P}(B)=c$.
    
        Define
        \[
            \mathscr{D}
            :=
            \left\{
                A\in\mathscr{F}_k:
                \mathbb{P}(A\cap B)=c\,\mathbb{P}(A)
            \right\}.
        \]
        We claim that $\mathscr{D}$ is a Dynkin system.
        First, $\Omega\in\mathscr{D}$ because $\mathbb{P}(B)=c$.
        If $A\in\mathscr{D}$, then
        \[
            \mathbb{P}(A^c\cap B)
            =
            \mathbb{P}(B)-\mathbb{P}(A\cap B)
            =
            c\bigl(1-\mathbb{P}(A)\bigr)
            =
            c\,\mathbb{P}(A^c),
        \]
        so $A^c\in\mathscr{D}$.
        If $(D_m)_{m\geqslant1}$ is a pairwise disjoint sequence
        in $\mathscr{D}$, then
        \[
            \mathbb{P}\left(
                \left(\bigcup_{m\geqslant1}D_m\right)\cap B
            \right)
            =
            \sum_{m\geqslant1}\mathbb{P}(D_m\cap B)
            =
            c\sum_{m\geqslant1}\mathbb{P}(D_m)
            =
            c\,\mathbb{P}\left(\bigcup_{m\geqslant1}D_m\right).
        \]
        Thus $\bigcup_{m\geqslant1}D_m\in\mathscr{D}$.
    
        By the induction hypothesis,
        $\overline{\mathscr{C}}_k\subset\mathscr{D}$.
        Since $\overline{\mathscr{C}}_k$ is $\cap$-stable,
        Theorem \ref{generators} yields
        \[
            \mathscr{F}_k
            =
            \sigma(\overline{\mathscr{C}}_k)
            \subset\mathscr{D}.
        \]
        This proves the induction step.
    
        Taking $k=n$, we obtain
        \eqref{eq:independence-generators} for every
        $A_i\in\mathscr{F}_i$, which proves independence.
    
        The corresponding factorization
        \[
            \mathbb{E}\left[\prod_{i=1}^n V_i\right]
            =
            \prod_{i=1}^n\mathbb{E}[V_i]
        \]
        for positive $\mathscr{F}_i$-measurable random variables
        $V_i:\Omega\rightarrow[0,\infty]$ follows first for
        positive simple functions by expanding the product,
        and then in general by increasing simple approximation
        and the monotone convergence theorem.
    \end{proof}
\subsection{Partition of an independency}
    \begin{thm}\label{partition of an independency}
        Let $(\mathscr{H}_t)_{t\in T}$ be an independent family
        of sub-$\sigma$-algebras of $\mathscr{H}$.
        Let $(T_i)_{i\in I}$ be a partition of $T$, and define for each $i\in I$
        \[
            \mathscr{H}_{T_i}
            :=
            \sigma\left(\bigcup_{t\in T_i}\mathscr{H}_t\right)
        \]
        Then $(\mathscr{H}_{T_i})_{i\in I}$ is independent.
    \end{thm}
    \begin{proof}
        For each $i\in I$, define
        \[
            \mathscr{C}_i
            :=
            \left\{
                \bigcap_{t\in J}A_t:
                J\subset T_i \text{ is finite},
                \quad A_t\in\mathscr{H}_t \text{ for every }t\in J
            \right\},
        \]
        where the intersection over the empty set is understood
        to be $\Omega$.
        Then $\Omega\in\mathscr{C}_i$, and $\mathscr{C}_i$ is a
        $\pi$-system: when intersecting two such events, events
        belonging to the same $\mathscr{H}_t$ can be combined
        using closure under finite intersections.
        Moreover,
        \[
            \bigcup_{t\in T_i}\mathscr{H}_t
            \subset \mathscr{C}_i
            \subset \mathscr{H}_{T_i},
        \]
        so $\sigma(\mathscr{C}_i)=\mathscr{H}_{T_i}$.
    
        Fix any finite collection of distinct indices
        $i_1,\dots,i_n\in I$, with $n\geqslant2$.
        By the theorem on independence of generated
        $\sigma$-algebras, it suffices to prove that
        \[
            \mathbb{P}\left(\bigcap_{k=1}^n C_k\right)
            =
            \prod_{k=1}^n\mathbb{P}(C_k)
            \qquad
            \forall C_k\in\mathscr{C}_{i_k},\quad k=1,\dots,n.
        \]
        Write
        \[
            C_k=\bigcap_{t\in J_k}A_t,
            \qquad
            J_k\subset T_{i_k}\text{ finite},
            \qquad
            A_t\in\mathscr{H}_t.
        \]
        Since the blocks $T_{i_1},\dots,T_{i_n}$ are pairwise
        disjoint, so are $J_1,\dots,J_n$.
        Independence of $(\mathscr{H}_t)_{t\in T}$ therefore gives
        \begin{align*}
            \mathbb{P}\left(\bigcap_{k=1}^n C_k\right)
            &=
            \mathbb{P}\left(
                \bigcap_{t\in J_1\cup\cdots\cup J_n}A_t
            \right)\\
            &=
            \prod_{t\in J_1\cup\cdots\cup J_n}\mathbb{P}(A_t)\\
            &=
            \prod_{k=1}^n
            \left(\prod_{t\in J_k}\mathbb{P}(A_t)\right)\\
            &=
            \prod_{k=1}^n\mathbb{P}(C_k),
        \end{align*}
        with empty products understood to be $1$.
    
        Hence $\mathscr{H}_{T_{i_1}},\dots,\mathscr{H}_{T_{i_n}}$
        are independent.
        Since every finite subfamily is independent,
        $(\mathscr{H}_{T_i})_{i\in I}$ is independent.
    \end{proof}
\subsection{Pairwise independence}
    \begin{thm}
        A sequence $(\mathscr{H}_i)_{i\in\mathbb{N}}$ of sub-$\sigma$-algebras of $\mathscr{H}$ is independent if and only if $\mathscr{H}_{\{1,\dots,n\}}$ and $\mathscr{H}_{n+1}$ are independent for all $n\in\mathbb{N}$.
    \end{thm}
    \begin{proof}
        Necessity follows immediately from Theorem \ref{partition of an independency}. For sufficiency, by Theorem \ref{independence of generated sub sigma algebras},
        \[
            \mathbb{P}(A_1\cap\cdots\cap A_n\cap A_{n+1})=\mathbb{P}(A_1\cap\cdots\cap A_n)\mathbb{P}(A_{n+1})=\cdots=\mathbb{P}(A_1)\cdots \mathbb{P}(A_n)\mathbb{P}(A_{n+1})
        \]
        for all $A_1\in\mathscr{H}_1,\dots,A_n\in\mathscr{H}_n,A_{n+1}\in\mathscr{H}_{n+1}$, $n\in\mathbb{N}$. Thus, $\mathscr{H}_1,\dots,\mathscr{H}_n,\mathscr{H}_{n+1}$ are independent for all $n\in\mathbb{N}$. But every finite subfamily of $(\mathscr{H}_i)_{i\in\mathbb{N}}$ is contained in $\mathscr{H}_1,\dots,\mathscr{H}_n,\mathscr{H}_{n+1}$ for some $n\in\mathbb{N}$. Theorem \ref{subfamily of an independency is independent}.
    \end{proof}
\subsection{Independence of random variables}
    \begin{define}
        A family $(X_i)_{i\in I}$ of random variables are \textit{independent} if $(\sigma X_i)_{i\in I}$ is an independency.
    \end{define}
    \begin{thm}\label{independence of random variables}
        Let $X_1,\dots, X_n$ be random variables taking values in $(E_1,\mathscr{E}_1),\dots,(E_n,\mathscr{E}_n)$, respectively. The following are equivalent:
        \begin{tight_enumerate}
            \item $X_1,\dots, X_n$ are independent,
            \item 
            for all $u_1\in \overline{\mathscr{E}_1}_+,\dots,u_n\in\overline{\mathscr{E}_n}_+$ holds
            \begin{equation}\label{independence of random variables in terms of expectation}
                \mathbb{E}((u_1\circ X_1)\cdots (u_n\circ X_n))=\mathbb{E}(u_1\circ X_1)\cdots\mathbb{E}(u_n\circ X_n)
            \end{equation}
            \item the joint distribution of $X_1,\dots, X_n$ is the product of their marginal distributions:
            \begin{equation}\label{independence of random variables in terms of distributions}
                (X_1,\dots,X_n)_*\mathbb{P}=((X_1)_*\mathbb{P})\times\cdots \times((X_n)_*\mathbb{P})
            \end{equation}
            \item for some $\cap$-stable generators $\mathscr{C}_1,\dots,\mathscr{C}_n$ of $\mathscr{E}_1,\dots,\mathscr{E}_n$, respectively, holds
            \begin{equation}\label{eq:criterion for independence of random variables}
                \mathbb{P}(\{X_1\in A_1\}\cap \cdots\cap \{X_n\in A_n\})=\mathbb{P}(\{X_1\in A_1\})\cdots\mathbb{P}(\{X_n\in A_n\})
            \end{equation}
            for all $A_1\in \mathscr{C}_1\cup\{E_1\},\dots,A_n\in\mathscr{C}_n\cup\{E_n\}$.
        \end{tight_enumerate}
    \end{thm}
    \begin{proof}
        $1\Leftrightarrow 2$ follows directly from definition and Theorem \ref{measurability w.r.t. the sigma algebra generated by a function}.

        $4\Rightarrow 3$. By Theorem \ref{Product of Generator is Generator of Product Sigma Algebra}, $(\mathscr{C}_1\cup\{E_1\})*\cdots *(\mathscr{C}_n\cup\{E_n\})$ is a $\cap$-stable generator of $\mathscr{E}_1\otimes\cdots\otimes\mathscr{E}_n$. Then \eqref{eq:criterion for independence of random variables} translates to that the joint distribution of $X_1,\dots,X_n$ and the product of their marginal distributions agree on this $\cap$-stable generator. Theorem \ref{uniqueness theorem for finite measures}.

        Alternatively, by Theorem \ref{sigma commutes with preimage}, $\sigma(X_i^{-1}(\mathscr{C}_i))=X_i^{-1}(\mathscr{E}_i)=\sigma X_i$ for $i=1,\dots,n$. Thus, $X_i^{-1}(\mathscr{C}_i)$ is a $\cap$-stable generator of $\sigma X_i$. Since $\{\{X_i\in A_i\}\,|\, A_i\in\mathscr{C}_i\cup\{E_i\}\}=X_i^{-1}(\mathscr{C}_i\cup\{E_i\})=X_i^{-1}(\mathscr{C}_i)\cup\{\Omega\}$ for each $i=1,\dots,n$, Theorem \ref{independence of generated sub sigma algebras} yields \eqref{eq:criterion for independence of random variables}.

        $3\Rightarrow 2$. If $u_1\in \overline{\mathscr{E}_1}_+,\dots,u_n\in\overline{\mathscr{E}_n}_+$, then $((x_1,\dots,x_n)\mapsto u_1(x_1)\cdots u_n(x_n))\in \overline{\mathscr{E}_1\otimes\cdots\otimes\mathscr{E}_n}_+$ as a composition of measurable functions. By Theorem \ref{the unique measure on the product of n spaces} and Theorem \ref{properties of integral of positive measurable functions},
        \begin{align*}
            \mathbb{E}((u_1\circ X_1)\cdots (u_n\circ X_n))&=\int_{(X_1,\dots,X_n)_*\mathbb{P}}u_1(x_1)\cdots u_n(x_n)\\
            &=\int_{(X_1)_*\mathbb{P}\,\times\cdots \times(X_n)_*\mathbb{P}}u_1(x_1)\cdots u_n(x_n)\\
            &=\int_{(X_1)_*\mathbb{P}}\cdots \int_{(X_n)_*\mathbb{P}}u_1(x_1)\cdots u_n(x_n)\\
            &=\mathbb{E}(u_1\circ X_1)\cdots\mathbb{E}(u_n\circ X_n)
        \end{align*}

        $2\Rightarrow 4$ is immediate by taking $u_i=\mathbbm{1}_{A_i}$ for $\mathscr{C}_i=\mathscr{E}_i$, $i=1,\dots,n$.
    \end{proof}
    \begin{thm}
        Measurable functions of independent random variables are independent.
    \end{thm}
    \begin{proof}
        Let $(X_i)_{i\in I}$ be an independent family of random variables taking values in $(E_i,\mathscr{E}_i)_{i\in I}$ respectively, and $u_i:(E_i,\mathscr{E}_i)\rightarrow (F_i,\mathscr{F}_i)$ be measurable functions. Then $\sigma(u_i\circ X_i)=X_i^{-1}(u_i^{-1}(\mathscr{F}_i))\subset X_i^{-1}(\mathscr{E}_i)=\sigma X_i$. Use Theorem \ref{sub-sigma-algebra of an independency is independent}.
    \end{proof}
    \begin{thm}
        If $X,Y\in L^2_\mathbb{R}$ are independent, then they are uncorrelated.
    \end{thm}
    \begin{proof}
        Let $\pi=(X,Y)_*\mathbb{P}$, $\mu=X_*\mathbb{P}$ and $\nu=Y_*\mathbb{P}$ be the distributions. By assumption, $x\in L^2_\mathbb{R}(\mu), y\in L^2_\mathbb{R}(\nu)$ and $xy\in L^1_\mathbb{R}(\pi)$. By \ref{independence of random variables in terms of distributions} and Fubini's theorem, $\mathbb{E}(XY)=\int_\pi xy=\int_{\mu\times\nu}xy=\int_\mu \int_\nu xy=\mathbb{E}X\,\mathbb{E}Y$. Thus, $\Cov(X,Y)=0$ by Theorem \ref{properties_of_covariance}.
    \end{proof}
\subsection{Kolmogolov's 0-1 law}
\subsection{Hewitt-Savage 0-1 law}
\section{Complements: Bernoulli sequences}
\chapter{\texorpdfstring{$L^p$}{Lp} Spaces}
    Let $(X,\mathscr{A},\mu)$ be a measure space. Let $\mathbb{K}\in\{\mathbb{R},\mathbb{C}\}$.
\section{The Banach Space \texorpdfstring{$L^p(\mu)$}{Lp}}
    % \begin{thm}
    %     The $p$-norm $u\mapsto \left\|u\right\|_p$ is a seminorm on the vector space $\mathfrak{L}^p_\mathbb{K}(\mu)$ for all $p\in[1,\infty]$.
    % \end{thm}
    % \begin{proof}
    %     Verify that for all $p\in[0,\infty]$
    %     \begin{enumerate}
    %         \item $\left\|u\right\|_p=0\Longleftrightarrow u=0\quad \mu\text{-a.e.}$
    %         \item $\forall \alpha\in\mathbb{R}\big(\left\|\alpha u\right\|_p=|\alpha|\left\|u\right\|_p\big)$.
    %         \item Minkowski's inequality. (See [4-13])
    %     \end{enumerate}
    %     %\par We can now conclude that $\mathcal{L}^p(\mu;\mathbb{R})$ is indeed a vector space for all $p\in[1,\infty]$, on which $\left\|\cdot\right\|_p$ is a semi-norm.
    % \end{proof}
    % $L^p_\mathbb{K}(\mu),p\in[1,\infty]$ is the quotient space
    % \[L^p(\mu;\mathbb{R}):=\mathcal{L}^p(\mu;\mathbb{R})/\mathcal{N}(\mu;\mathbb{R})=\left\{[u]\in M(\mathscr{A};\mathbb{R})\,\big|\,\int [u]\odif \mu\in\mathbb{R}\right\}\]
    % where $\mathcal{N}(\mu;\mathbb{R})=\{u\in\mathcal{M}(\mathscr{A};\mathbb{R})\,|\,u=0\quad\mu\text{-a.e.}\}=\{u\in \mathcal{L}^p(\mu;\mathbb{R})\,|\,\left\|u\right\|_p=0\}$
    % \par Be $[u]\in L^p(\mu;\mathbb{R})$ the coset of $u\in \mathcal{L}^p(\mu;\mathbb{R})$. Then 
    %     \[
    %         [v]=[u]\Longleftrightarrow v=u\quad \mu\text{-a.e.}\Longrightarrow \left\|u\right\|_p=\left\|v\right\|_p
    %     \]
    %     and we thus define $\left\|[u]\right\|_p:=\left\|u\right\|_p$ and turn $L^p(\mu;\mathbb{R}),p\in[1,\infty]$ into a bona fide normed vector space.
    % \par By Markov's inequality, $(u\in\mathcal{M}(\mathscr{A};\overline{\mathbb{R}}))\wedge(\left\|u\right\|_p<\infty)\Longrightarrow u$ is $\mu$-a.e. $\mathbb{R}$-valued, for all $p\in[1,\infty]$.
    \begin{lem}\label{p-norm triangle inequality for series}
        Let $(u_n)\in\mathcal{L}^p(\mu;\mathbb{R})^\mathbb{N}$ such that $u_n\geqslant 0$. Then
        \[\left\|\sum_{n=1}^\infty u_n\right\|_p\leqslant \sum_{n=1}^\infty\left\|u_n\right\|_p\]
    \end{lem}
    \begin{thm}
        $L^p(\mu;\mathbb{R})$ is a Banach space, $p\in [1,\infty]$.
    \end{thm}
    \begin{proof}
        Let $p\in[1,\infty)$. Be $([u_n])\in L^p(\mu;\mathbb{R})^\mathbb{N}$ a Cauchy sequence. We find a subsequence $([u_{n(k)}])$ such that
        \[\left\|[u_{n(k+1)}]-[u_{n(k)}]\right\|_p=\left\|u_{n(k+1)}-u_{n(k)}\right\|_p<2^{-k},\quad k\in\mathbb{N}\]
        \par Turn the sequence $(u_{n(k)})$ into a series by
        \[u_{n(k+1)}=\sum_{i=0}^k(u_{n(i+1)}-u_{n(i)}),\quad u_{n(0)}:=0\]
        \par By Lemma \ref{p-norm triangle inequality for series} we conclude that
        \[
            \left\|\sum_{i=0}^\infty|u_{n(i+1)}-u_{n(i)}|\right\|_p\leqslant \left\|u_{n(1)}\right\|_p+1<\infty
        \]
        which implies by Corollary \ref{integrable functions are a.e. real valued} that $N=\left\{\sum\limits_{i=0}^\infty|u_{n(i+1)}-u_{n(i)}|=\infty\right\}\in\mathscr{N}_\mu$.
        %Thus, $\sum\limits_{i=0}^\infty (u_{n(i+1)}-u_{n(i)})$ is $\mu$-a.e. convergent in $\mathbb{R}$.
        Thus $\lim\limits_{k\to\infty}u_{n(k)}=\sum\limits_{i=0}^\infty (u_{n(i+1)}-u_{n(i)})(x)$ converges (absolutely) in $\mathbb{R}$ at all points of $X\setminus N$. Then
        \[
            u:=\lim_{i\to\infty} (u_{n(i)}\cdot\mathbbm{1}_{N})=\sum_{i=0}^\infty (u_{n(i+1)}\cdot\mathbbm{1}_{N}-u_{n(i)}\cdot\mathbbm{1}_{N})\in\mathcal{L}^p(\mu;\mathbb{R})
        \]
        \par Verify that $[u_{n(k)}]\underset{k\to\infty}\to[u]$. Use \autoref{necessary and sufficient conditions for metric completeness}.
        %\begin{align*}
        %    \left\|[u]-[u_{n(k)}]\right\|_p&=\left\|[u]-[u_{n(k)}\cdot \mathbbm{1}_{N}]\right\|_p\\
        %    &=\left\|u-u_{n(k)}\cdot \mathbbm{1}_{N}\right\|_p\\
        %    &=\left\|\sum_{i=k}^\infty(u_{n(i+1)}\cdot\mathbbm{1}_{N}-u_{n(i)}\cdot\mathbbm{1}_{N})\right\|_p\\
        %    &=\left\|\left|\sum_{i=k}^\infty(u_{n(i+1)}\cdot\mathbbm{1}_{N}-u_{n(i)}\cdot\mathbbm{1}_{N})\right|\right\|_p\\
        %    &\leqslant \left\|\sum_{i=k}^\infty|u_{n(i+1)}\cdot\mathbbm{1}_{N}-u_{n(i)}\cdot\mathbbm{1}_{N}|\right\|_p\\
        %    &\leqslant \sum_{i=k}^\infty\left\|u_{n(i+1)}\cdot\mathbbm{1}_{N}-u_{n(i)}\cdot\mathbbm{1}_{N}\right\|_p\\
        %    &=\sum_{i=k}^\infty\left\|[u_{n(i+1)}]-[u_{n(i)}]\right\|_p\\
        %    &\leqslant 2^{-k+1}
        %\end{align*}
    \end{proof}
\section{Convergence Theorems in \texorpdfstring{$L^p_\mathbb{R}(\mu)$}{Lp}}
    \begin{crl}\label{Cauchy sequence in p-norm has a.e. pointwise convergent subsequence}
        Let $([u_n])\in L^p(\mu;\mathbb{R})^\mathbb{N},p\in[1,\infty)$. Let $[u]\in L^p(\mu;\mathbb{R})$.
        \[[u_n]\to[u]\Longrightarrow u_{n(k)}(x)\to u(x)\quad \mu\text{-a.e.}\quad\text{for some subsequence }(u_{n(k)})\]
         %If $[u_n]\to[u]$, then there exists a subsequence $(u_{n(k)})$ such that $u_{n(k)}(x)\to u(x)\quad \mu$-a.e.
        %\[\exists N\in\mathscr{N}_\mu\quad\exists u\in\mathcal{L}^p(\mu;\mathbb{R})\quad u_{n(k)}\cdot\mathbbm{1}_N\to u\]
        %There exists a $\mu$-a.e. pointwise convergent subsequence of $(u_n)$.
    \end{crl}
    \begin{thm}
        Let $(u_n)\in \mathcal{L}^p(\mu;\mathbb{R})^\mathbb{N}, p\in[1,\infty)$. If
        \begin{enumerate}
            \item $\exists w\in \mathcal{L}^p(\mu;\mathbb{R})\,\,\forall n\in\mathbb{N}\,\,(|u_n|\leqslant w\quad\mu\text{-a.e.})$
            \item $\exists\lim\limits_{n\to\infty}u_n(x)\in\mathbb{R}\quad\mu\text{-a.e.}$
        \end{enumerate}
        then
        \[N=\left\{x\in X:\nexists\lim_{n\to\infty}u_n(x)\right\}\cup\bigcup_{n\in\mathbb{N}}\{|u_n|>w\}\in\mathscr{N}_\mu\]
        and $u=\lim_{n\to\infty}\limits u_n\cdot \mathbbm{1}_{N^c}\in \mathcal{L}^p(\mu;\mathbb{R})$ and $[u_n]\to[u]$.
    \end{thm}
    \begin{proof}
        $|u|\leqslant w\Longrightarrow |u|^p\leqslant w^p\Longrightarrow u\in \mathcal{L}^p(\mu;\mathbb{R})$. Since
        \begin{enumerate}
            \item $|u_n\cdot \mathbbm{1}_{N^c}-u|^p\leqslant (|u_n\cdot \mathbbm{1}_{N^c}|+|u|)^p\leqslant 2^p w^p$.
            \item $|u_n\cdot \mathbbm{1}_{N^c}-u|^p\to 0$.
            \item $|u_n\cdot \mathbbm{1}_{N^c}-u|^p\in \mathcal{L}^1(\mu;\mathbb{R})$.
            \item $2^p w^p\in \mathcal{L}^1(\mu;\mathbb{R})$.
        \end{enumerate}
        By dominated convergence, we conclude that
        \[
            \int |u_n\cdot \mathbbm{1}_{N^c}-u|^p\odif \mu\to 0\Longleftrightarrow \|u_n\cdot \mathbbm{1}_{N^c}-u\|_p\to 0\Longleftrightarrow [u_n\cdot \mathbbm{1}_{N^c}]\to [u]\Longleftrightarrow [u_n]\to [u]
        \]
    \end{proof}
    \begin{thm}
        Let $(u_n)\in \mathcal{L}^p(\mu;\mathbb{R})^\mathbb{N}, p\in[1,\infty)$. 
        \par If $u_n(x)\to u(x)\quad\mu$-a.e. for some $u\in \mathcal{L}^p(\mu;\mathbb{R})$, then
        \[[u_n]\to [u]\Longleftrightarrow \|u_n\|_p\to\|u\|_p\]
    \end{thm}
    \begin{thm}
        Let $(u_n)\in (L^1)^\mathbb{N}$. If $u_n\to u\in L^1$ as $n\to\infty$, then
        \[
            \int_\mu u_n\to\int_\mu u
        \]
    \end{thm}
    \begin{proof}
        This is actually trivial, since $\int_\mu$ is a bounded linear functional on $L^1$.
    \end{proof}
\section{Dense subspaces}
    Recall that \textit{sequential density} implies density and they are equivalent for first-countable spaces.
\subsection{Integrable simple functions}
    \begin{thm}
        A simple function $e$ is $\mu$-integrable if and only if $\mu(\{e\neq 0\})<\infty$.
    \end{thm}
    \begin{thm}
        $\{\text{integrable simple functions}\}\subset \bigcap_{1\leqslant p\leqslant \infty}L^p$.
    \end{thm}
    \begin{thm}
        For $1\leqslant p<\infty$, the subspace of integrable simple functions is dense in $L^p$.
    \end{thm}
    \begin{proof}
        If $u$ is integrable, so is $u^+,u^-$, to which by sombrero lemma there exist increasing sequences $(u^+_n),(u^-_n)$ of positive simple integrable functions tending. By dominated convergence, $\lim_{n\to\infty}\|u^+_n-u^+\|_p=\lim_{n\to\infty}\|u^-_n-u^-\|_p=0$. By triangle inequality, $(u_n^+-u_n^-)$ converges to $u$ in $L^p$.
    \end{proof}
    \begin{thm}
        For $1\leqslant p<\infty$, for a vector subspace $V\subset L^p$ to be dense in $L^p$ it suffices that for every measurable set $A$ with $\mu(A)<\infty$ and for all $\varepsilon>0$ there exists $v\in V$ such that $\|v-\mathbbm{1}_A\|_p<\varepsilon$.
    \end{thm}
    \begin{proof}
        Since $V$ is a vector subspace, by triangle inequality, every integrable simple function $e$ can be approximated arbitrarily closely by elements of $V$. Given $u\in L^p$, choose $(e_n)\in E^\mathbb{N}$ such that $\|e_n- u\|_p\to 0$. For each $e_n$ choose $v_n\in V$ such that $\|v_n-e_n\|_p<\|e_n- u\|_p$. Then $\|v_n-u\|_p\to 0$.
    \end{proof}
\subsection{Bounded continuous functions}
    \textbf{Let $(X,d)$ be a finite metric space and $\mathscr{A}=\mathscr{B}(X)$ for the remainder of this section.}
    \begin{lem}
        Be $G$ an open set in $X$ with $\mu(G)<\infty$.
        \[f\in\mathcal{E}(X)\wedge (\{f\neq 0\}\subset G)\Longrightarrow\exists (\varphi_n)\in(\mathcal{C}_b(X;\mathbb{R})\cap \mathcal{L}^p_\mathbb{R}(\mu))^\mathbb{N}\quad [\varphi_n]\to[f]\]
        %where $\widetilde{f}\in\mathcal{E}(X;\mathbb{R})$ is the extension of $f$ onto $X$ such that $\widetilde{f}|_{G^c}=0$. 
    \end{lem}
    % \begin{thm}
    %     If $\mu$ is finite, then $\mathcal{C}_b(X;\mathbb{R})\subset\mathcal{L}^p(\mu;\mathbb{R})$ and $\{[f]\in L^p(\mu;\mathbb{R})\,|\,f\in\mathcal{C}_b(X;\mathbb{R})\}$ is dense in $L^p(\mu;\mathbb{R})$.
    % \end{thm}
    \begin{thm}
        For $C_b(X;\mathbb{R})\cap L^p_\mathbb{R}(\mu)$ to be dense in $L^p_\mathbb{R}(\mu)$, it suffices that one of the following holds:
        \begin{tight_enumerate}
            \item $\mu$ is finite, in which case, $C_b(X;\mathbb{R})\subset L^p_\mathbb{R}(\mu)$.
            \item $\mu(B(0,r))<\infty$ for all $r>0$.
            \item $\mu$ is outer regular.
        \end{tight_enumerate}
    \end{thm}
    % \begin{crl}
    %     If $\mu(B(0,r))<\infty$ for all $r>0$, then $\{[f]\in L^p(\mu;\mathbb{R})\,|\,f\in\mathcal{C}_b(X;\mathbb{R})\cap\mathcal{L}^p(\mu;\mathbb{R})\}$ is dense in $L^p(\mu;\mathbb{R})$.
    % \end{crl}
    % \begin{thm}
    %     If $\mu$ is outer regular, then $\{[f]\in L^p(\mu;\mathbb{R})\,|\,f\in\mathcal{C}_b(X;\mathbb{R})\cap\mathcal{L}^p(\mu;\mathbb{R})\}$ is dense in $L^p(\mu;\mathbb{R})$.
    % \end{thm}
\subsection{Compactly supported continuous functions}
    \begin{thm}
        If $X$ is $\sigma$-compact and $\mu(K)<\infty$ for all compact subspaces $K$, then $C_c(X;\mathbb{R})\subset L^p_\mathbb{R}(\mu)$ and for $C_c(X;\mathbb{R})$ to be dense in $L^p(\mu;\mathbb{R})$, it suffices that either of the following holds:
        \begin{tight_enumerate}
            \item $\mu(B(0,r))<\infty$ for all $r>0$.
            \item $\mu$ is outer regular.
        \end{tight_enumerate}
    \end{thm}
    \begin{thm}
        If $X$ is $\sigma$-compact and $\mu(K)<\infty$ for all compact subspaces $K$, then $C_c(X;\mathbb{R})\subset L^p_\mathbb{R}(\mu)$ and for $C_c(X;\mathbb{R})$ to be dense in $L^p(\mu;\mathbb{R})$, it suffices that either of the following holds:
        \begin{tight_enumerate}
            \item $\mu(B(0,r))<\infty$ for all $r>0$.
            \item $\mu$ is outer regular.
        \end{tight_enumerate}
    \end{thm}
    \begin{crl}
        Be $\mu$ a measure on $(\mathbb{R}^n,\mathscr{B}(\mathbb{R}^n))$ such that $\mu(K)<\infty$ for all compact subspaces $K$. Then $C^\infty_c(\mathbb{R}^n;\mathbb{R})$ is dense in $L^p_\mathbb{R}(\mu)$.
    \end{crl}
\subsection{Continuous functions vanishing at infinity}
    There seem to be two different notions of vanishing at infinity. Check out.
    \begin{thm}
        $C_0\subset L^\infty$, at least for the locally compact space definition of vanishing at infinity.
    \end{thm}
    \begin{thm}
        $C_c$ is dense in $C_0$.
    \end{thm}
    \begin{proof}
        Proof required.
    \end{proof}
    \begin{thm}
        $C_0$ is the completion of $C_c$ under the $\infty$-norm. %$\|\ph\|_\infty$.
    \end{thm}
    \begin{proof}
        Proof required.
    \end{proof}
\section{Riesz-Markov-Kakutani Representation Theorem}
\subsection{Positive Linear Functional}
    A positive linear functional on an ordered vector space $(V,\leqslant)$ is a linear functional $\Lambda$ such that
    \[\Lambda(v)\geqslant 0\quad\forall v\geqslant 0\]
    \begin{lem}
        Be $\Lambda:V\rightarrow\mathbb{R}$ a positive linear functional.
        \begin{tight_enumerate}
            \item $\Lambda$ is increasing, i.e. $u\leqslant v\Longrightarrow \Lambda(u)\leqslant\Lambda(v)$.
            \item $|\Lambda u|\leqslant \Lambda |u|$.
            \item If $V$ is a Banach space of functions or equivalence classes of functions such that $f\in V\Longrightarrow |f|\in V\wedge \||f|\|=\|f\|$, then $\Lambda$ is bounded.
        \end{tight_enumerate}
    \end{lem}
    \begin{thm}
        Let $V\in\{L^p_{\mathbb{R}}(\mu),(C_c(X;\mathbb{R}),\|\cdot\|_\infty),(C_b(X;\mathbb{R}),\|\cdot\|_\infty),(C_\infty(X;\mathbb{R}),\|\cdot\|_\infty)\}$. Be $\Lambda:V\rightarrow \mathbb{R}$ a bounded linear functional. Then there exists positive linear funcitonals $\Lambda^{\pm}:V\rightarrow\mathbb{R}$ such that $\Lambda=\Lambda^+-\Lambda^-$.
    \end{thm}
\subsection{Duality of the Spaces \texorpdfstring{$L^p_\mathbb{R}(\mu),1\leqslant p<\infty$}{Lp, p<infty} }
    Be $(X,\mathscr{A},\mu)$ a $\sigma$-finite measure space. Let $p\in[1,\infty),q\in(1,\infty]$ such that $p^{-1}+q^{-1}=1$.
    \[I_f:L^p_\mathbb{R}(\mu)\rightarrow\mathbb{R},u\mapsto \int uf\odif \mu,\quad f\in L^q_{\mathbb{R}_+}(\mu)\]
    defines a positive linear functional on $L^p_\mathbb{R}(\mu)$ bounded by $|I_fu|\leqslant\|f\|_q\|u\|_p$.
    \begin{thm}
        For every positive linear functional $\Lambda:L^p_\mathbb{R}(\mu)\rightarrow\mathbb{R}$, there exists a unique $f\in L^q_{\mathbb{R}_+}(\mu)$ such that $\Lambda=I_f$.
    \end{thm}
\chapter{Conditioning}
    Let $(\Omega,\mathscr{H},\mathbb{P})$ be a probability space.
\section{Conditional expectation}
    Let $\mathscr{E}\subset\mathscr{H}$ be a sub-$\sigma$-algebra of $\mathscr{H}$.
    \begin{define}
        The \textit{conditional expectation} of a positive random variable $X\in\overline{\mathscr{H}}_+$ given $\mathscr{E}$ is any positive random variable $X'\in \overline{\mathscr{E}}_+$ such that for all $V\in \overline{\mathscr{E}}_+$,
        \[
            \mathbb{E}(V X)=\mathbb{E}(V X')
        \]
    \end{define}
    \begin{thm}
        $X'\in \overline{\mathscr{E}}_+$ is the conditional expectation of $X\in\overline{\mathscr{H}}_+$ if and only if for all $A\in\mathscr{E}$
        \[
            \mathbb{E}(\mathbbm{1}_AX)=\mathbb{E}(\mathbbm{1}_AX')
        \]
    \end{thm}
\subsection{Uniqueness}
    \begin{thm}
        If both $X'$ and $X''$ are conditional expectations of $X\in\overline{\mathscr{H}}_+$ given $\mathscr{E}$, then $X'=X''$ almost surely.
    \end{thm}
    \begin{proof}
        
    \end{proof}
\chapter{Transformation Theorems}
\section{The general transformation theorem}
    Let $(E,\mu)$ be a measure space and $F$ a measurable space. Let $\hat{\mathbb{K}}\in\{\mathbb{R},\mathbb{C},\er,[0,\infty]\}$.
    \begin{thm}
        Let $f:E\rightarrow F$ and $u:F\rightarrow \er_+$ be measurable functions.
        \[\int_\mu f^* u=\int_{f_*\mu} u\]
        where $f^* u=u\circ f$ is the pullback of $u$ by $f$.
    \end{thm}
    \begin{proof}
        It suffices to prove for $\hat{\mathbb{K}}=[0,\infty]$ since $(u\circ f)^\pm=u^\pm\circ f$.
    \end{proof}
    \begin{thm}
        Let $f:T\rightarrow X$ be a \cat{Mes} isomorphism and $u\in L^1_{\mathbb{C}}(X)$.
        \[
            \int u\odif \mu=\int f^*u\odif f^{-1}_*\mu
        \]
    \end{thm}
    \begin{crl}\label{transformation theorem for affine automorphisms}
        Let $f\in\Aut_\cat{Aff}(\mathbb{R}^r)$ and $u:\mathbb{R}^r\rightarrow \hat{\mathbb{K}}$ a $\mathscr{L}^r/\mathscr{B}\hat{\mathbb{K}}$-measurable function.
        \begin{gather*}
            \int f^*u\odif\lambda^r=|\det f|^{-1}\int u \odif \lambda^r\\
            \int f^*u\odif\beta^r=|\det f|^{-1}\int u \odif \beta^r\\
            \int u\odif\lambda^r=|\det f|\int f^*u\odif \lambda^r\\
            \int u\odif\beta^r=|\det f|\int f^*u \odif \beta^r
        \end{gather*}
    \end{crl}
    \begin{proof}
        Recall that $f$ is $\mathscr{L}^r/\mathscr{L}^r$-measurable and $f_*\lambda^r=|\det f|^{-1}\lambda^r$.
    \end{proof}
\subsection{*Transforming \texorpdfstring{$\mu$}{mu}-Integral into Stieltjes Integrals}
\subsection{Transformation of Measures with Densities}
    \begin{thm}
        For all measurable $g:Y\rightarrow [0,\infty]$,
        \[
            g\odot f_*\mu=f_*((f^*g)\odot \mu)
        \]
    \end{thm}
    \begin{crl}
        If $f:X\rightarrow Y$ is an \cat{Mes} isomorphism, then for all measurable $h:X\rightarrow [0,\infty]$,
        \[
            f_*(h\odot \mu)=(h\circ f^{-1})\odot f_*\mu
        \]
    \end{crl}
\section{Jacobi's Transformation Theorem}
    \begin{thm}\label{Jacobi Transformation Theorem}
        Be $\varphi:U\rightarrow V$ a $\mathcal{C}^1$-diffeomorphism between open sets $U,V\subset\mathbb{R}^n$. Then
        \[
            \lambda^n|_V=\varphi_*(|\det \varphi'(\cdot)|\lambda^n|_U)
        \]
        or, equivalently,
        \[\int u\odif \lambda^n|_V=\int (u\circ \varphi)\cdot|\det\varphi'|\odif \lambda^n|_U,\quad\forall u\in\mathcal{M}(V;\overline{\mathbb{R}}_+)\]
    \end{thm}

    % \begin{lem}
    %     Norm $\mathbb{R}^n$ and $\mathbb{R}^d$ with the maximum norm $\|\cdot\|_\infty$.
    %     Be $\varphi=(\varphi_1,\cdots,\varphi_d):\mathbb{R}^n\rightarrow \mathbb{R}^d$ an $\alpha$-H\"{o}lder continuous map of constant $C$, where $\alpha\in(0,1]$. Then for all $x\in\mathbb{R}^n$ and $s>0$ one has
    %     \[\varphi\left(\bigtimes_{i=1}^n[x_i-s,x_i+s]\right)\subset\bigtimes_{j=1}^d[\varphi_j(x)-Cs^\alpha,\varphi_j(x)+Cs^\alpha]\]
    %     In particular, for all $n$-squares $I$,
    %     \[
    %         \lambda^d(\varphi(I))\leqslant (2^{1-\alpha}C)^d(\lambda^n(I))^{\alpha d/n}
    %     \]
    % \end{lem}

    \begin{lem}
        \autoref{Jacobi Transformation Theorem} holds for affine linear maps $\varphi$.
    \end{lem}
    \begin{lem}
        Norm $\mathbb{R}^n$ and $\mathbb{R}^d$ with the maximum norm $\|\cdot\|_\infty$.
        Be $\varphi=(\varphi_1,\cdots,\varphi_d):B(x,r)\rightarrow \mathbb{R}^d$ an $\alpha$-$K$-H\"{o}lder continuous map, where $\alpha\in(0,1]$. Then
        \[\varphi\left(B(x,r)\right)\subset B(\varphi(x),Kr^\alpha)\]
        In particular,
        \[
            \lambda^d(\varphi(B(x,r)))\leqslant (2^{1-\alpha}K)^d(\lambda^n(B(x,r)))^{\alpha d/n}
        \]
    \end{lem}
    \begin{lem}
        Be $\mu,\nu$ measures on the measurable space $(X,\sigma(\mathscr{S}))$, where $\mathscr{S}$ is a semiring.
        \[(\mathscr{S}\text{ exhausts }X)\wedge (\mu|_\mathscr{S}\leqslant\nu|_\mathscr{S}<\infty)\Longrightarrow \mu\leqslant\nu\]
    \end{lem}
    \begin{lem}
        Be $\varphi:U\rightarrow V$ a $\mathcal{C}^1$-diffeomorphism between open sets $U,V\subset\mathbb{R}^n$. Then
        \[\lambda^n|_V(\varphi(I))\leqslant \int_I |\det \varphi'(x)|\odif \lambda^n|_U\]
        for all $I\in\llbracket \mathbb{Q}^n,\mathbb{Q}^n\llbracket$ such that $\cl I\subset U$.
    \end{lem}

    % \begin{proof}
    %     Let $\|x\|_\infty:=\max_{1\leqslant i\leqslant n}|x^i|$ for $x\in\mathbb{R}^n$. Since $\cl I$ is compact, $\varphi'|_{\cl I}$ is uniformly continuous. Fix $\varepsilon>0$, we find $\delta>0$ such that $\sup_{\|x-x'\|_\infty\leqslant\delta,\,\,x,x'\in\cl I}\|\varphi'(x)-\varphi'(x')\|\leqslant {\varepsilon}/{C}$, where $C:=\sup_{x\in\cl I}\|(\varphi')^{-1}(x)\|$.
    %     \par Partition $I=\bigcup_{j=1}^NI_j$ into disjoint half-open $n$-squares $I_j\in \llbracket \mathbb{Q}^n,\mathbb{Q}^n\llbracket$ with side-length $<\delta$. Since $\det\varphi'$ is continuous, we find for each $j=1,\cdots,N$ a point $t_j\in \cl I_j$ such that $|\det\varphi'(t_j)|=\inf_{t\in I_j}|\det\varphi'(t)|$.
    %     \par Let $A_j:=\varphi'(t_j)$. Then
    %     \[(A_j^{-1}\circ \varphi)'=A_j^{-1}\circ\varphi'=I+A_j^{-1}\circ(\varphi'(\cdot)-\varphi'(t_j))\]
    %     Thus,
    %     \[\sup_{t\in\cl I_j}\|(A_j^{-1}\circ \varphi)'(t)\|\leqslant \sup_{t\in\cl I_j}\|I\|+\|A_j^{-1}\|\cdot\|\varphi'(t)-\varphi'(t_j)\|\leqslant 1+C\cdot\frac{\varepsilon}{C}=1+\varepsilon\]
    % \end{proof}

\section{Generalization of the Transformation Theorem}
    \begin{lem}
        Be $\varphi:I\rightarrow \mathbb{R}^d$ an $\alpha$-H\"{o}lder continuous map, where $I$ is a closed rectangle with rational endpoints, $\alpha\in(0,1]$. If $\alpha d\geqslant n$, then
        %\[\varphi(N^*)\in\mathscr{N}_{\overline{\lambda^d}},\quad\forall N^*\in\mathscr{N}_{\overline{\lambda^n}}\]
        \[N^*\in\mathscr{N}_{\overline{\lambda^n}|_I}\Longrightarrow \varphi(N^*)\in\mathscr{N}_{\overline{\lambda^d}}\]
    \end{lem}
    \begin{thm}
        Be $\varphi:U\rightarrow V$ a $\mathcal{C}^1$-diffeomorphism between open sets $U,V\subset\mathbb{R}^n$. Then
        \[
            \overline{\lambda^n}|_V=\varphi_*(|\det \varphi'|\overline{\lambda^n}|_U)
        \]
    \end{thm}
    \begin{proof}
        Let $N^*\in\mathscr{N}_{\overline{\lambda^n}|_V}$. Since $\varphi^{-1}$ is $\mathcal{C}^1$, it is Lipschitz continuous on closed rectangles $I_n$ contained in $V$ with rational endpoints, but the (countable) union of which is $V$ and $\varphi^{-1}(N^*\cap I_n)\in\mathscr{N}_{\overline{\lambda^n}|_U}$. Then $\varphi^{-1}(N^*)=\varphi^{-1}(\bigcup_{n\in\mathbb{N}}(N^*\cap I_n))=\bigcup_{n\in\mathbb{N}}\varphi^{-1}(N^*\cap I_n)\in\mathscr{N}_{\overline{\lambda^n}|_U}$.
    \end{proof}
    \begin{thm}
        Let $\varphi\in\mathcal{C}^1(U;\mathbb{R}^n)$, where $U\subset\mathbb{R}^n$. If $\varphi|_{U\setminus \{\det\varphi'\neq 0\}}$ is injective, then
        \[\int_{\varphi(U)}u\odif{\overline{\lambda^n}}=\int_U (u\circ\varphi)\cdot|\det\varphi'|\odif{\overline{\lambda^n}}\quad\forall u\in\mathcal{M}(\varphi(U);\overline{\mathbb{R}}_+)\]
    \end{thm}
\section{Image of Borel Sets}
    \begin{lem}
        Let $\varphi\in\mathcal{C}(\mathbb{R}^n;\mathbb{R}^d)$.
        \[F\text{ is an $F_\sigma$-set in }\mathbb{R}^n\Longrightarrow \varphi(F)\text{ is an $F_\sigma$-set in }\mathbb{R}^d\]
    \end{lem}
    \begin{thm}
        Let $B\in\mathscr{B}(\mathbb{R}^n)$. There exists an $F_\sigma$-set $F$ and a $G_\delta$-set $G$ such that
        \[F\subset B\subset G\quad\wedge\quad\lambda^n(F)=\lambda^n(B)=\lambda^n(G)\]
    \end{thm}
    \begin{thm}\label{Lusin}
        Let $\varphi\in\mathcal{C}(\mathbb{R}^n;\mathbb{R}^d)$. The following assertions are equivalent
        \begin{enumerate}
            \item $\varphi(N^*)\in\mathscr{N}_{\overline{\lambda^d}},\quad\forall N^*\in\mathscr{N}_{\overline{\lambda^n}}$.
            \item $\varphi(B^*)\in\overline{\mathscr{B}(\mathbb{R}^d)},\quad\forall B^*\in\overline{\mathscr{B}(\mathbb{R}^n)}$.
        \end{enumerate}
    \end{thm}
    \begin{crl}
        Be $\varphi:\mathbb{R}^n\rightarrow \mathbb{R}^d$ an $\alpha$-H\"{o}lder continuous map, where $\alpha\in(0,1]$. If $\alpha d\geqslant n$, then $\varphi\left[\overline{\mathscr{B}(\mathbb{R}^n)}\right]\subset \overline{\mathscr{B}(\mathbb{R}^d)}$ and
        \[\overline{\lambda^d}(\varphi(B^*))\leqslant (2^{1-\alpha}C)^d\overline{\lambda^n}(B^*),\quad\forall B^*\in\overline{\mathscr{B}(\mathbb{R}^n)}\]
    \end{crl}
\section{Polar Coordinates and the Volume of the Unit Ball}
    Let $\psi:]0,\infty[\,\times\, ]0,\pi[^{n-2}\,\times\, ]-\pi,\pi[\,\rightarrow \mathbb{R}^n\setminus\{x\in\mathbb{R}^n\,|\,x^n=0,x^{n-1}\leqslant 0\}$,
    \[
    \begin{pmatrix}
        r\\
        \bm{\theta}
    \end{pmatrix}=\begin{pmatrix}
        r\\
        \theta_1\\
        \theta_2\\
        \vdots\\
        \theta_{n-2}\\
        \theta_{n-1}
    \end{pmatrix}\mapsto\begin{pmatrix}
        r\cos\theta_1\\
        r\sin\theta_1\cos\theta_2\\
        r\sin\theta_1\sin\theta_2\cos\theta_3\\
        \vdots\\
        r\sin\theta_1\sin\theta_2\cdots\sin\theta_{n-2}\cos\theta_{n-1}\\
        r\sin\theta_1\sin\theta_2\cdots\sin\theta_{n-2}\sin\theta_{n-1}
    \end{pmatrix}\]
    \begin{lem}
        $\psi$ is a diffeomorphism and
        \[\det\psi'(r,\bm{\theta})=r^{n-1}\sin^{n-2}\theta_1\sin^{n-3}\theta_2\cdots\sin\theta_{n-2}\]
    \end{lem}
\section{Surface Measure on the Sphere}
    Be $\mathbb{S}^{n-1}:=\partial B(0,1)$ the $(n-1)$-dimensional unit sphere in $\mathbb{R}^n$. Set
    \[\rho:\mathbb{R}^n\setminus\{0\}\rightarrow \,]0,\infty[\,\times\mathbb{S}^{n-1},\quad x\mapsto (|x|,e(x)):=(|x|,x/|x|)\]
    \begin{lem}
        $\rho$ is a diffeomorphism.
    \end{lem}
    Introduce the surface measure $\sigma_{n-1}$ on $\mathbb{S}^{n-1}$
    \[\sigma_{n-1}(A):=n\rho_*(\lambda^n)(\{1\}\times A)=n\lambda^n(e^{-1}(A)\cap B(0,1)),\quad\forall A\in\mathscr{B}(\mathbb{S}^{n-1})\]
    \begin{thm}
        $\sigma_{n-1}$ is invariant under rotations and
        \[\lambda^n=\rho^{-1}_*(f_{n-1}\lambda^1\times \sigma_{n-1}),\quad f_{n-1}(r)=r^{n-1}\]
        or, equivalently,
        \[\int u\odif\lambda^n=\iint r^{n-1}u(rs)\sigma_{n-1}(\odif s)\lambda^1|_{]0,\infty[}(\odif r)\]
    \end{thm}
    

\chapter{Faltung und Fourier-Transformation}
\section{Die Faltung auf \texorpdfstring{$L^p_\mathbb{K}(\beta^r)$}{L\^p(beta\^r)}}
    \begin{thm}
        If $f,g:\mathbb{R}^r\rightarrow \mathbb{K}$ are Borel-measurable, then the function
        \[
            \varphi:\mathbb{R}^{2r}\rightarrow \mathbb{K},\,\,(x,y)\mapsto f(x-y)g(y)
        \]
        is Borel-measurable, and thus $\varphi(x,\bullet):\mathbb{R}^r\rightarrow \mathbb{K}$ is Borel-measurable.
    \end{thm}
    \begin{define}
        The \textit{Faltung} of two Borel-measurable functions $f,g:\mathbb{R}^r\rightarrow \mathbb{K}$ is the partial function
        \[
            f*g:\mathbb{R}^r\rightarrow \mathbb{K},\,\,x\mapsto \int f(x-y)g(y)\,\beta^r(dy)
        \]
    \end{define}
    \begin{thm}\label{cases when the Faltung is defined everywhere}
        Let $p,q\in [1,\infty]$ such that $p^{-1}+q^{-1}=1$. If $f\in L^p_\mathbb{R}(\beta^r)$ and $g\in L^q_\mathbb{R}(\beta^r)$, then $f*g$ is defined everywhere on $\mathbb{R}^r$. (what about $\mathbb{K}=\mathbb{C}$?)
    \end{thm}
    \begin{proof}
        By H\"{o}lder's inequality, $\|f(x-\bullet)g(\bullet)\|_1\leqslant \|f(x-\bullet)\|_p\cdot \|g\|_q=\|f\|_p\cdot \|g\|_q<\infty$.
    \end{proof}
    \begin{thm}\label{Faltung of two L^1 functions is L^1}
        For $f,g\in L^1_\mathbb{K}(\beta^p)$, the Faltung $f*g$ is a $\beta^p$-almost function, $f*g\in L^1_\mathbb{K}(\beta^p)$, and
        \[
            \|f*g\|_1\leqslant \|f\|_1\cdot \|g\|_1
        \]
    \end{thm}
    \begin{proof}
        By Corollary \ref{transformation theorem for affine automorphisms}, with $\tau(x)=x-y$,
        \[
            \int \left(\int |f(x-y)g(y)|\,\beta^p(dx)\right)\beta^p(dy)=\|f\|_1\int |g(y)|\,\beta^p(dy)=\|f\|_1\cdot \|g\|_1<\infty
        \]
        By Fubini's theorem, $f*g$ is a measurable $\beta^p$-almost function that is in $L^1_\mathbb{K}(\beta^p)$, and $\|f* g\|_1\leqslant \int \left(\int |f(x-y)g(y)|\,\beta^p(dy)\right)\beta^p(dx)=\int \left(\int |f(x-y)g(y)|\,\beta^p(dx)\right)\beta^p(dy)=\|f\|_1\cdot \|g\|_1$.
    \end{proof}
    \begin{thm}
        Let $f,g,h\in L^1_\mathbb{K}(\beta^p)$.
        \begin{tight_enumerate}
            \item (commutativity) $f*g=g*f$.
            \item (distributivity) $(f+g)*h=f*h+g*h=h*f+h*g=h*(f+g)$.
            \item (associativity) $(f*g)*h=f*(g*h)$.
        \end{tight_enumerate}
        Thus, $(L^1_\mathbb{K}(\beta^p),*)$ is a commutative Banach algebra.
    \end{thm}
\section{Faltung of measures}
    Let $\mu,\nu$ be $\sigma$-finite measures on $(\mathbb{R}^n,\mathscr{B}(\mathbb{R}^n))$. Let $f,g\in\mathcal{M}(\mathbb{R}^n;\overline{\mathbb{R}})$.
    \begin{enumerate}
        \itemsep0em
        \item $f*g:\mathbb{R}^n\rightarrow \overline{\mathbb{R}},\quad \displaystyle x\mapsto \int f(x-y)g(y)\lambda^n(\odif y)$.
        \item $f*\mu:\mathbb{R}^n\rightarrow \overline{\mathbb{R}},\quad \displaystyle x\mapsto \int f(x-y)\mu(\odif y)$.
        \item $\mu*\nu:\mathscr{B}(\mathbb{R}^n)\mapsto \overline{\mathbb{R}}_+,\quad B\mapsto \displaystyle \iint \mathbbm{1}_B(x+y)\mu(\odif x)\nu(\odif y)$.
    \end{enumerate}
    are called Faltung (resp. of two functions, of a function with a measure and of two measures).
    \begin{thm}
        \begin{enumerate}
            \itemsep0em
            \item The Faltung of two functions, of a function with a measure, or of two measures are bilinear.
            \item $\mu*\nu=S_*(\mu\times \nu)$, where $S:\mathbb{R}^n\times \mathbb{R}^n\rightarrow \mathbb{R}^n,\,\, (x,y)\mapsto x+y$.
            \item $\mu*\nu=\nu*\mu$.
            \item $\mu*\nu(B)=\displaystyle \int \mu(B-y)\nu(\odif y)=\int \nu(B-x)\mu(\odif x)$, where $\tau_y:x\mapsto x-y$.
            \item $\displaystyle \int f\odif\mu*\nu=\iint f(x+y)\mu(\odif y)\nu(\odif x)$.
            \item $f\in\mathcal{L}^1(\mu*\nu;\mathbb{R})\Longleftrightarrow f\circ S\in \mathcal{L}^1(\mu\times \nu;\mathbb{R})$.
            \item Let $f,g\geqslant 0$. Then $f\lambda^n*g\lambda^n=(f*g)\lambda^n$.
            \item $f*g=g*f$.
            \item $|f*g|\leqslant |f|*|g|$.
        \end{enumerate}
    \end{thm}
    \begin{thm}
        Let $f,g\in\mathcal{M}(\mathbb{R}^n;\overline{\mathbb{R}})$. What if $f*g$ is not almost defined? What assumptions can guarantee the almost definedness of $f*g$?
        %Let $f\in\mathcal{L}^1(\lambda^n;\mathbb{R})$ and $g\in\mathcal{L}^p(\lambda^n;\mathbb{R}),p\in[1,\infty)$. Then $f*g\in \mathcal{L}^p(\lambda^n;\mathbb{R})$ and
        \[\|f*g\|_p\leqslant \|f\|_1\cdot \|g\|_p\]
    \end{thm}
    % \begin{thm}\label{continuity of p-norm w.r.t. infinitesimal translation}
    %     Let $u\in L^p_\mathbb{R}(\beta^r),\,p\in[1,\infty[$. The map
    %     \[x\mapsto \int |u(x+y)-u(y)|^p\,\beta^r(\odif y)=
    %     \|u(x+\bullet)-u(\bullet)\|_p^p\]
    %     is continuous. In particular, as $x\to 0$,
    %     \[
    %         u(x+\bullet)\to u(\bullet)\text{ in }L^p_\mathbb{R}(\beta^r)
    %     \]
    % \end{thm}
\section{Strong continuity of the translation group}
    \begin{thm}\label{continuity of p-norm w.r.t. infinitesimal translation}
        Let $u\in L^p(\lambda^r),\,1\leqslant p<\infty$. The map
        \[\mathbb{R}^r\rightarrow L^p(\lambda^r),\,\, t\mapsto
        u(\ph+t)\]
        is uniformly continuous. That is, denote by $\tau_t$ the translation by $t$ in $\mathbb{R}^r$, the translation group $(\tau_t^*)_{t\in\mathbb{R}^r}$ on $L^p$ is strongly continuous.
    \end{thm}
    \begin{proof}
        Recall that $C_c(\mathbb{R}^r)$ is dense in $L^p(\lambda^r)$ for $1\leqslant p<\infty$. For $\phi\in C_c(\mathbb{R}^r)$, if $\supp \phi\subset B_R(0)$ then $|\phi(\ph+t)-\phi(\ph)|^p$ is majorized by $2^p\|\phi\|_\infty^p\one_{B_{R+1}(0)}$ for all $t\in B_1(0)$, so $x\mapsto \int |\phi(\ph+t)-\phi(\ph)|^p\odif\beta^r=\|\phi(\ph+t)-\phi(\ph)\|_p^p$ is continuous at $0$. Let $\phi_n\to u$ in $L^p(\lambda^r)$. $%|\|u(x+\bullet)-u(\bullet)\|_p-\|u(x'+\bullet)-u(\bullet)\|_p|\leqslant 
        \|u(\ph+t)-u(\ph+t')\|_p\leqslant \|u(\ph+t)-\phi_n(\ph+t)\|_p+\|\phi_n(\ph+t)-\phi_n(\ph+t')\|_p+\|\phi_n(\ph+t')-u(\ph+t')\|_p=\|\phi_n(\ph+t-t')-\phi_n(\ph)\|_p+2\|\phi_n-u\|_p$. Given $\varepsilon>0$, there exists $n(\varepsilon)\in\mathbb{N}$ such that $2\|\phi_{n(\varepsilon)}-u\|_p<\varepsilon/2$ and then for $t,t'$ sufficiently close $\|\phi_{n(\varepsilon)}(\ph+t-t')-\phi_{n(\varepsilon)}(\ph)\|_p<\varepsilon/2$.
    \end{proof}
    \begin{crl}\label{continuity and boundedness of the Faltung}
        $f\in L^1_\mathbb{R}(\beta^r)\wedge g\in L^\infty_\mathbb{R}(\beta^r)\Longrightarrow f*g\text{ is bounded and uniformly continuous}$.
    \end{crl}
    \begin{proof}
        By Theorem \ref{cases when the Faltung is defined everywhere}, $f*g$ is defined everywhere.
    \end{proof}
\section{Regularization}
    \begin{thm}\label{unit generalized L^1 functions}
        Let $(k_n)\in L^1_\mathbb{R}(\beta^p)^\mathbb{N}$ with $k_n\geqslant 0$, $\|k_n\|_1=1$ and $\diam(\{0\}\cup \,\supp k_n)\to 0$ as $n\to \infty$. Then for all $f\in L^1_\mathbb{R}(\beta^p)$,
        \[
            k_n*f{\rightarrow} f\text{ in }L^1_\mathbb{R}(\beta^p)
        \]
        where convergence is in the topology induced by the norm.
    \end{thm}
    \begin{proof}
        The proof relies on Theorem \ref{continuity of p-norm w.r.t. infinitesimal translation}.
    \end{proof}
    \begin{thm}
        For $f\in L^1_\mathbb{R}(\beta^r)$ and $g\in C_c^\infty(\mathbb{R}^r)$ holds that $f*g\in C^\infty(\mathbb{R}^r)$ and
        \[
            \pdv{(f*g)}{x_k}=f*\pdv{g}{x_k}
        \]
    \end{thm}
    \begin{proof}
        $C_c^\infty(\mathbb{R}^r)\subset C_c(\mathbb{R}^r)\subset L^\infty_\mathbb{R}(\beta^r)$, so $f*g$ is defined everywhere. Since $\pdv{g}{x_k}$ is uniformly continuous, $|((f*g)(x+te_k)-(f*g)(x))/t-(f*\pdv{g}{x_k})(x)|\to 0$ by finite increment theorem. Since $\pdv{g}{x_k}\in C_c^\infty(\mathbb{R}^r)$, $f*\pdv{g}{x_k}$ is (uniformly) continuous by Corollary \ref{continuity and boundedness of the Faltung}.
    \end{proof}
    \begin{crl}
        $C_c^\infty(\mathbb{R}^r)$ is dense in $L^1_\mathbb{R}(\beta^r)$.
    \end{crl}
    \begin{lem}
        There exists a function $\phi\in \mathcal{C}^\infty(\mathbb{R}^n;\mathbb{R})$ such that
        \[\phi\geqslant 0\quad\wedge\quad \supp\phi\subset\cl B(0,1)\quad\wedge\quad\int\phi\odif \lambda^n=1\]
        Set $\phi_\varepsilon:x\mapsto \varepsilon^{-n}\phi(x/\varepsilon)$ for $\varepsilon>0$. The function $J_\varepsilon(u):=\phi_{\varepsilon}*u$ is called the Friedrichs mollifier of $u\in\mathcal{L}^p(\mu;\mathbb{R}),p\in[1,\infty)$.
    \end{lem}
    \begin{lem}
        Let $u\in\mathcal{L}^p(\mu;\mathbb{R}),p\in[1,\infty)$.
        \begin{enumerate}
            \item $\phi_\varepsilon\in\mathcal{C}_c^\infty(\mathbb{R}^n;\mathbb{R})\quad\wedge\quad\supp\phi_\varepsilon\subset\cl B(0,\varepsilon)\quad\wedge\quad\|\phi_\varepsilon\|_1=1$.
            \item $\supp u*\phi_\varepsilon\subset \supp u+\supp \phi_\varepsilon\subset \{y\in\mathbb{R}^n\,|\,\exists x\in\supp u(|x-y|\leqslant\varepsilon)\}$.
            \item $\|\phi_\varepsilon*u\|_p\leqslant \|u\|_p$ for all $\varepsilon>0$.
        \end{enumerate}
    \end{lem}
    \begin{thm}
        Let $u\in L^p_\mathbb{R}(\beta^r),p\in[1,\infty[$. Then $\phi_\varepsilon*u\in C^\infty(\mathbb{R}^r)$ and
        \[[\phi_\varepsilon*u]\to[u],\quad\text{as }\varepsilon\to 0\]
    \end{thm}
\section{The Fourier transform}
    Let $r\in\mathbb{N}$. Throughout this section, $L^p$ is short for $L^p(\mathbb{R}^r,\mathscr{B}(\mathbb{R}^r),\beta^r)$, $1\leqslant p\leqslant \infty$.
\subsection{The Fourier transform of integrable functions}
    \begin{define}
        The \textit{Fourier transform} of $u\in L^1$ is the function
    \[
        \mathcal{F}u%=\widehat{u}
        :\mathbb{R}^r\rightarrow \mathbb{C},\,\, \xi\mapsto \int_{\beta^r} e^{-i2\pi\braket{\xi}{x}}u(x)\odif x
    \]
    \end{define}
    \begin{define}
        A function $f:\mathbb{R}^r\rightarrow\mathbb{C}$ is said to \textit{vanish at infinity} if 
        \[
            f(x)\to 0\text{ as }\|x\|\to\infty
        \]
        i.e., for any $\varepsilon>0$ there exists $R>0$ such that $|f(x)|<\varepsilon$ for all $\|x\|>R$.

        The set of continuous functions $\mathbb{R}^r\rightarrow\mathbb{C}$ vanishing at infinity is denoted by $C_0$.
    \end{define}
    \begin{thm}\label{Fourier transform of a L^1 function is bounded by its 1-norm}
        The Fourier transformation is a short % $\mathbb{C}$-
        linear operator $\mathscr{F}:L^1\rightarrow C_0\subset L^\infty$.
        % with
        % \[
        %     \|\mathscr{F}u\|_\infty\leqslant \|u\|_1
        % \]
    \end{thm}
    \begin{proof}
        Let $u\in L^1$. For all $\xi\in\mathbb{R}^r$, $|\mathcal{F}u(\xi)|\leqslant \int |e^{-i2\pi\braket{\xi}{x}}u(x)|\odif x= \|u\|_1$. Ergo $\|\mathcal{F}u\|_\infty\leqslant \|u\|_1$. To show $\mathcal{F}u:\mathbb{R}^n\rightarrow \mathbb{C}$ is continuous, apply Theorem \ref{continuity of parameter-dependent integral}. It remains to show\footnote{This is the Riemann-Lebesgue lemma.} $\lim_{\|\xi\|\to\infty} \mathcal{F}u(\xi)=0$. $\mathcal{F}f=-\int e^{-i\left<t,x\right>}f(x+\frac{\pi}{\|t\|^2}t)\odif\rho^r$. Thus, $2|\mathcal{F}f(t)|\leqslant \int |f(x)-f(x+\frac{\pi}{\|t\|^2}t)|\odif\rho^r\to 0$ as $\|t\|\to\infty$ by Theorem \ref{continuity of p-norm w.r.t. infinitesimal translation}.
    \end{proof}
\subsection{The convolution theorem}
    \begin{thm}
        Let $u,v\in L^1$.
        \[
            \mathscr{F}(u*v)=\mathscr{F}u\cdot \mathscr{F}v
        \]
    \end{thm}
    \begin{proof}
        By Theorem \ref{Faltung of two L^1 functions is L^1}, $\|e^{-i2\pi\braket{\xi}{\bullet}}(u*v)\|_1<\infty$. Use Fubini's theorem.
    \end{proof}
    Similar results hold for convolution of measures (see Schilling).
    \begin{crl}
        There exists no $k\in L^1$ such that $k*f=f$ for all $f\in L^1$.
    \end{crl}
\subsection{Scaling and shifting properties}
    \begin{thm}
        Let $u\in L^1$. Let $\tau_a:\mathbb{R}^r\rightarrow \mathbb{R}^r,\,\,x\mapsto x+a$ be the translation by $a\in\mathbb{R}^r$. %and $\gamma_\alpha:\mathbb{R}^r\rightarrow \mathbb{R}^r,\,\,x\mapsto \alpha x$ the scaling by $\alpha \in\mathbb{R}\setminus\{0\}$.
        \begin{tight_enumerate}
            \item $\mathscr{F}(u(\ph+a))=e^{i2\pi\braket{\sph}{a}}\cdot \mathscr{F}u$.
            \item $\mathscr{F}(e^{i2\pi\braket{a}{\sph}}\cdot u)=(\mathscr{F}u)(\ph-a)$.
            \item $\widehat{u\circ\alpha}=|\alpha|^{-r}(\widehat{u}\circ \alpha^{-1})$. % To make sense of $\alpha^*f$ recall that a scalar for a vector space is identified via a ring homomorphism with an endomorphism of the underlying abelian group.
            \item $\mathcal{F}(u\circ R)=\mathcal{F}u\circ R$, where $R$ is an orthogonal linear tranformation.
        \end{tight_enumerate}
    \end{thm}
\subsection{Injectivity}
    \begin{thm}\label{Fourier inversion transform}
        % The Fourier transformation $\mathcal{F}:W^r_\mathbb{C}\rightarrow L^1_\mathbb{C}(\rho^r)$ has a left inverse:
        % \[
        %     \mathcal{G}:L^1_\mathbb{C}(\rho^r)\rightarrow L^\infty_\mathbb{C}(\rho^r),\,\,f\mapsto \left(t\mapsto % \mathcal{F}(f(-x))(t)=
        %     \int e^{i\left<t,\bullet\right>}f(\bullet)\odif\rho^r\right)
        % \]
        Let $u\in L^1$. If $\mathcal{F}u\in L^1$, then for almost all $x\in \mathbb{R}^r$,
        \[
            u(x)=\int_{\beta^r} e^{i2\pi\braket{x}{\xi}}\mathcal{F}u(\xi)\odif\xi
        \]
    \end{thm}
    \begin{proof}
        Let $k_n:x\mapsto (2\pi)^{r/2}\prod_{i=1}^r 0\vee (n-n^2x_i)$. For $f\in L^1_\mathbb{C}(\rho^r)$, let $\mathcal{G}f:t\mapsto (\mathcal{F}f)(-t)=\mathcal{F}(f(-\bullet))t=\int e^{i\langle t,\bullet\rangle}f(\bullet)\odif\rho^r$. Verify that $\mathcal{F}k_n(t)=\prod_{i=1}^r(\frac{\sin (t_i/2n)}{t_i/2n})^2$ and $\mathcal{G}\mathcal{F}k_n=k_n$ for all $n\in\mathbb{N}$. Since $\mathcal{F}k_n\cdot \mathcal{F}f=\mathcal{F}(k_n*f)\in L^1_\mathbb{C}(\rho^r)$, by Fubini's theorem, $\mathcal{G}(\mathcal{F}k_n\cdot \mathcal{F}f)=f*k_n$. By dominated convergence ($|\mathcal{F}k_n\cdot \mathcal{F}f|\leqslant |\mathcal{F}f|$ since $\sin x/x<1$), $\mathcal{F}k_n\cdot \mathcal{F}f\to \mathcal{F}f$ in $L^1_\mathbb{C}(\rho^r)$ and thus $f*k_n=\mathcal{G}(\mathcal{F}k_n\cdot \mathcal{F}f)$ converges uniformly to $\mathcal{G}(\mathcal{F}f)$ since $|\mathcal{G}(\mathcal{F}k_n\cdot \mathcal{F}f-\mathcal{F}f)|\leqslant \|\mathcal{F}k_n\cdot \mathcal{F}f-\mathcal{F}f\|_1\to 0$. Thus, for $R>0$, $\|(\mathcal{G}\mathcal{F}f-f)\cdot \one_{B_R(0)}\|_1\leqslant \|(\mathcal{G}\mathcal{F}f-f*k_n)\cdot \one_{B_R(0)}\|_1+\|(f*k_n-f)\cdot \one_{B_R(0)}\|_1\to 0$ (
        Recall Theorem \ref{unit generalized L^1 functions}: $f*k_n\to f$ in $L^1_\mathbb{C}(\rho^r)$). Thus, $\mathcal{G}\mathcal{F}f=f$ a.e. on $B_R(0)$. But $R$ is arbitrary.
    \end{proof}
    \begin{thm}
        The Fourier transformation $\mathcal{F}:L^1\rightarrow C_0$ is injective.
    \end{thm}
    \begin{proof}
        If $\mathcal{F}f=\mathcal{F}g$, then $\mathcal{F}(f-g)=0$ and $(f-g)(x)=\int_{\beta^r} 0=0$ for a.a. $x\in\mathbb{R}^r$, i.e., $f=g$ a.e.
    \end{proof}
\subsection{The Fourier algebra}
    \begin{define}
        The \textit{Fourier algebra} is the image of the Fourier transformation $L^1\rightarrow C_0$.
    \end{define}
    \begin{thm}
        The Fourier algebra with pointwise multiplication is indeed a $\mathbb{C}$-algebra.
    \end{thm}
    \begin{thm}\label{density of Fourier algebra in C_0}
        The Fourier algebra is dense, uniformly, in $C_0$.
    \end{thm}
    \begin{proof}
        Apply Stone-Weierstrass theorem.
    \end{proof}
\subsection{Fourier transform in the Schwartz space}
    \begin{thm}
        $S(\mathbb{R}^r)$ with pointwise multiplication is a $\mathbb{R}$-algebra.
    \end{thm}
    \begin{thm}
        If $f$ is a Schwartz function (i.e., rapidly decreasing), then for all $x\in \mathbb{R}^r$,
        \[
            f(x)=\int e^{i\langle x,\bullet\rangle}\mathcal{F}f(\bullet)\odif\rho^r
        \]
    \end{thm}
\subsubsection{Differential properties}
    \begin{thm}
        Let $f(x)\in L^1_\mathbb{C}(\rho^r)\cap C^\infty$ and assume $x^k\cdot f\in L^1$, where $x=(x^1,\dots,x^r)$.
        \[
            \pdv{}{t_k}\mathcal{F}f=(-i)\mathcal{F}(x^k\cdot f)
        \]
    \end{thm}
% \subsection{The Wiener algebra, see Schilling, the term Wiener algebra is not quite standard}
%     \begin{define}
%         The \textit{Wiener algebra} is 
%         \[W:=%\mathcal{F}^{-1}(L^1_\mathbb{C}(\rho^r))=
%         \{f\in L^1\,|\,\mathcal{F}f\in L^1\}\]
%     \end{define}
%     \begin{thm}
%         The Fourier transformation $\mathcal{F}$ restricts to a $\mathbb{C}$\cat{-Vect} automorphism of $W$.
%     \end{thm}
%     \begin{proof}
%         For $w\in W$, $\mathcal{F}w$ and $u:t\mapsto \mathcal{F}w(-t)$ are $L^1$. By Theorem \ref{Fourier inversion transform} $\mathcal{F}u=w$, and in particular $u\in W$. Thus, $W\subset \mathcal{F}(W)$. Moreover, $\mathcal{F}\mathcal{F}w=w(-\bullet)$. Thus, $\mathcal{F}w\in W$ and $\mathcal{F}(W)\subset W$.
%     \end{proof}
%     \begin{thm}
%         Let $u,v\in L^1$.
%         \begin{tight_enumerate}
%             \item $W\subset C_0$.
%             \item $W\subset \bigcap_{1\leqslant p<\infty}L^p$.
%             \item $u,v\in W\Longrightarrow u*v\in W$.
%             \item $u,v\in W\Longrightarrow u\cdot v\in W$ and
%             \[
%                 \mathcal{F}(u\cdot v)=\mathcal{F}u*\mathcal{F}v
%             \]
%         \end{tight_enumerate}
%     \end{thm}
%     \begin{proof}
%         $\mathcal{F}^{-1}(u*v)=\mathcal{F}^{-1}u\cdot \mathcal{F}^{-1}v$ since $\mathcal{F}^{-1}f=\widetilde{\mathcal{F}f}$, where $\widetilde{f}:t\mapsto f(-t)$.
%     \end{proof}
%     \begin{lem}
%         $u\in L^1\Longrightarrow u*g_t\in W$.
%     \end{lem}
%     \begin{lem}
%         If $u\in C_b$ is uniformly continuous, then
%         \[\lim_{t\to 0} [u*g_t]=[u]\]
%         where the convergence is taken in $L^\infty$.
%     \end{lem}
%     \begin{lem}
%         $W$ is dense in $L^p$ for $p\in[1,\infty)$.
%     \end{lem}
%     \begin{lem}
%         $W$ is dense in $C_0$.
%     \end{lem}
\subsection{Plancherel's theorem}
    \begin{thm}
        Let $f\in L^1_{\mathbb{C}}(\rho^r)\cap L^2_\mathbb{C}(\rho^r)$.
        \[\|\mathcal{F}f\|_2=\|f\|_2\]
        In particular, $\mathcal{F}:L^1_{\mathbb{C}}(\rho^r)\cap L^2_\mathbb{C}(\rho^r)\rightarrow L^2_\mathbb{C}(\rho^r)$ and there is a continuous extension \[\overline{\mathcal{F}}:L^2_{\mathbb{C}}(\rho^r)\rightarrow L^2_{\mathbb{C}}(\rho^r)\]
    \end{thm}
    \begin{proof}
        $W^r_\mathbb{C}\subset L^1_{\mathbb{C}}(\rho^r)\cap L^2_\mathbb{C}(\rho^r) \subset L^2_\mathbb{C}(\rho^r)$ and $W^r_\mathbb{C}$ is dense in $L^2_\mathbb{C}(\rho^r)$.
    \end{proof}
    The Fourier transform of $f\in L^2_\mathbb{C}(\rho^r)$ can be computed by means of an improper integral (see https://math.stackexchange.com/questions/2857971/compute-the-fourier-transform-of-a-l2-function).
\subsection{Parseval's formula}
    \begin{thm}
        $L^2_\mathbb{C}(\mu)$ is a Hilbert space with inner product given by
        \[
            \left<f,g\right>:=\int f\cdot \overline{g}\odif\mu
        \]
    \end{thm}
    \begin{thm}
        The Fourier transformation is a unitary operator on $L^2_\mathbb{C}(\rho^r)$: for $f,g\in L^2_\mathbb{C}(\rho^r)% L^1_{\mathbb{C}}(\rho^r)\cap L^2_\mathbb{C}(\rho^r)
        $.
        \[
            \left<\mathcal{F}f,\mathcal{F}g\right>=\left<f,g\right>
        \]
    \end{thm}
    \begin{proof}
        $\left<f,g\right>%=(\|f+g\|_2^2-\|f\|_2^2-\|g\|_2^2)/2
        =(\|f+g\|_2^2-\|f-g\|_2^2)/4$.
    \end{proof}
\subsection{The Fourier transform of finite measures}
    \begin{define}
        The Fourier transform of a finite measure $\mu$ on $(\mathbb{R}^r,\mathscr{B}(\mathbb{R}^r))$ is
        \[
            \widehat{\mu}:\mathbb{R}^r\rightarrow \mathbb{C},\,\,\xi\mapsto \int_\mu e^{i\braket{x}{\xi}}\odif x
        \]
    \end{define}
    \begin{thm}
        Let $u\in L^1_{\mathbb{R}}(\lambda^r)$. Then 
        \[
        \mathcal{F}u(\xi)=(\widehat{u^+\odot \lambda^r}-\widehat{u^-\odot \lambda^r})(-2\pi\xi)
        \]
    \end{thm}
    \begin{lem}
        Let $t\in\mathbb{R}^*_+$.
        \[g_t:x\mapsto (2\pi t)^{-n/2}e^{-|x|^2/2t}\Longrightarrow \widehat{g_t}:\xi\mapsto (2\pi)^{-n}e^{-t|\xi|^2/2}\]
    \end{lem}
    \begin{thm}
        The Fourier transform $\widehat{\mu}$ is a continuous function bounded by
        \[
            |\widehat{\mu}|\leqslant \mu(\mathbb{R}^r)
        \]
    \end{thm}
    \begin{lem}
        Let $u\in L^1_{\mathbb{R}}(\lambda^n)$ such that $u(x)=f(|x|)$.
        \[\mathcal{F}\{f\circ |\cdot|\}(\xi)=\frac{(2\pi)^{n/2}}{|\xi|^{n/2-1}}\int_{0}^\infty f(r)r^{n/2}J_{n/2-1}(r|\xi|)\odif r\]
        where $J_\nu$ is a Bessel function of the first kind of order $\nu$.
    \end{lem}
\subsection{Injectivity}
    \begin{thm}
        For $\mu,\nu$ two finite measures on $(\mathbb{R}^r,\mathscr{B}(\mathbb{R}^r))$,
        \[
            \widehat{\mu}=\widehat{\nu}\Longrightarrow \mu=\nu
        \]
    \end{thm}
    \begin{proof}
        For $f$ in the Fourier algebra, there exists $u\in L^1$ with $f=\mathcal{F}u$. Thus, $\int_\mu f=\int_\mu \mathcal{F}u=\int_\mu\int_{\beta^r}e^{-i2\pi\braket{\xi}{x}}u(x)\odif x\odif\xi=\int_{\beta^r}u(x)\int_\mu e^{-i2\pi\braket{\xi}{x}}\odif\xi\odif x=\int_{\beta^r}u(x)\widehat{\mu}(-2\pi x)\odif x$. Similarly, $\int_\nu f=\int_{\beta^r}u(x)\widehat{\nu}(-2\pi x)\odif x$. Thus, $\int_\mu f=\int_\nu f$. But the Fourier algebra is uniformly dense in $C_0$. Thus, $\int_\mu$ and $\int_\nu$ agree on $C_0$. It remains to apply Riesz's theorem.
    \end{proof}
\subsection{Smoothness}
    \begin{thm}
        Let $\mu$ be a finite measure on $(\mathbb{R}^r,\mathscr{B}(\mathbb{R}^r))$. Let $1\leqslant p\leqslant \infty$.
        
        If $x^1,\dots,x^r\in L^p(\mu)$, then $\widehat{\mu}\in C^p$, and the partial derivative
        \[
            \frac{\partial^p}{\partial \xi^{k_1}\cdots\partial \xi^{k_p}}\widehat{\mu}:\xi\mapsto i^p\int_\mu x_{k_1}\cdots x_{k_p}e^{i\braket{\xi}{x}}\odif x
        \]
        is in $C_b$ for all $k_1,\dots,k_p\in\{1,\dots,r\}$.
    \end{thm}
    \begin{proof}
        Assume $p<\infty$. Let $\alpha$ be an $r$-dimensional multi-index with $\alpha_1+\cdots+\alpha_r=q\leqslant p$. Then by Theorem \ref{AM-GM}, $|(x^1)^{\alpha_1}\cdots (x^r)^{\alpha_r}|=|(x^1)^q|^{\frac{\alpha_1}{q}}\cdots |(x^r)^q|^{\frac{\alpha_r}{q}}\leqslant \frac{\alpha_1}{q}|(x^1)^q|+\cdots+\frac{\alpha_r}{q}|(x^r)^q|\in L^1(\mu)$ since $x^1,\dots,x^r\in L^p(\mu)\subset L^q(\mu)$. Thus, $(x^1)^{\alpha_1}\cdots (x^r)^{\alpha_r}e^{i\braket{\xi}{\bullet}}\in L^1(\mu)$. By Theorem \ref{differentiability of parameter-dependent integral}, the partial derivatives of $\widehat{\mu}$ exist up to order $p$, and by Theorem \ref{continuity of parameter-dependent integral}, are continuous.

        If $p=\infty$, then $(x^1)^{\alpha_1}\cdots (x^r)^{\alpha_r}\in L^1(\mu)$ is immediate, for $\alpha$ any $r$-dimensional multi-index.
    \end{proof}
    It is claimed that for $p$ even, in the preceding theorem, if the partial derivatives up to order $p$ exist at $\xi =0$, then $x^1,\dots,x^r\in L^p(\mu)$ and hence $\widehat{\mu}\in C^p$. How?
    \begin{thm}
        Let $u\in L^1$. If $(x^1)^p\cdot u,\cdots,(x^r)^p\cdot u\in L^1$, then $\mathcal{F}u\in C^p$, and
        \[
            \frac{\partial^p}{\partial \xi^{k_1}\cdots\partial \xi^{k_p}}\mathcal{F}u=(-2\pi i)^p\mathcal{F}(x_{k_1}\cdots x_{k_p}\cdot u)
        \]
    \end{thm}
\section{Weak Convergence}
    Be $(\mu_n)_{n\in\mathbb{R}}$ and $\mu$ finite measures on $(\mathbb{R}^n,\mathscr{B}(\mathbb{R}^n))$. Then $\mu_n$ converges weakly to $\mu$ if
    \[\lim_{n\to\infty}\int u\odif{\mu_n}=\int u\odif{\mu}\quad\forall u\in \mathcal{C}_b(\mathbb{R}^n;\mathbb{R})\] 
    \begin{thm}
        \[\mathcal{F}\{\mu_n\}\xrightarrow[n\to\infty]{\text{pointwise}}\mathcal{F}\{\mu\}\Longleftrightarrow \mu_n\xrightarrow[n\to\infty]{\text{weakly}}\mu\]
    \end{thm}
\section{The Laplace transform}
\subsection{of functions}
    The Laplace transform of a function $f\in L^1(\lambda^n|_{\mathbb{R}_+})$ is
    \[
        \mathcal{L}\{f\}:=\mathcal{L}_t\{f(t)\}:\mathbb{C}\rightarrow\mathbb{C},\quad s\mapsto \int_{\mathbb{R}_+} e^{-st}f(t)\odif t
    \]
    \begin{exam}
        \begin{tight_enumerate}
            \item $\mathcal{L}\{1\}(s)=\dfrac{1}{s}$.
            \item $\mathcal{L}_t\{e^{\alpha t}\}(s)=\dfrac{1}{s-\alpha}$.
            \item $\mathcal{L}_t\{\cos\omega t\}(s)=\dfrac{s}{s^2+\omega^2}$.
            \item $\mathcal{L}_t\{\sin\omega t\}(s)=\dfrac{\omega}{s^2+\omega^2}$.
        \end{tight_enumerate}
    \end{exam}
    \begin{lem}
        If $f$ is of exponential type, i.e. as $t\to\infty$, the inequality
        \[
            |f(t)|\leqslant Ae^{kt}
        \]
        holds ultimately, where $A,k\in\mathbb{R}$, then the Laplace transform $\mathcal{L}\{f\}$ exists.
    \end{lem}
    \begin{lem}
        If $f'\in L^1(\lambda^n|_{\mathbb{R}_+})$ and $f$ is of exponential type, then $\mathcal{L}\{f'\}$ exists and
        \[
            \mathcal{L}\{f'\}=s\mathcal{L}\{f\}-f(0)
        \]
    \end{lem}
    Under appropriate conditions, we have
    \[
        \mathcal{L}\{f^{(n)}\}=s^n\mathcal{L}\{f\}-\sum_{i=1}^ns^{n-i}f^{(i-1)}(0)
    \]
    \begin{lem}
        \[
            \mathcal{L}_t\{f(t-a)u(t-a)\}(s)=e^{-as}\mathcal{L}\{f\}(s)
        \]
    \end{lem}
    \[
        \mathcal{L}^{-1}\{e^{-\tau s}F(s)\}=[\mathcal{L}^{-1}\{F(s)\}(t-\tau)]u(t-\tau)
    \]
    \begin{align*}
        tu(t-2)&=(t-2)u(t-2)+2u(t-2)\\
        \mathcal{L}\{tu(t-2)\}&=\mathcal{L}\{(t-2)u(t-2)+2u(t-2)\}\\
        &=\mathcal{L}\{(t-2)u(t-2)\}+2\mathcal{L}\{1\cdot u(t-2)\}\\
        &=e^{-2s}\frac{1}{s^2}+2e^{-2s}\frac{1}{s}\\
        \mathcal{L}\{tu(t-2)\}&=e^{-2s}\mathcal{L}\{t+2\}\\
        &=e^{-2s}\left(\frac{1}{s^2}+\frac{2}{s}\right)\\
        \mathcal{L}\{y'\}&=s\mathcal{L}\{y\}-y(0)=s\mathcal{L}\{y\}-4\\
        \mathcal{L}\{y\}&=e^{-2s}\left(\frac{1}{s^3}+\frac{2}{s^2}\right)+\frac{4}{s}\\
        y&=\mathcal{L}^{-1}\left\{e^{-2s}\left(\frac{1}{s^3}+\frac{2}{s^2}\right)\right\}+\mathcal{L}^{-1}\left\{\frac{4}{s}\right\}\\
        y&=\left(\frac{1}{2}(t-2)^2+2(t-2)\right)u(t-2)+4\\
        y&=\left(\frac{1}{2}t^2-2\right)u(t-2)+4
    \end{align*}
\subsection{The Laplace transform of finite measures}
    \begin{define}
        The \textit{Laplace transform} of a finite measure $\mu$ on $([0,\infty]^r,\mathscr{B}([0,\infty]^r))$ is
        \[
            \mathcal{L}\mu:\mathbb{R}_+^r\rightarrow \mathbb{R}_+,\,\,\xi\mapsto \int_\mu e^{-\braket{\xi}{x}}\odif x
        \]
    \end{define}
    \begin{thm}
        If $\mu,\nu$ are finite measures on $(\mathbb{R}^r_+,\mathscr{B}(\mathbb{R}^r_+))$, then 
        \[
            \mathcal{L}\mu=\mathcal{L}\nu\Longrightarrow \mu=\nu
        \]
    \end{thm}
    \begin{proof}
        Apply Stone-Weierstrass theorem to show $\int_\mu u=\int_v u$ for $u\in C_0$. Use Riesz's theorem.
    \end{proof}
    \begin{thm}
        Of $\mu$ a finite measure on $([0,\infty],\mathscr{B}([0,\infty]))$, the Laplace transform $\mathcal{L}\mu:\mathbb{R}_+\rightarrow \mathbb{R}_+$ is a decreasing continuous function with
        \[
            \mathcal{L}u(0)=\mu([0,\infty]),\quad\lim_{\xi\to 0}\mathcal{L}\mu(\xi)=\mu([0,\infty[)
        \]
        % By Fubini's theorem, 
        % \[
        %     \mathcal{L}\mu(\xi)=\int_\xi^\infty \odif q \int_\mu xe^{-qx} \odif x
        % \]

        If $\id \in L^p(\mu)$, then $\mathcal{L}\mu\in C^p(]0,\infty[,\mathbb{R}_+)$ with
        \[
            \mathcal{L}\mu^{(p)}(\xi)=\int (-x)^p e^{-\xi x}\odif x
        \]
        and
        \[
            \lim_{\xi\downarrow 0}\mathcal{L}\mu^{(p)}(\xi)=(-1)^p\int_\mu x^p \odif x
        \]
    \end{thm}
    \begin{proof}
        $e^{-\xi x}\uparrow \mathbbm{1}_{[0,\infty[}(x)$ as $\xi\downarrow 0$ for every $x\in [0,\infty]$.

        By mean value theorem or convexity, $|(e^{-\xi x}-e^{-\xi_0 x})/(\xi-\xi_0)|\leqslant x $ for all $\xi_0,\xi,x>0$. Apply Theorem \ref{differentiability of parameter-dependent integral}.
    \end{proof}
\section{Application to signals \& systems}
\subsection{Energy signals, power signals and transients}
    \begin{define}
        An \textit{energy signal} is a function in $L^2_\mathbb{C}(\mathbb{R})$.
    \end{define}
    \begin{define}
        

        A \textit{power signal} is a function $u\in L_{\text{loc},\mathbb{C}}^2(\mathbb{R})$ for which the limit
        \[
            \|u\|_P^2:=\lim_{\tau\to\infty}\frac{1}{2\tau}\int_{-\tau}^\tau |u|^2\odif x
        \]
        exists, and is called the \textit{average power} of $u$. A \textit{transient} is a power signal with zero average power.
    \end{define}
    \begin{thm}
        Power signals form a vector subspace of $L_{\text{loc},\mathbb{C}}^2(\mathbb{R})$. 
        
        Transients form a vector subspace of the space of power signals.
    \end{thm}
    Let $P$ be the vector space of power signals modulo the subspace of transients.
    \begin{thm}
        $P$ is an inner product space with
        \[
            \braket{u}{v}_P:=\lim_{T\to\infty}\frac{1}{2T}\int_{-T}^T \overline{u}v\odif x
        \]
    \end{thm}
    \begin{thm}
        In $P$, the complex exponentials $(e^{i\omega t})_{\omega\in\mathbb{R}}$ are orthonormal.
    \end{thm}
    \begin{thm}
        Every periodic $u\in L_{\text{loc},\mathbb{C}}^2(\mathbb{R})$ function is a power signal with
        \[
            \|u\|^2=\frac{1}{T}\int_{0}^T |u|^2\odif x
        \]
        where $T$ is the period of $u$.
    \end{thm}
    \begin{define}
        The \textit{DC value} of a power signal $u\in P$ is $\braket{1}{u}_{P}$.
    \end{define}
    Define the following functions $\mathbb{R}\rightarrow \mathbb{R}$
    \begin{align*}
        \operatorname{u}&:=\mathbbm{1}_{]0,\infty[}\\
        \operatorname{rect}&:=\mathbbm{1}_{]-\frac{1}{2},\frac{1}{2}[}\\
        \operatorname{tri}&:=\operatorname{rect}*\operatorname{rect}=(1-|t|)\cdot \mathbbm{1}_{]-1,1[}
        % ,\,\,t\mapsto\begin{cases}
        %     t+1 &\text{for }-1<t<0\\
        %     1-t &\text{for }0<t<1\\
        %     0 &\text{otherwise}
        % \end{cases}
    \end{align*}
    fow which 
    \[
        \int \operatorname{rect} \odif\lambda=1=\int \operatorname{tri}\odif\lambda
    \]
    Let $\operatorname{sinc}:\mathbb{R}\rightarrow \mathbb{R}$ be the continuous function $(\sin \pi x)/\pi x$ and $\omega=2\pi f$ a change of variable.
    \begin{thm}
        \begin{align*}
            \widehat{\operatorname{rect}}(\omega)&=\frac{1}{\sqrt{2\pi}}\operatorname{sinc}\left(\frac{\omega}{2\pi}\right)=\frac{1}{\sqrt{2\pi}}\operatorname{sinc}f\\
            \widehat{\operatorname{tri}}(\omega)&=(\widehat{\operatorname{rect}}*\widehat{\operatorname{rect}})(\omega)=\frac{1}{2\pi}\operatorname{sinc}^2f
        \end{align*}
    \end{thm}
\subsection{Fourier analysis}
    So
    \[
        \widehat{\widehat{u}}=\widetilde{u},\quad \mathcal{G}u=\widehat{\widetilde{u}}=\widetilde{\widehat{u}}
    \]
    \begin{thm}
        $\overline{\mathcal{F}u}(\omega)=\mathcal{F}\overline{u}(-\omega)=\mathcal{G}\overline{u}$.

        In particular, $|\mathcal{F}u(\omega)|=|\mathcal{F}\overline{u}(-\omega)|$ and $-\arg {\mathcal{F}u}(\omega)=\arg \mathcal{F}\overline{u}(-\omega)$.

        \[
            \overline{\widehat{u}}=\check{\overline{u}}
        \]
        If $u$ is real, then
        \[
            \overline{\widehat{u}}=\check{u}=\widetilde{\widehat{u}}
        \]
        and clearly $|\widehat{u}|$ is even and $\operatorname{arg} \widehat{u}$ is odd.
        
        If $u$ is even, then $\widehat{u}$ is even:
        \[
            \widetilde{\widehat{u}}=\check{u}=\widehat{\widetilde{u}}=\widehat{u}
        \]
        If $u$ is odd, then $\widehat{u}$ is odd:
        \[
            \widetilde{\widehat{u}}=\check{u}=\widehat{\widetilde{u}}=-\widehat{u}
        \]
        Thus, the Fourier transformation preserves parity.

        If $u$ is real and even, then $\widehat{u}$ is real and even:
        \[
            \overline{\widehat{u}}=\check{u}=\widehat{u}
        \]
        If $u$ is real and odd, then $\widehat{u}$ is imaginary and odd:
        \[
            \overline{\widehat{u}}=\check{u}=-\widehat{u}
        \]
    \end{thm}
\subsection{Linear time-invariant systems}

\chapter{Martingales}
\section{sigma-finite filtrations}
    \begin{define}
        A filtered measure space $(X,(\mathscr{X}_t)_{t\in T},\mu)$ is \textit{$\sigma$-finitely filtered} if $\mu|_{\mathscr{X}_t}$ is $\sigma$-finite for all $t\in T$.
    \end{define}
    Clearly, if the toset $T$ admits a least element $t_\text{min}$, then $(\mathscr{X}_t)_{t\in T}$ is a $\sigma$-finite filtration if and only if $\mathscr{X}_{t_{\text{min}}}$ is $\sigma$-finite.
\section{Martingales}
    \begin{define}
        A stochastic process 
        \[
            u:T\times X\rightarrow \mathbb{R}
        \]
        on a filtered measure space $(X,(\mathscr{X}_t)_{t\in T},\mu)$ such that $u_t\in L^1(\mu|_{\mathscr{X}_t})$ for all $t\in T$ is called
        \begin{tight_enumerate}
            \item a \textit{submartingale} if for all $t_1,t_2\in T$, $t_1\leqslant t_2$ and any $A\in\mathscr{X}_{t_1}$
            \[
                \int_{\mu|_{\mathscr{X}_{t_2}}} u_{t_2}\mathbbm{1}_A\geqslant\int_{\mu|_{\mathscr{X}_{t_1}}} u_{t_1}\mathbbm{1}_A
            \]
            \item a \textit{supermartingale} if for all $t_1,t_2\in T$, $t_1\leqslant t_2$ and any $A\in\mathscr{X}_{t_1}$
            \[
                \int_{\mu|_{\mathscr{X}_{t_2}}} u_{t_2}\mathbbm{1}_A\leqslant\int_{\mu|_{\mathscr{X}_{t_1}}} u_{t_1}\mathbbm{1}_A
            \]
            \item a \textit{martingale} if it is both a submartingale and a supermartingale, i.e., for all $t_1,t_2\in T$, $t_1\leqslant t_2$ and any $A\in\mathscr{X}_{t_1}$
            \[
                \int_{\mu|_{\mathscr{X}_{t_2}}} u_{t_2}\mathbbm{1}_A=\int_{\mu|_{\mathscr{X}_{t_1}}} u_{t_1}\mathbbm{1}_A
            \]
        \end{tight_enumerate}
    \end{define}
    \begin{thm}
        On a filtered measure space $(X,(\mathscr{X}_t)_{t\in T},\mu)$,
        \begin{tight_enumerate}
            \item if $u:T\times X\rightarrow \mathbb{R}$ is a submartingale, then $-u$ is a supermartingale.
            \item if $u:T\times X\rightarrow \mathbb{R}$ is a supermartingale, then $-u$ is a submartingale.
            \item if $u_1,u_2:T\times X\rightarrow \mathbb{R}$ are martingales, so is $\alpha u_1+\beta u_2$ for all $\alpha,\beta\in \mathbb{R}$.
            \item if $u_1,u_2:T\times X\rightarrow \mathbb{R}$ are submartingales, so is $\alpha u_1+\beta u_2$ for $0\leqslant \alpha,\beta\in \mathbb{R}$.
            \item if $u_1,u_2:T\times X\rightarrow \mathbb{R}$ are supermartingales, so is $\alpha u_1+\beta u_2$ for $0\leqslant \alpha,\beta\in \mathbb{R}$.
            \item if $u:T\times X\rightarrow \mathbb{R}$ is a submartingale, then $u^+:T\times X\rightarrow\mathbb{R}$ is a submartingale.
            \item if $u:T\times X\rightarrow \mathbb{R}$ is a martingale, then $|u|$ is a submartingale.
            \item if $u:T\times X\rightarrow \mathbb{R}$ is a martingale and $u_t\in L^p(\mu|_{\mathscr{X}_t})$ for all $t\in T$ for some $1<p<\infty$, then $|u|^p$ is a submartingale.
        \end{tight_enumerate}
        \begin{proof}
            6. $\{u_t\geqslant 0\}\in\mathscr{X}_t$. For $t_2\geqslant t_1$ and $A\in\mathscr{X}_{t_1}$, $\int_{\mu|_{\mathscr{X}_{t_2}}}u_{t_2}^+\mathbbm{1}_A\geqslant \int_{\mu|_{\mathscr{X}_{t_2}}}u_{t_2}^+\mathbbm{1}_A\mathbbm{1}_{\{u_{t_1}\geqslant 0\}}\geqslant \int_{\mu|_{\mathscr{X}_{t_2}}}u_{t_2}\mathbbm{1}_{A\cap \{u_{t_1}\geqslant 0\}}\geqslant \int_{\mu|_{\mathscr{X}_{t_1}}}u_{t_1}\mathbbm{1}_{A\cap \{u_{t_1}\geqslant 0\}}=\int_{\mu|_{\mathscr{X}_{t_1}}}u_{t_1}^+\mathbbm{1}_{A}$. \quad
            7. $|u|=2u^+-u$.
        \end{proof}
    \end{thm}
\chapter{Appendix}
\section{Increasing functions}
    \begin{define}
        For $f:\mathbb{R}\rightarrow \mathbb{R}$ a function, $x\in \mathbb{R}$ let
        \[
            f(x+):=\lim_{\varepsilon\downarrow 0}f(x+\varepsilon),\quad f(x-):=\lim_{\varepsilon\downarrow 0}f(x-\varepsilon)
        \]
    \end{define}
    \begin{thm}
        If $f:\mathbb{R}\rightarrow \mathbb{R}$ is increasing, then $f(x+)$ and $f(x-)$ exist for all $x\in\mathbb{R}$, and
        for $a,b\in\mathbb{R}$ with $a< b$ holds
        \[
            f(a-)\leqslant f(a)\leqslant f(a+)\leqslant f(b-)\leqslant f(b)\leqslant f(b+)
        \]

        Furthermore, $f$ has at most countably many points of discontinuity.
    \end{thm}
\end{document}
